% =====================================================================
%  GENERATED; DO NOT EDIT BY HAND.  Built by
%    python3 ebook-method/psykologiske-undersoegelser/tr_parts.py build --write
%  from ebook-method/psykologiske-undersoegelser/tr_template.tex and the translators' parts
%  (texts/hoeffding/psykologiske-undersoegelser/.parts/tr/*.texfrag).  To change the English, edit the part
%  (or this template) and rebuild.
%
%  Harald Høffding, PSYKOLOGISKE UNDERSØGELSER (Vidensk. Selsk. Skr., 6. Række, hist.-philos. Afd. III, 1;
%  Kjøbenhavn 1889): English translation, "Psychological Investigations".  Translated from the Danish,
%  2026-10-09/10, from transcription.tex (the text of record; see its header) by translator agents working
%  from one brief and one glossary (ebook-method/psykologiske-undersoegelser/TRANSLATOR-BRIEF.md,
%  GLOSSARY.md), each part checked by script against the Danish (tr_parts.py build): every page marker once
%  and in order, footnotes, counter resets and paragraph breaks per printed page, the Greek character for
%  character.  The margin numbers [N] are the pages of the 1889 printing.
%
%  CONVENTIONS.  Quotations Høffding gives in German, Latin, French or English stay in the original; what
%  he gives in Danish is translated from his Danish.  Titles of books and journals as printed.
%  Letterspacing (\emph in the Danish) -> \emph.  Footnotes at the same anchor, numbered 1), 2) per printed
%  page as in the print.  Misprints the transcription marks PRINTED AS IS are translated in their evident
%  sense, with a "% print:" comment at the site.  Terms: GLOSSARY.md.  American spelling.
% =====================================================================
\documentclass[12pt,a4paper,openany]{book}
\usepackage[utf8]{inputenc}
\usepackage[T1]{fontenc}
\usepackage{textalpha}
\usepackage{libertinus}
\usepackage{libertinust1math}
\usepackage{graphicx}
\DeclareUnicodeCharacter{0254}{\reflectbox{c}}   % ɔ (should not occur in the English)
\usepackage[english]{babel}
\usepackage[protrusion=true,expansion=true]{microtype}
\usepackage{setspace}
\usepackage[margin=1.2in]{geometry}
\usepackage{fancyhdr}
\setlength{\headheight}{28pt}
\addtolength{\topmargin}{-13.5pt}
\usepackage{hyperref}
\hypersetup{pdftitle={Psychological Investigations (1889)}, pdfauthor={Harald Høffding}, hidelinks}
\pagestyle{fancy}
\fancyhf{}
\fancyhead[LE,RO]{\thepage}
\fancyhead[RE]{\textit{Harald Høffding: Psychological Investigations}}
\fancyhead[LO]{\textit{1889}}
\fancypagestyle{plain}{\fancyhf{}\fancyhead[LE,RO]{\thepage}}
\setstretch{1.3}
\usepackage{marginnote}
\newcommand{\opage}[1]{\marginnote{\footnotesize\textit{[#1]}}}
% Footnotes as the print marks them, numbered per page: 1), 2), 3) ...
\renewcommand{\thefootnote}{\arabic{footnote})}
\newcommand{\tocline}[3]{\noindent\makebox[2.2em][r]{#1}\hspace{0.6em}#2\dotfill\ #3\par}
\pdfstringdefDisableCommands{\renewcommand{\footnote}[2][]{}\renewcommand{\opage}[1]{}}

\begin{document}

% --- p. 1 --- (title)
\opage{1}
\begin{center}
\vspace*{3em}
{\Huge Psychological Investigations.}\\[4em]
By\\[3em]
{\large\textbf{Harald Høffding.}}\\[4em]
{\small Vidensk. Selsk. Skr., 6. Række, historisk og philosophisk Afd.\ \ III.\ 1.}\\[6em]
{\large Kjøbenhavn.}\\[2pt]
{\small Bianco Lunos Kgl. Hof-Bogtrykkeri (F. Dreyer).}\\[2pt]
{\small 1889.}
\end{center}
\clearpage

% --- p. 3 --- (contents)
\opage{3}
\begin{center}{\large\textbf{Contents.}}\end{center}
\medskip
\hfill{\small Page}\par
\tocline{}{Introduction}{5.}
\tocline{I.}{Immediate Recognition}{8.}
\tocline{II.}{The Presuppositions of Association by Contiguity}{36.}
\tocline{III.}{Free Association by Similarity}{41.}
\tocline{IV.}{The Association of Ideas and Comparison}{59.}
\tocline{V.}{On the Concept of Psychical Activity}{67.}
\tocline{VI.}{Psychical and Physical Activity}{81.}
\tocline{VII.}{The Conformity to Law of Psychical Activity}{92.}
\clearpage

% --- p. 5 ---
\setcounter{footnote}{0}%
\opage{5}

\begin{center}\large Introduction.\end{center}

1. The doctrine of the association of ideas occupies a prominent place in
recent psychology. In the endeavor to discover definite laws for the origin and
mutual relations of psychical phenomena, attention was naturally directed to
ideas, whose origin and connections are the easiest to demonstrate. The
conditions and relations of sensations cannot so easily be found by way of
self-observation, but are often discovered only with the help of the physical
conditions to which the sensations are tied. But the appearance and connection
of ideas depend essentially on conditions that are given in consciousness
itself. As for feeling and will, they are not so easy to analyze; although some
principal phenomena stand out with striking clearness, it is nonetheless
attended with great difficulties for untrained observation and reflection to
find the individual elements and their relations. This in large part explains
the one-sidedness with which several schools within psychology have emphasized
knowledge at the expense of the other sides of psychical life. The question of
the laws for the connection of ideas had, moreover, a great practical
importance on account of the undeniable influence that ideas have on feeling
and will: through laws for the connection or association of ideas one would
thus also gain insight into how the character of feeling and will could be
changed. Both the degree and the direction in which psychical life develops are
determined in large part by ideas and their mutual connection. This influence
has been made greater than it really is; the constant and significant reaction
exerted from the side of feeling and will on the development and connection of
ideas has been overlooked; sufficient attention has not been paid to the
peculiar relations that hold in the case of feeling and will.

When in this treatise, which is to discuss various psychological questions
disputed in the most recent times, I begin with an investigation of some main
points in the doctrine of association, it is not the one-sidednesses just
mentioned in older treatments of this doctrine that I mean to oppose. This side
of the matter I hope to have illuminated in my ``Psykologi i Omrids paa
Grundlag af Erfaring''. Instead I shall try to determine
% --- p. 6 ---
\setcounter{footnote}{0}%
\opage{6}what I shall call the \emph{lower} and the \emph{upper limit of the
concept of the association of ideas}. --- The question of the lower limit of
this concept arises because by the expression association of ideas one
understood from the first such a relation between ideas, each of which came
forward on its own as a free and independent element in consciousness, that
when the one was given, the other also appeared. The question is whether the
concept of association is not also rightly applied in cases where we cannot
immediately distinguish several elements, but nonetheless have occasion --- in
any case theoretically or hypothetically --- to conceive the matter as several
elements having fused, elements which under other conditions would assert
themselves freely, each by itself. It is recognition, or rather a certain kind
of recognition, that is here to be investigated more closely. --- The question
of the upper limit arises when we discuss the relation between the association
of ideas and psychical activity. Association has been conceived --- especially
in the older English school and in the Herbartian school --- in such a way that
it excluded any genuine psychical activity, and when association was then also
made the fundamental phenomenon of psychology, it is clear that psychical life
was conceived in a highly one-sided way. It was assumed that association took
place between independent elements, just as in the physical world attraction
connects the independent elements; this analogy gained an unjustified and
overwhelming influence on psychological thinking. When other psychological
schools, in opposition to this --- without denying the factual existence of
association --- maintained the unity and activity of psychical life, the
question becomes: are there two kinds of psychical phenomena, one in which mere
association rules, and another in which psychical unity and activity appear,
--- or is association itself perhaps only a form of psychical activity, so that
the difference between what we call psychical passivity and activity is only a
difference of degree or rests on abstraction?

But the question of the limits of the application of the concept of association
is in turn more or less clearly connected with the question \emph{how many laws
of the association of ideas} must be \emph{set up}. --- The oldest inquirers,
Plato and Aristotle, who mention the association of ideas, already name several
kinds, of which the most important are association by virtue of the likeness of
the contents of ideas (association by similarity) and association by virtue of
the coincidence in time or space of the contents of ideas (association by
contiguity). The third law that Aristotle names, association by contrast, it is
now agreed not to regard as a separate law; the phenomena it was supposed to
fit find their explanation by the two other laws\footnote[1]{Cf. my Psykologi.
2nd ed. p. 185 f.}. Now in recent psychology attempts have asserted themselves
to reduce association by similarity to
% --- p. 7 ---
\setcounter{footnote}{0}%
\opage{7}association by contiguity, so that the latter became the only law of
the association of ideas. This theory appears, as should be noted, in different
contexts in different psychological inquirers. In the older English school it
appears in connection with the attempt to make the association of ideas the
fundamental law of the whole life of consciousness. All processes of
consciousness are first reduced to associations, and thereby to reciprocal
relations between different elements into which consciousness is thought of as
dissolved; next all association is reduced to being a connection by virtue of
external connectedness. Such a standpoint is represented in typical fashion by
\emph{James Mill} in his ``Analysis of the human mind'' (1829). But this latter
reduction has also been attempted by inquirers who make a very sharp
distinction between lower, mechanical and higher, ideal psychical processes,
and who alongside sensation and the association of ideas maintain a special
faculty for comparing or evaluating the elements of consciousness. Here one
perhaps even sees a danger in recognizing the independence of association by
similarity, since thereby phenomena that are thought to testify to pronounced
psychical activity seem to be explained in a purely mechanical way. In any case
the difference is emphatically maintained between association as a lower and
comparing and ``relating'' thinking as a higher kind of psychical process. Such
a standpoint is represented by \emph{Hermann Lotze}.

It is then of interest to investigate the relation from this side as well. But
this in turn leads further. The concept of activity is all too often used in
psychology in a somewhat obscure way. An attempt will therefore be made to give
it greater definiteness than it has had hitherto. An attempt in this direction
would suggest itself all the more if it should turn out that the association of
ideas is through and through a form of psychical activity; for then there will
be all the greater need to have this latter concept investigated. And from
there we are in turn led naturally to another large problem: what is the
relation between psychical activity and physical activity? Finally, it will
also serve to illuminate fully the nature of psychical activity to discuss how
far it is at all points subject to definite laws, or whether it should be
justified to assume that it is to a greater or lesser extent withdrawn from the
dominion of the law of causality.

2. Although this treatise discusses a whole series of different problems, the
foregoing development nonetheless shows that there is a certain connection
between them. The course of the exposition is also given by the foregoing. I
shall 1) first investigate \emph{immediate recognition} and its relation to the
association of ideas. Next I shall 2) go on to investigate \emph{the
presuppositions of association by contiguity} and try to show that among these
presuppositions is also immediate recognition, which in any case is a process
akin to association by similarity. It then remains 3) to decide whether
\emph{association by similarity} is rightly set up as an independent form of
association
% --- p. 8 ---
% print: „nogenside“ for „nogensinde“ (misprint)
\setcounter{footnote}{0}%
\opage{8}alongside association by contiguity. After these main points of the
doctrine of association have been illuminated, I ask 4) about \emph{the
relation between the association of ideas} (\emph{in particular recognition and
association by similarity}) \emph{and comparing activity of thought}, and am
thereby led to consider 5) \emph{psychical activity} as a whole and 6) \emph{its
relation} to \emph{physical activity}, with which it is connected in
experience. While this question leads us beyond the purely psychological domain
to the theory of knowledge and metaphysics, the last question to be raised, 7)
about \emph{the subjection of psychical activity to law}, will lead us from
psychology over into ethics.

\begin{center}\large I. Immediate Recognition.\end{center}

3. Let us begin with quite simple and indubitable observations. Something I see
seems familiar to me. Let it be a face, or, to take something still simpler, a
single feature of a face, a wink of the eye. Or I see in the evening sky a shade
of color that is indeed unusual but nonetheless seems familiar to me. Or a
foreign word is mentioned that I cannot translate, but whose sound nonetheless
has a familiar ring for me. Someone asks me: ``Have you ever been in Les
Plans?'' The name Les Plans is familiar to me, and yet I cannot connect any
idea whatsoever with it. Or let us take examples from inner experience. A
certain organic sensation, a certain mood within the vital feeling, which
surfaces in me, stands before me with a certain stamp of the intimate and the
homelike. I have a quite immediate consciousness that I have had them before,
although I cannot attach a single closer determination to them.

Closely related cases are those in which we try to call up something in our
memory without succeeding, while we recognize it at once when it presents
itself to us by another route, or those in which we are not sure whether we
have perceived a single phenomenon or not, but become sure of the matter by
going through the whole series of phenomena to which it belongs. Here too we
get a quite immediate impression of something familiar. This appears in an
instructive way in \emph{Cattell's} experiments on how many lines, numbers or
letters can be perceived simultaneously. ``When'', he says\footnote[1]{Ueber
die Trägheit der Netzhaut und des Sehcentrums. (Wundt's Philosophische
Studien. III. Leipzig 1886) p. 127.}, ``a certain number was exceeded, the
traces perceived were too weak and of too short duration for me to guess from
them the objects that served as impressions. Although it was now impossible for
me to state these objects straightway, I could nonetheless answer correctly
when I was asked whether a certain written character was there or not. This is
especially the case with long
% --- p. 9 ---
% print: „Fordel“ for „Forskel“ (misprint)
\setcounter{footnote}{0}%
\opage{9}sentences; there I have the peculiar feeling as if I had known the
sentences and forgotten them again''. With this a droll example from the
numerous experiments of recent years with hypnotized persons may be compared.
\emph{Charles Richet} hypnotized one of his friends and, despite some
resistance, forced on him the idea that his name was Durand, so that he
answered with this name when he was asked what his name was. When Richet now
asked him what his name was in the waking state, he was unable to think of it,
although he knew that he ``knew'' it. But when it was mentioned to him, he knew
it at once and repeated it several times with great
satisfaction\footnote[1]{Le somnambulisme provoqué (in the work: L'homme et
l'intelligence. Paris 1885). p. 540. --- As Cattell's experiments seem to
establish, one need not have had a distinct perception of something, let alone
one connected with attention, for it to be able to be recognized. It may indeed
sometimes even be open to question whether the sensation that is recognized was
above the threshold of consciousness the first time. It is the same here as
with the sensations that we infer we have had from our getting certain
sensations of contrast; such an inference is not entirely certain. Cf. my
``Psykologi'', 2nd ed. p. 87 and, besides the facts discussed there,
\emph{Mach's} fine experiments in his ``Beiträge zur Analyse der
Empfindungen''. Jena 1886. p. 107. --- Already \emph{Ludvig Vives} mentions the
recognition of something that was not attentively perceived. ``Qvædam recipit
intelligentia illa simplex, qvæ est de iis, qvæ extrinsecus obveniunt per
videndi vel audiendi sensus, qvæ non \emph{animadversa} tradit statim memoriæ.
Attentio vero velut resurgens considerat in memoria et intelligit interdum
continuo post''. De anima. Brugis 1538. p. 63. That ``intelligentia simplex''
is obviously a pure hypothesis.}.

What is given in such states of consciousness as those discussed here is the
immediate perception of the difference between something familiar and intimate
and something new and strange. This difference is so simple and clear that it
can no more be described more closely than, e.g., the difference between
pleasure and pain or the difference between yellow and blue. It is an immediate
difference of quality that we face here. The peculiar quality with which the
familiar enters consciousness in contrast to the new I shall in what follows
call \emph{the quality of familiarity}. The question is how this quality is to
be explained, whether it is to be taken to be a simple and uncompounded
psychological element, or whether it can be traced back to more compound
relations.

Before I leave the cases cited, which are to form the starting points of the
investigation, I shall draw attention to some further peculiarities in them.

The phenomena that are recognized are here \emph{not uncompounded}, but
nonetheless so simple that they come before consciousness
\emph{simultaneously}. A wink of the eye, a shade on the evening cloud, a short
series of written letters, a short word: these are phenomena with which there
can be no question of successive perception. Experiments have moreover
established that we can simultaneously perceive 4 to 5 unconnected visual
impressions (lines, numbers, letters), and about three times as many when the
impressions are parts of an idea already familiar as a whole. It has also been
shown that practice in this respect extends the
% --- p. 10 ---
\setcounter{footnote}{0}%
\opage{10}range of consciousness. Words, indeed even sentences, are taken up
into consciousness as wholes\footnote[1]{\emph{Cattell} (in the treatise
cited) p. 127. He says expressly: ``Als an mir selbst Versuche angestellt
wurden, bemerkte ich, dass die Eindrucke simultan in das Bewusstsein
gelangten''. --- \emph{Berger}: Ueber den Einfluss der Uebung auf geistige
Vorgänge (Wundt's Studien, V. 1888). p. 177.}. --- This is the one peculiarity
that must be firmly held onto if immediate recognition is to be distinguished
from other kinds of recognition. By it immediate recognition is distinguished
on the one hand from cases where quite uncompounded sensations, which do not
occur as isolated in our ordinary experience, are to be recognized, on the
other hand from cases where the perception of the phenomenon, and consequently
also the recognition, must take place through a successive psychological
process.

The second peculiarity is that the recognition is \emph{unprepared}, i.e., it
is not shortly preceded by any idea, voluntarily or involuntarily called up, of
the phenomenon that is the object of the recognition. --- In some of the cases
cited there was indeed in a way a preparation, in that one had tried to bring
forth the idea of the phenomenon, but since this did not succeed, this
preparation is at most indirect, although it is by no means insignificant on
that account\footnote[2]{Cf. my Psykologi, p. 83 f. --- B. \emph{Erdmann}
(Zur Theorie der Apperception. Vierteljahrsschrift fur wissenschaftliche
Philosophie. 1886. p. 343) therefore distinguishes between ``ein erregtes und
ein unerregtes Unbewusste'', of which the first forms the immediate basis of
consciousness.}. This question of preparation will on the whole prove to be
much more intricate than has usually been assumed. There are countless degrees
here between complete unpreparedness and complete preparedness. The most
important differences we shall have occasion to bring out in what follows.

Thirdly, it must further be emphasized that self-observation in the cases
cited \emph{does not show the slightest trace of other ideas} \emph{that are
aroused} by the \emph{recognized phenomenon} and could be assumed to play a
role in the recognition itself. Insofar, then, as anyone would assume that all
recognition presupposes such ideas, the burden of proof lies on him, and if
immediate recognition, as it appears in the cases cited, can be explained
without such an assumption, that will be the only scientific explanation.

4. A way to the simplest explanation will appear when we compare more closely a
new sensation with a recognized sensation. For brevity's sake I use, in what
immediately follows, the expression sensation instead of complex of sensations
or group of sensations. It is, after all, as already emphasized, as a rule not
single sensations that are at issue here.

By a new \emph{sensation} we need not here understand an \emph{absolutely} new
sensation, i.e., one that we have never at any time had before. Such absolutely
% --- p. 11 ---
\setcounter{footnote}{0}%
\opage{11}new sensations we can, when we keep to the simplest cases, have only
at the very first beginning of the life of consciousness or --- as far as the
general sensations are concerned --- at the beginning of new periods of life
(the years of puberty, e.g.) or at the outbreak of new illnesses. We can also be
said to have a new sensation when a certain, not too short, time has passed
since we last had a similar sensation. When I come home one evening and, after
undressing, wonder whether I have in fact also locked my front door, I shall be
able to decide this without trouble only if I can reproduce the sensations that
accompany my locking of the door, and \emph{find} them fresher and more
distinct than such reproduced sensations as correspond to earlier lockings. If
I believe I have shut the door without really having done so, the reason is
that the reproductions of earlier vivid sensations of this action have produced
in me the semblance that a repetition of the action took place this evening. An
illusion of memory then arises. --- But under ordinary conditions a sensation
is relatively new when it occurs after a certain interval. How long this must
be for the sensation still to be called new will differ in the different cases
and for the different sense modalities. A tone was recognized with absolute
certainty after an interval of two seconds; even with an interval of five
seconds the influence of time was not distinctly noticed; with intervals of 10
and 15 seconds, on the other hand, a considerable increase of incorrect
estimates was noticed, and with an interval of 30 seconds certainty had quite
vanished\footnote[1]{\emph{Wolfe}: Untersuchungen über das Tongedächtniss
(Wundts Studien III. Leipzig 1886). p. 522. --- Cf. the facts cited in my
Psykologi p. 141, which show that the time interval can also bring it about
that the differences are not perceived, which is confirmed by the results
reported by Wolfe ib. p. 541 f.}. In more compound cases the influence of time
will probably not make itself felt so soon, since there are more points of
support for attention, the relation between the elements, namely, as a rule
playing a larger role than any of the single elements by itself alone.

If we now compare a (relatively) new \emph{sensation} with a repeated and
recognized one, there is one condition that they have in common in any case,
namely that an impression has come to the brain (consciousness), and that this
impression stands in such a relation to the state of the brain (of
consciousness) that a certain definite sensation (with a certain quality and
independence) can appear. What, now, can it be that \emph{recognition}
presupposes beyond this?

In answering this question we must first and foremost keep to the only
circumstance that differs in the two cases: in the one the impression is
(relatively) new, in the other repeated. The only effect this circumstance can
have is that a reproduction becomes possible. But this reproduction need not
lead to what is reproduced appearing as an independent member in
consciousness, and in the present cases this does not happen either. Their
peculiarity consisted, among other things, precisely in their
% --- p. 12 ---
\setcounter{footnote}{0}%
\opage{12}seemingly uncompounded character. Besides the recognized feature or
features there is not the least thing in consciousness that has to do with the
recognition. The word ``Les Plans'' sounds familiar, and this quality of
familiarity in the sound is the \emph{whole phenomenon}. Although I at once
become aware of the psychological interest of the phenomenon, I can nonetheless
not discover the very least trace of contributing sensations or ideas. Now when
it nonetheless stands firm that the quality of familiarity must be connected
with my having heard the word before, this quality can only be explained by the
fact that the earlier state into which I came when I heard the word for the
first time can in one way or another assert itself anew together with the
repeated sensation itself, without however appearing as an independent element
in consciousness alongside it. This we can very well express by the word
reproduction, if we assume that the impression now arriving in the brain
(consciousness) arouses a state that is \emph{also} determined by the
after-effect of the earlier impression. And this after-effect, we can simply
assume, consists in the greater ease with which the state sets in upon
repetition. We appeal here to the law of practice that holds for all organic
tissue, not least for nervous and muscular tissue. The more frequently a
process has gone on in an organ, the more easily it sets in anew. Even outside
the organic domain one could speak of practice, since, as is well known, the
sounding boards of musical instruments are the more easily set into vibration
the more they are used. Every new attempt, every new impression profits from
the earlier impressions, these having, as it were, paved the way for it. ---
\emph{Psychologically} we can express what happens in recognition by the
formula (\textit{A} + \textit{a}), or perhaps better $\left({a\atop
A}\right)$, where \textit{A} means the sensation, \textit{a} its reproduction
(the idea), and where the parenthesis signifies that these two do not here
appear as independent members in consciousness, but are \emph{thought
theoretically} as factors in the seemingly uncompounded phenomenon. The
justification for this mode of expression lies in the fact that the
after-effect of \textit{A} which makes recognition possible would under other
conditions have made possible \textit{a'}s appearance as a free (independent)
idea, e.g. if by association by contiguity \textit{B} called up \textit{a.} The
meaning, then, is: the \emph{same} thing that is a \emph{condition} for
\textit{B being able under certain conditions to recall} \textit{a,} is
\emph{also a condition} for \textit{A} being able to be \emph{recognized}. The
name ``Les Plans'' could perhaps have been called up by the sight or the idea of
the place; but the name can, when I hear it, seem familiar to me even if it
cannot call up the idea of the place. In the quality of familiarity, then, in
the latter case the same thing is at work that in the former case would have
appeared as an independent idea. --- How one will conceive \emph{physiologically}
what goes on through practice in the minute parts of the organism, here of the
brain, is a matter by itself, into which we need not go further here. The most
natural assumption would presumably be that the first impression brings about a
rearrangement of the molecules, which after the cessation of the impression is
again replaced by the earlier state, but in such a way that this
% --- p. 13 ---
\setcounter{footnote}{0}%
\opage{13}is now more unstable, \emph{more easily brought} out of
\emph{equilibrium}. To that extent a certain disposition to the same
rearrangement can be said to have been produced, so that it takes place more
easily when the same impression comes again\footnote[1]{This assumption is so
far from being anything modern that it is already stated with great clarity by
\emph{Charles Bonnet} in his Essai analytique sur l'âme (Copenhague 1760),
p. 61, where it says: ``La souplesse ou la mobilité des fibres augmente par le
retour des mèmes ébranlements. Le \emph{sentiment attaché} à \emph{cette
augmentation de souplesse} ou de \emph{mobilité}, \emph{constitue la
réminiscence}.''}. Recognition, or rather the quality of familiarity, then
forms the psychological correlate of the greater ease with which a change in
the arrangement of the brain molecules concerned is produced. No other
circumstance presents itself by which the peculiar difference between sensation
and recognition can be explained. Now the difference between greater and lesser
possibility of rearrangement is admittedly only a difference of degree, while
psychologically regarded there is a difference of quality between sensation and
recognition. But this cannot surprise us when we consider that to the
psychological qualities there always correspond in the physical world only
differences of quantity. The color qualities have their physical correlates in
quantitative differences between the vibrations of the ether, and probably also
their physiological correlates in quantitative differences between the motions
of the molecules in the visual centers of the brain. Our explanation of the
phenomenon of recognition in its simplest form is thus in good analogy with
what seems to hold for other simple psychical phenomena.

I only note further that the theory set out here is a further development of
what I have already set out in my Psykologi (2nd ed. p. 140---142). I have here
sought to give closer determinations and explanations, which have proved to be
necessary.

5. Sensation and recognition can often move so close to each other that no
decisive difference can be made between them. This is connected with the fact
that we hardly have any sensation without a certain attention being turned
toward it. And this in turn comes from the fact that with every sensation there
follows a feeling of pleasure or pain and therefore an urge to activity and
movement, which is naturally directed toward what calls up the sensation. But
the attention thus aroused does not work entirely continuously. It lets go of
its object and seizes it again shortly after, when the impression, and with it
the motive of attention, continues to act. Wherever one may seek the cause of
this, the attention we direct toward a continuous impression is discontinuous
and periodic. This is seen most clearly with very weak sense impressions. When
a pressure on the tip of the finger continuously increases, we do not notice
this continuous increase; attention holds together different stages of the
rise of the impression and thereby we become conscious of
it\footnote[2]{\emph{Stanley Hall}, cited in \emph{Ribot}: La Psychologie de
l'Attention. Paris 1889. p. 16.}. With very small sense sensations lying just
above the
% --- p. 14 ---
\setcounter{footnote}{0}%
\opage{14}threshold of consciousness, experience shows that at one moment they
appear in consciousness, but at the next moment vanish; they are in a constant,
rhythmic movement which, when the impression acts continuously, can only be
explained by the periodicity of attention. When one and the same image can be
perceived differently, the different perceptions tend to alternate with one
another at certain intervals. An outline of a staircase (the so-called Schröder
staircase figure) one can also ``see'' as a wall from which stones have fallen
out at the bottom. If one regards this image attentively and without
interruption, staircase and wall alternate periodically. It is not the image,
nor indeed properly the sensation, but it must be attention that
alternates\footnote[1]{N. \emph{Lange}: Beiträge zur Theorie der sinnlichen
Aufmerksamkeit und der activen Apperception (Wundt's Studien. IV. Leipzig
1888). p. 395. 406.}. It seems as if the full sum of force that the act of
attention requires cannot be present at every moment, which is why at certain
intervals some time must be spent in replacing the energy consumed. But as
attention thus seizes its object again after having let go of it, the
reproduction of the preceding grasp of the object always comes to play a part
as well. When an attentive regarding makes a simple object clearer to us, this
can only happen by such reproduction (conceived in the manner set out in 4)
strengthening the sensation. Through this interaction of sensation and
reproduction, thus through a series of acts of recognition, the seemingly
continuous contemplation runs its course. This is seen clearly when we
\emph{voluntarily} try to hold fast a definite interpretation of the staircase
figure mentioned; we try earnestly to persuade ourselves that it is to be a
staircase, which happens through voluntary reproduction of staircases seen
earlier\footnote[2]{Cf. N. \emph{Lange} ib. p. 407. See also Wundt:
Physiologische Psychologie. 3. Aufl. Leipzig 1887. II. p. 238. 253 f.}, which,
however, do not come to appear as independent elements of consciousness, but
only serve to strengthen the sensation. In \emph{involuntary} attention it must
of course be the impression itself which, as soon as the inner condition,
sufficient energy, is present, always anew releases attention and reproduction.
--- In this way, then, a greater or smaller series of acts of recognition is
woven into our seemingly quite simple and uncompounded experiences.

6. An immediate recognition can take place just as well with respect to an idea
(or a group of ideas) as with respect to a sensation (or a group of
sensations). We can recognize memory images and images of the imagination. When
in a waking dream an image involuntarily pushes itself forward before me, it
can come with quite the same quality of familiarity as that which a landscape
has when I see it anew. The memory image then profits from the earlier
impressions. It is a not uncommon experience that we do not
% --- p. 15 ---
\setcounter{footnote}{0}%
\opage{15}recognize memory images, that they stand before us as strange, and
that we must first laboriously convince ourselves that they correspond to real
experiences. In such cases they lack that quality. And even if we are sure that
we have experienced before what we represent to ourselves, we can be unsure
\emph{where and when}. With images of the imagination it is not earlier real
experiences but imagined experiences that are at issue. I may, e.g., have
accustomed myself to picturing a poetic figure in a certain definite way, and I
then know this figure when my imagination involuntarily brings it forward
before me. This is a phenomenon that will be particularly frequent with poets
who produce instinctively and with a lively imagination. Cf. \emph{Goethe's}
famous lines:

Ihr naht euch wieder, schwankende Gestalten!

Die früh sich einst dem trüben Blick gezeigt.

Here the same law is at work as in the recognition of something really
experienced, namely the law of practice. It holds just as well for the brain
functions to which memory and imagination are tied as for the one to which
sense sensation is tied.

7. Although immediate recognition appears to self-observation as a quite simple
phenomenon, we can nonetheless, then, regard it \emph{theoretically} as
compound. There is a moment of memory, of after-effect from the past, which in
it immediately combines with the sensation aroused at the moment. What appears
in consciousness is the product of both moments; their fusion has taken place
below the threshold of consciousness. I have therefore in my Psykologi (V B, 1)
designated the phenomenon as \emph{bound memory}, in contrast to such phenomena
of memory in which what is reproduced appears as an independent element in
consciousness; here a free \emph{memory} takes place.

If it is justified to call immediate recognition a phenomenon of memory, it must
also be correct to call it a phenomenon of association. Everyone is of course
free to use words as he likes; and if by association one will understand only
such a connection of ideas in which both of the connected members assert
themselves in a free and independent way in consciousness, immediate
recognition cannot be called association. Thus \emph{David Hartley}, the second
founder of the theory of association, clearly describes both immediate
recognition and free association by similarity, although he uses the word
association only of association by contiguity\footnote[1]{Observations of man.
The London edition of 1791. I. p. 291 ff.}. And --- to take an example from the
most recent times --- \emph{James Ward}, who rejects free association by
similarity, nonetheless acknowledges the immediate cooperation, as described,
of sensation
% --- p. 16 ---
\setcounter{footnote}{0}%
\opage{16}and reproduction, but declares that it is assimilation and not
association. If it is to be called association, it is automatic
association\footnote[1]{Art. ``Psychology'' in the 9th edition of Encyclopædia
Britannica. Vol. XX. p. 60 b.}. --- Immediate recognition (or, as I also call
it: perception) is a border phenomenon which has the seemingly uncompounded
character in common with sensation, and which has in common with free
association that reproduction clearly enough plays a part. From free
association it is distinguished by the complete fusion of sensation and
reproduction. In immediate recognition elements are combined which we can not
only \emph{think} of as separate, but which can also occur each by itself under
other conditions. (Cf. 4.) Recognition does not always set in at the first
moment. \textit{A} can come up as mere sensation and only after a short
interval become $\left({a\atop A}\right)$; and on the other hand \textit{a}
can be called up in consciousness by free association, thus as an independent
element. Immediate recognition is related to association as the tangent to the
secant. In the tangent the two points of intersection with the circle have
fused into one; it is a secant whose points of intersection lie infinitely
close to each other. It would be association by similarity, if anything, with
which one would have to compare immediate recognition according to this
analogy. --- The question whether there is a real association by similarity we
shall go into more closely later.

8. Just as immediate recognition as a whole is an example of the psychological
influence of practice, so it will also undergo new changes with continued
repetition. With repetition it takes place faster and faster, and can take
place so fast that it requires no separate act of consciousness, but merely
serves to release subsequent acts of consciousness. It then approaches the
threshold of consciousness, if it does not cross it. Experiments have shown
(what ordinary experience moreover already establishes) that with acquaintance
with a language the speed increases with which one can read in this language
letters that form words and words that form sentences. The familiar experience
that foreigners seem to us to speak faster than we ourselves do is explained by
this. If on the other hand we take words that do not form sentences, and
letters that do not form words, then the time required for the sentence is
about doubled\footnote[2]{\emph{Cattell}: Ueber die Zeit der Erkennung und
Benennung von Schriftzeichen, Bildern und Farben (Wundt's Studien. II. Leipzig
1885). p. 646 ff.}. When the physiological time, the time taken up in receiving
an impression and answering it by a movement, is shortened by repetition and
practice, the whole process can probably now run its course through
subordinate brain centers, without the higher centers, to which the combining
processes of consciousness are tied, needing to take part. The impression then
immediately releases the
% --- p. 17 ---
\setcounter{footnote}{0}%
\opage{17}movement, without one's being conscious of having received the
impression. Recognition thus fades over into the unconscious by virtue of the
same causes through which it originally came into being. But it has been a
necessary condition for the processes concerned now being able to take place
with such great certainty and speed. It is a first stage of adaptation and
conditions the onset of the subsequent stages of adaptation. In everyday speech
we do not always distinguish between these different stages. We say, e.g., of
the practiced musician that he \emph{knows} the notes, although he need not at
all apply any separate act of consciousness at the sight of them, since they can
immediately release associations by contiguity or movements. The usage is
within its rights; for this \emph{secondary} unconsciousness, which is due to
practice, must not without further ado be put together with the \emph{primary}
unconsciousness that is present when no impressions and sensations have arisen.
And to this it must be added that the boundary between conscious and
unconscious ``knowing'' can be difficult to draw. When an act of recognition
stands as the introducer of a series of significant, affect-arousing ideas, we
easily overlook it, our whole interest being turned toward what follows. What
has only a minimum of interest vanishes (especially for untrained
self-observation) in relation to what sets interest in stronger motion, and is
not preserved in memory. Over the explosion we forget the spark that caused it.
This is connected with the fact that it is overwhelmingly practical
considerations that determine the development of our ideas and of our
knowledge. We have no practical interest in paying particular attention to the
sensations of movement that play so large a role in the perception of distance,
whereas it is of great interest to be sure in judging the distance between us
and the things around us.

When instincts are released by mere impressions and sensations, it is surely in
large part secondary recognitions, become automatic, that we have before us
here. The young secretary bird, which at once strikes at a snake (so eagerly
that one can get it to strike at a gut instead), presumably does not pay much
attention to what calls up its attack. But the impression of the snake's shape
must nonetheless, as it were, set old strings in motion, must be able to
propagate itself into its interior along previously paved paths, if it is to
have this effect. When a rat starts at a cat's cry even if it has lost the
cerebrum, it is because it has adapted itself to this impression in such a way
that the impression can have its effect even if the highest central organs are
missing. Proceeding from these highest organs, the centers of movement have
been so trained that they can be set functioning by an impression or a
sensation that otherwise would not be sufficient for this. The adaptation here
benefits not only the single individuals but also their descendants, which can
be seen, e.g., from the way in which the nutritive instinct of newborn animals
and children expresses itself.

% --- p. 18 ---
\setcounter{footnote}{0}%
\opage{18}9. It must be well noted that only \emph{immediate} recognition has
been discussed here, the phenomenon of recognition in its simplest form. If the
explanation of this phenomenon that I have given in the foregoing is correct,
it has the interest of presenting a transitional form between mere sensations
and ideas proper (free ideas). I have therefore in my Psykologi (V B, 1) set it
out after the conclusion of the doctrine of ideas. --- But this interest would
be founded on illusion if the objections that have been directed in the most
recent times against the view set out were justified.

Recognition, it has been held, is never an entirely simple phenomenon, but must
always be explained according to the usual (free) association of ideas. And
since it has then also been denied that there is association by similarity
proper, it must be the law of association by contiguity that is laid at the
basis. The process that is supposed to take place in all recognition must then
be the following. I approach a house that I have seen before, and fix my gaze
on a single part of it, e.g. a window. This window now arouses in me (by
association by contiguity) the idea of an adjacent part of the house, e.g. an
ornament. If I call the sight of the window \textit{A,} the idea of the
ornament \textit{b} (capital letters being used throughout to designate
sensations or groups of sensations, and small letters to designate ideas or
groups of ideas), then I now have \textit{A} + \textit{b.} With the idea of the
ornament there also naturally arises an expectation of getting to see it, and
as this expectation is fulfilled, when the gaze leaves the window and moves on
along the house, there sets in the sensation \textit{B.} It is then now,
according to this theory, the agreement between the idea \textit{b} reproduced
\emph{beforehand} and the sensation \textit{B} setting in \emph{afterwards} that
conditions the recognition, or in which it consists. And when this process
repeats itself with the different parts of the house, we say at last that we
recognize the house.

With respect to this theory I shall first of all remark that the psychological
process it depicts is entirely natural and constantly finds application in our
actual life of consciousness. I have myself brought it out in my Psykologi (V
B, 2 and 4); but it certainly did not occur to me that it could be used to
explain all recognition and to dispute that there was \emph{immediate}
recognition in the stated sense of the word. From the circumstance that there
is \emph{successive} recognition, which takes place in the manner described, it
surely cannot without further ado be inferred that there \emph{cannot} be
\emph{immediate} recognition.

For it is clear that the theory cited does not fit the simple phenomena of
recognition cited earlier, those cases in which one or more features appear to
us with the stamp of familiarity without any other idea whatsoever connected
with them being traceable in consciousness. These simplest cases of recognition
may very well concern \emph{relatively} compound objects; recognition of
compound objects is, namely, not the same as successive recognition, as long as
it is possible to comprehend the given content by one act of consciousness.
(Cf. 3.) The main thing is that the
% --- p. 19 ---
\setcounter{footnote}{0}%
\opage{19}quality of familiarity immediately asserts itself and stands isolated
in our consciousness, without arousing any idea. There is here, then, no
question of any successive movement of observation from part to part. It would
be an entirely unjustified construing to put such a successive process into
those simple phenomena.

But even where such a process goes on, namely with respect to such objects as
cannot be comprehended by one act of consciousness, the theory cited leads to
conflict with experience. For it presupposes that of a compound object
\textit{X} = \textit{A} + \textit{B + C} (where the letters mean the parts or
properties of the object) one \emph{never} recognizes the first part that is
perceived, but \emph{always} first the second or third. \textit{A} can be
recognized only when one later returns from one of the other parts to
\textit{A.} Yet does experience not show that one can very well recognize the
very first member?

10. The assertion cited, that all recognition rests on an idea aroused
beforehand being followed by the corresponding sensation, is grounded on a
\emph{confusion of immediate recognition itself with the verification of this
recognition}. Immediate recognition, which happens in a now and is a quite
simple mental process, cannot in and of itself establish its own validity. Both
sense illusions and illusions of memory establish sufficiently that the quality
of familiarity can be present without there being a real repetition of the
phenomenon we believe we have before us. If the correctness of the recognition
is to be tested, this can only happen by \emph{investigating the ideas}
\emph{that the recognized phenomenon arouses} in me. \emph{This is impossible}
in \emph{such cases} as those \emph{discussed} in 3, in \emph{which something
stood as familiar} and yet \emph{caused no associations}. \emph{Precisely such
cases are therefore of interest as experimenta crucis for testing} the
\emph{different theories}. But for the most part what is recognized will at
once, by association by contiguity, arouse ideas of features connected with it,
and one can then test whether these accompanying features are also really
present. (Cf. my Psykologi V B, 4 on the experimenting method by which
consciousness involuntarily subjects the ideas that surface to an ordeal by
fire and thereby learns to draw the boundary between the possible and the
actual. Consciousness here proceeds practically in the same way as science: it
sets up a provisional hypothesis and seeks to verify it.) In the perception of
very compound objects the verification can thus accompany the perception
itself step by step. But it \emph{need} not go so. This is seen already from
the fact that we can have perceived a series of phenomena without its being
possible for us to reproduce freely the single members of the series, although
we \emph{know} them again as soon as they present themselves anew. (Cf. 3.)
When I read a good book a second time, the peculiar enjoyment (which can be
called the enjoyment of recognition) surely does not rest on my uninterruptedly
thinking of what is to come, and constantly having this not particularly
% --- p. 20 ---
\setcounter{footnote}{0}%
\opage{20}purposeful\footnote[1]{\emph{David Hume} would in any case have
declined such anticipations, and there are surely many who share his taste. Cf.
his ``Essay on Simplicity and Refinement'' (Essays. London 1777. Vol. I,
p. 211): If the merit of the composition lies in a point of wit, it may strike
at first, but the mind anticipates the thought in the \emph{second perusal},
and is no longer affected by it. When I read an epigram of Martial, the first
line recalls the whole, and I \emph{have no pleasure in repeating in myself
what} I \emph{know already}. But each line, each word in Catullus has its
merit; and I am never tired with the perusal of him.'' --- It is obvious that
Hume did not subscribe to the expectation theory. If it held unconditionally,
Hume would not have read any book a second time.} expectation confirmed. But
it rests on each single feature that presents itself gliding more easily into
my consciousness, or arising in it with the quality of familiarity. We dwell
on the single features; they stand immediately before us as old acquaintances.
There was more expectation the first time we read the book; then we hurried
forward because we knew that something came after; now this curiosity is
satisfied, and we can dwell securely at each place as if we belonged there.

That verification, where it is possible, most often goes on involuntarily comes
from the fact that, when consciousness or the brain is in a healthy and
vigorous state, there is an urge to apply energy to ideational activity. There
is a (relatively) spontaneous ideational activity, just as there is a
(relatively) spontaneous movement. Every sensation or idea that surfaces leads
the ideational urge in a definite direction according to the laws of
association. Every idea will therefore, when the energy of consciousness is not
weakened, set series of ideas in motion. Therefore immediate recognition also
easily passes over into successive recognition through aroused and verified
ideas and recognitions.

Illusions can sometimes arise here which cannot be explained according to the
expectation theory, but on the contrary seem definitely to conflict with it. I
can believe I know a face, although it is only a single feature, e.g. the eye,
that resembles that of a person known to me. This is explained easily and
simply if we assume that the immediate recognition of the single feature has
made the other features recede into shadow for us. But if the eye had aroused
in me ideas of the forehead, mouth, etc. of the person I really know, then these
ideas would of course at once come into conflict with the images of the
forehead and mouth that the present person has; the expectation would be
disappointed; --- and how then was the illusion of recognition possible? If I
began the perception of the face with the eye (\textit{A}), this would arouse,
e.g., ideas of the forehead (\textit{b}) --- but when there then comes as the
next member in the successive perception --- not \textit{B}, but ---
\textit{P}, then according to the theory mentioned the imagined recognition
seems unable to come about at all. ---

Verification does not always go on involuntarily; it must sometimes, if it is
to happen at all, be undertaken with a definite intention and with voluntary
attention. We then say: ``This is the window \textit{A}; if this window is now
to be the one I have seen on the house \textit{X}, the ornament \textit{B} must
be found beside it''. The expectation is here formulated as a hypothesis,
% --- p. 21 ---
\setcounter{footnote}{0}%
\opage{21}whose validity rests on certain conditions, and the recognition (if
it sets in) is formulated as a judgment by which the agreement between
\textit{A} and\textit{ a}, the image of the present window and the idea of the
window seen earlier, is established. In this judgment what in an immediate
recognition would be immediately fused appears as separated. We get
\textit{A = a}, instead of $\left({a\atop A}\right)$. In the judgment of
recognition sensation and idea are put in a definite and conscious relation to
each other. It would not be correct to call the immediate act of recognition
itself a judgment. One would thereby (as when one would call effects of
contrast and sense illusions inferences) overlook the difference between a
reflected and an entirely immediate phenomenon of consciousness.

In many cases an inference can be drawn from the logical context in which a
phenomenon occurs, without any recognition brought about by association by
contiguity seeming to take place. --- I recently read a little boy's German
lesson with him. Some hours later I asked him: ``What does `Baarschaft'
mean?'' This word had occurred in his lesson, but it was now entirely strange
to him, and he could not at all remember having heard it before. But when I
mentioned the sentence: ``Er hatte alle seine Waaren in Baarschaft angelegt'',
he knew at once what the word meant; but that was because he \emph{inferred}
its meaning from the context; the word did not on that account sound more
\emph{familiar} in his ears; for that, several more repetitions were needed.
--- We shall in what follows have use for this distinction between knowing a
word's meaning (being able to name another word that expresses the same) and
recognizing it in the immediate and proper sense (so that it appears with the
quality of familiarity like ``Les Plans'' in 3).

11. We cannot do without theoretical constructions in psychology. I have used
one such construction when I represent immediate recognition as a fusion of a
sensation and an idea: $\left({a\atop A}\right)$ or (\textit{A} + \textit{a}).
But the expectation theory must make use of constructions that go much further,
since in \emph{every} act of recognition, even the simplest, it must
distinguish between four different operations: 1) the appearance of
\textit{A}; 2) the appearance of \textit{b} by association; 3) the appearance
of \textit{B}; 4) the comparison of \textit{b} and \textit{B.} If all this
really took place, we should, however, undoubtedly notice something of it. We
should in any case notice that something was going on in consciousness, even
if we did not quite know what it was. (In any case, those authors who reject
the opposite explanation as grounded in ``superficial'' observation must
assume that the process is noticed.) For one can very well have a sense that
more elements of consciousness are present than one can perceive with full
attention. Thus \emph{Wundt}\footnote[1]{Physiologische Psychologie. 3. Aufl. II, p. 247.} remarks: ``When more
impressions present themselves than can be apperceived, one is clearly
% --- p. 22 ---
\setcounter{footnote}{0}%
\opage{22}conscious that still other impressions were present; but one is not
able, or able only with the help of a distinct succession, to make them
definitely present to oneself.'' What holds of impressions must also hold of
ideas. But it is certain that there are cases of recognition in which no more
elements of consciousness are noticed at all than one can survey, and in which
precisely this gives the whole state its peculiar character.

Even in the successive perception of a known object there seems to be nothing
to prevent the recognition from consisting in a series of immediate acts of
recognition according to the formula $\left({a\atop A}\right)$, without ideas
aroused beforehand \emph{needing} to play any part. The imagined recognition
discussed in 10 points decidedly in that direction.

It is on the whole an altogether arbitrary assertion that recognition should
always presuppose an idea aroused \emph{beforehand}. All recognition rests on
a connection between an idea and a sensation, but this idea can very well be
aroused by the sensation itself. In illusions of sense this is evidently the
case. On a walk I once thought I saw a small lake in the distance, but soon
became aware that it was a shining white house that I had perceived as
glittering water. Shortly before, I had looked at the map, which showed a
small lake in the direction in which the house lay; but I had not thought
further about it and was not at all looking for the lake when it suddenly
seemed to show itself out there. The illusion was here altogether momentary
and had so little objective ground in the impression that, had the idea of the
lake not been so easily aroused, the illusion would hardly have arisen.

Perhaps it will be said, on the occasion of this case, that the idea was
nevertheless present beforehand, even if it was not conscious. At least some
of the authors who hold that the expectation theory explains all recognition
declare expressly that by the ideas which they assume to precede every
recognition they mean ``unconscious ideas'' or nervous processes to which no
elements of consciousness correspond\footnote[1]{Thus \emph{Arne Löchen} (Spørgsmaal vedkommende de afasiske
Sygdomme. Kristiania 1888. p. 128 f.) and \emph{Alfr}. \emph{Lehmann} (Om
Genkendelse. Det kgl. danske Videnskabernes Selskabs Skrifter, 6. Række, II, 2.
p. 15). On the other hand, from \emph{Kroman's} exposition in ``Kortfattet Tænke-
og Sjælelære''. 2nd ed. Copenhagen 1888. p. 146 ff. one gets the impression
that he conceives the whole composite process as conscious. In any case no
reservation is made.}. As soon as one conceives the contributing ideas in this
way, the whole dispute really comes to an end. For ideas that are not above
the threshold of consciousness exist only as dispositions to ideas, and the
decisive question is how these dispositions are aroused to actuality. Here I
now maintain that this can happen, and often does happen, through the
repetition of the corresponding sensation. It goes without saying that this
repetition would in itself have no significance for recognition if there were
no disposition (trace, residuum, after-echo or whatever one wants to call it)
that could be aroused. \emph{The interval} during \emph{which}
% --- p. 23 ---
\setcounter{footnote}{0}%
\opage{23}\emph{the idea has been present merely as a disposition}
(``\emph{below the threshold}''), \emph{can of course be of very different
length}, and the whole question in dispute can here easily reduce itself to a
pure question of degree. When after the lapse of 20 years I meet a person
again and at once recognize him, the disposition time (if I may call it so)
has been very long. In the illusion of sense just mentioned, the disposition
time was only about a quarter of an hour. In the recognition of the ornament
(in 9) the disposition time was hardly a minute. Even if we cannot voluntarily
call up an idea, and even if it does not come forth in consciousness on
occasions that are in and of themselves natural, it still cannot be said with
certainty to have altogether disappeared as a disposition\footnote[1]{\emph{Leibniz} rightly says: ``Je ne vois aucune nécessité, qui
nous oblige d'assurer, qu'il ne reste aucune trace d'une perception, quand il
n'y en a pas assez pour se souvenir qu'on l'a eue.'' (Nouveaux Essais. I, 3, 24.
The remark is directed against \emph{Locke}, who denied that one could be said
to have an idea when it could not be called up in memory.)}. It may be that a
quite peculiar combination of circumstances is required to arouse the
disposition to actuality. Just as we are on the whole too inclined to derive a
phenomenon from a single cause or, perhaps more correctly, to find the cause in
a single circumstance, so we also easily let ourselves be caught by our
schematic way of conceiving things in the opinion that an association of ideas
can be brought about starting from one given element of consciousness. What a
single element would not be able to release, several might perhaps be able to
effect if they followed one another so quickly that their effects could
accumulate. We may have received impulses from several different sides in the
direction of a certain association of ideas, so that this association is
thereby so prepared that now an insignificant occasion is enough to bring the
idea up over the threshold of consciousness\footnote[2]{R. \emph{Wahle} has drawn attention to such
``supplementary and competing'' relations in association in his treatise
``Bemerkungen zur Beschreibung und Eintheilung der Ideenassociationen''
(Vierteljahrsschrift für wissenschaftliche Philosophie IX) p. 414---420.}. In
all recognitions and associations such relations surely play a part. The
threads of our psychical life are so intimately interwoven or entangled that
the principles and laws we can formulate in our psychology can give only the
very barest outlines.

An \emph{altogether unprepared} recognition does not exist in the strictest
sense, since all recognition presupposes a disposition which stems from the
first perception of the phenomenon that is now to be recognized. But there are
all possible degrees of preparation. In some of the cases mentioned in 3 the
recognition was preceded by a voluntary searching, which did not, to be sure,
lead to the goal, but which has hardly been altogether without significance,
since experience frequently shows that even if we do not find the idea sought
at once, it can later, without external occasion, emerge in consciousness. We
have a more definite preparation when the idea has been in consciousness
shortly before through voluntary or involuntary recall, even if we do not hold
it fast uninterruptedly until the phenomenon comes again (
% --- p. 24 ---
\setcounter{footnote}{0}%
\opage{24}which, after all, we are not able to do either --- cf. 5). --- But
all this is only differences of degree. When the expectation theory lets its
``expectation'' be unconscious, it has really given itself up at once. It
will not be possible to deny that it is the \emph{subsequent sensation}
\emph{that helps the} ``\emph{unconscious}'' \emph{idea over the threshold}.
The \textit{unconscious b aroused by A} \textit{would remain forever} in the
\emph{unconscious world} \emph{if} \textit{B did not} come along and
\emph{help it forth}. And since \textit{b} and \textit{B} do not always
appear on this occasion as independent elements of consciousness, we get the
often-mentioned fusion or immediate recognition\footnote[1]{Already in my Psykologi p. 141 I have pointed out that
recognition in many cases becomes possible only through a preparation that
lies a certain \emph{small} interval of time beforehand.}.

An example will show this clearly. I copy out some lines that are found in a
book on p. 270. Is the page reference now also correct? I look in the book and
find that the page on which the quoted lines stand is, quite rightly, p. 270.
As my eye falls on the number of the page in the book, I get the sensation
that is expressed by the words ``quite right''; the quality of familiarity
mentioned several times makes itself felt. Yet in the short interval I have
not expressly held fast the number 270; on the contrary, various ideas from
quite other regions passed quickly through consciousness. Nor, when my eye
fell on the printed number 270, did I come to think of the number I had
written. The printed number was recognized immediately, and if I had not seen
it, the idea of the written number would have continued to be latent (unless
it should be drawn forth by some other cause, e.g. by the wish to verify it).

With respect to this relation one finds in almost all psychological authors a
misunderstanding, which is connected with a confusion of the
\emph{psychological explanation} of ideational activity with its
\emph{logical significance}. The idea that makes another idea familiar and
intelligible to us is usually regarded as the active one; the one that finds
its determination or characterization is regarded as the passive one. In
Herbart's usage: the former is called the \emph{apperceiving} idea, the latter
the \emph{apperceived}. When I see an object before me and say: ``that is a
fish!'' then the idea ``fish'' is called the apperceiving idea. But it is
nevertheless in reality the momentary group of sensations that sets the whole
process going; it is this that arouses the general idea ``fish''; it is thus
this that is the active party in the process that takes place. And so the
relation must also be conceived in every recognition (in so far as we can here
theoretically distinguish between two elements). It goes without saying that
absolute activity belongs to neither of the terms; the significance of the
given group of sensations is only to \emph{release} the reproduction, to set
free accumulated energy. But that usage does not seem to be a happy one, since
it starts from the logical yield, and not from the psychical process itself.
% --- p. 25 ---
\setcounter{footnote}{0}%
\opage{25}Only in thinking proper, bound up with a conscious search for a
logical yield, do we begin with the ``apperceiving'' (interpreting and
measuring) element, since we then have a provisional idea of what we are
seeking, or of what demands the thing sought is to satisfy.

I add further what has, incidentally, been sufficiently emphasized above (4),
that the theory of immediate recognition really makes no use at all of
``unconscious ideas''. What I express \emph{theoretically} as a fusion of a
sensation and an idea is the change a sensation can undergo through
repetition. It is the effect of practice, and nothing else. ---
\emph{Secondarily}, however (8), the whole recognition can become an --- in
any case as good as --- unconscious act.

It still remains to discuss the attempts that have been made, partly by the
use of pathological cases, partly by the experimental route, to disprove the
existence of immediate recognition.

12. Dr. \emph{Arne Löchen}, in his already mentioned work, has sought to
disprove the assumption of immediate recognition through an analysis of the
remarkable pathological cases called word blindness and word deafness. These
are cases in which the patient does indeed \emph{perceive} written or heard
words as mere phenomena of sense, but does not \emph{understand} them, i.e.
does not perceive them as words or signs of ideas. On the question of how such
cases are to be explained, two views stand opposed to each other. The one
(Kussmaul, Charcot and others) holds that recognition has dropped out, while
the mere sensation of sense is unimpaired (in the expression used earlier:
\textit{A} is perceived but can no longer be combined with \textit{a}). The
other (Stricker, Lichtheim) holds that what is lacking is the possibility of
arousing such associated ideas from other domains of sense as those by which
the meaning of the words is understood. For on this view a visible word is
understood by its arousing an idea of the corresponding sound or of the
movement of the vocal muscles necessary for its pronunciation.

Against the first view there speaks, according to Löchen, in the first place
the circumstance that the word-blind seem to be able to copy. For this ability
proves --- he thinks --- that the visual ideas (the optical memory images)
have suffered nothing. --- This objection seems to me to carry very little
weight when one remembers the secondary unconsciousness with which practiced
actions can be performed. In practiced copyists there hardly takes place any
longer any conscious recognition of the written characters; the letters go,
so to speak, straight from the eyes down into the fingers. This acquired and
practiced skill thus by no means needs to drop out because the ability to
recognize has dropped out. Through practice a shortcut has been established,
which may very well be passable because the
% --- p. 26 ---
\setcounter{footnote}{0}%
\opage{26}original detour is impassable. --- Dr. Löchen himself, indeed, also
admits (p. 57) the possibility ``that the letters may seem strange, even if
one can copy them''.

There is in the second place --- according to the author --- an important
circumstance that the theory of immediate recognition has overlooked, namely
that word blindness is as a rule connected with greater or less difficulty in
finding the names of objects. The correct explanation will, he thinks, be
found by starting from this fact. Word blindness is then thought to arise from
damage to the path from the sensory center of sight to the center of the
ideas of sound or the ideas of movement. The presupposition of this
explanation is that all recognition of a visible word takes place by the
visible sign arousing the idea, attached to it by association by contiguity,
of the sound of the word or of the movement we perform with the vocal muscles
when we pronounce it. On the arousal of this association rests ``the
subjective feeling of recognition'' (by which expression Löchen designates
what I have called the quality of familiarity). It is clear that if
recognition consists in nothing else, it must drop out when the connection
between sight and hearing or the movement of speech is interrupted.

But Dr. Löchen has certainly made it very easy for himself in showing that
``the subjective feeling of recognition'' can be explained in this way. In his
opinion it is to be enough that a sensation of sense arouses ideas in me, no
matter which, and then it is recognized. ``If we came upon a flower,'' he says
(p. 103), ``and this called up the name of a species that was wrong, the
calling up of the name would certainly be enough for us to declare the flower
known; if, however, someone comes along who has better knowledge and tells us
that we have mistaken the name, we of course acknowledge that our recognition
was mistaken... But \emph{the recognition} in and of itself \emph{lies
certainly} in \emph{the calling up of the associations} and not in the
confirmation of its correctness. Here, after all, the only question is
wherein the subjective feeling of recognition lies, not its objective
truth.'' This last remark I fully endorse; but I cannot approve the use the
author makes of it. I have already distinguished definitely between
recognition and its verification (10). But the mere calling up of an idea, no
matter which, is not enough for recognition. The idea thus called up may
itself, after all, sometimes be altogether strange to us, namely if I have
forgotten it to such a degree that it does not itself appear with the quality
of familiarity. The question is thus only pushed further back. (I here
disregard what I shall come to later, whether association by contiguity does
not itself always presuppose a recognition.)
--- It seems as if Dr. \emph{Löchen here parts company} with the
\emph{expectation theory}, to which he \emph{otherwise adheres}. The
expectation theory, after all, started precisely from the view that it was
not enough for recognition that an idea was aroused, but sought to show that
recognition set in only when the ideas aroused by the parts of the object
first perceived were found to agree with the sensations aroused by the parts
% --- p. 27 ---
\setcounter{footnote}{0}%
\opage{27}appearing later. According to the expectation theory an
uninterrupted verification takes place. (Error therefore seems almost to be
excluded by this theory, which confuses recognition with its verification.
Cf. 10.) --- It is in reality a third theory that comes forward here before
us, but a theory that seems to be less happy than either of the others.

Should recognition not be able to take place without any free association by
contiguity making itself felt? Should one not be able to recognize a flower
without fastening any name whatever, right or wrong, upon it? The cases
adduced in 3 seem to speak for this. I know two little girls whose names I
constantly confuse and am therefore constantly unsure of, although I never
confuse the girls themselves, but recognize each of them with certainty, and
in recognizing them I am never conscious of any helping idea whatever. ---
The author has let himself be misled by the strong urge we as a rule have to
pass over from a given sensation or idea to other ideas or to movement (among
other things, speech). The spontaneous ideational urge and urge to movement,
as already mentioned (see 10), seldom lets immediate recognition stand
isolated and closed, but leads beyond it to verifications or to new acts of
recognition. --- It is a meritorious work that Dr. Löchen has undertaken, to
submit the pathological phenomena belonging here to a psychological analysis,
and I owe his calm and objective exposition a good deal of instruction. But
occupation with the pathological phenomena must not divert attention from
the phenomena that normal life presents.

As regards the pathological cases themselves, there is reason to wonder that
the author has not sought out such phenomena as could give information about
the state of consciousness of the patient during word blindness. It was worth
investigating whether in word blindness (in word deafness the author seems to
admit it for most of the cases) a different perception of the signs does not
take place than before, whether they do not stand in a different light, with
a different quality, than when they are recognized. At least one case, in
which an intelligent patient has described his state during the affliction,
points decidedly in that direction. The well-known physician \emph{Lordat} of
Montpellier suffered for a time from word blindness and word deafness. He
described his state more closely in the following way: ``En perdant le
souvenir de la signification des mots entendus, j'avais perdu celui de leurs
signes visibles. La syntaxe avait disparu avec les mots, \emph{l'alphabet seul
m'était resté}, mais la jonction des lettres pour la formation des mots était
une étude à faire. Lorsque je voulus jeter un \emph{coup} d'œil sur le livre
que je lisais, quand ma maladie m'avait atteint, je me vis dans
l'impossibilité d'en lire le titre. Il m'a fallu épeler lentement la plupart
des mots.'' Lordat describes, says \emph{Kussmaul}\footnote[1]{Die Störungen der Sprache. Leipzig 1877. p. 175 f. ---
The emphases are mine.}, from whom I take this detail, ``with deep emotion
% --- p. 28 ---
\setcounter{footnote}{0}%
\opage{28}the happy moment when, after some sorrowful weeks, he let his gaze
wander through his library, and when \emph{unexpectedly}, in a corner, the
words on the spine of a folio, `Hippocratis opera', \emph{again shone} out to
\emph{meet} him.'' This description does not seem to agree with Dr. Löchen's
assertion (p. 119) that the word-blind person sees the letters just as a
normal person does. What Lordat complains of is that he sees letters, but not
words. And this agrees excellently with what was adduced above, that words and
sentences are recognized more easily than combinations of letters without
meaning, and that the impressions, when they thus belonged together, came to
consciousness \emph{simultaneously} (3). Lordat's description both of what he
lost through the illness and of what came forth for him again at its end shows
that it was precisely the ability for immediate recognition that he had lost.
We have here a good example of unprepared recognition (3, 5). It is not easy
to understand how Dr. Löchen (p. 119) can make out that the word-blind see the
letters in completely the same way as healthy people: it surely makes no small
difference whether I merely see them as a heap of figures, or whether they
``shine out to meet me'', as the title of Lordat's favorite book did for him in
the moment of his recovery, with the whole warmth of the quality of
familiarity. ---

It is not the case, as Dr. Löchen (p. 11) seems to assume, that according to
the theory of immediate recognition the sensation of sense always comes first,
and the memory image after an interval of time. The theory maintains precisely
that in immediate recognition we do not distinguish between sensation and
reproduction, but that these two activities come forward before us only in
their common product, the element of consciousness furnished with the quality
of familiarity. Nor is it necessary to assume that sensing and memory belong
to two different centers of the brain. It seems to be so; but the theory can
stand without it.

--- In taking leave of Dr. Löchen, I must further mention that, in contrast
to other authors who have expressed themselves on this matter with great
self-confidence, he declares that ``at the present standpoint of science he
has believed he must take a somewhat more cautious position''.

13. In his already mentioned treatise (Om Genkendelse. Forsøg paa en
experimental Verifikation af Forestillingsassociationernes Teori.
Videnskabernes Selskabs Skrifter. 6. Række. II) Dr. \emph{Alfr}.
\emph{Lehmann} has believed he could show by the experimental route that there
is no immediate recognition in the often stated sense.

Such a belief must arouse some surprise already beforehand. If there are
different kinds or forms of recognition, how then will one really be able, by
experiments, which must always take place \emph{under certain definite
conditions}, to show that recognition cannot take place in another way
\emph{under other conditions}? It is, after all, a clear and simple rule for
all experimental research that its results hold only for
% --- p. 29 ---
\setcounter{footnote}{0}%
\opage{29}the conditions under which the experiments were made, and that the
warrant for extending them to \emph{all} conditions must in any case be shown
specially. One cannot without further ado infer from a hypothetical judgment
to a categorical judgment; for it might, after all, be that the conditions
were what was decisive. It is the fundamental error of Dr. Lehmann's treatise
that he has not seen the relation between the two theories that stood opposed
to each other. The one says that all recognition rests on expectations that
are aroused by association by contiguity and then confirmed; the other says
that \emph{besides} the (successive) recognition that takes place in this way,
there is also an immediate recognition. It is then clear enough that one does
not refute this latter theory, and does not furnish any proof of the former,
by showing that in certain cases, under certain definite conditions,
recognition takes place with the help of ideas aroused beforehand, --- if these
conditions were precisely such as excluded the possibility of an immediate
recognition! Such experiments prove what no one has denied, and disprove ---
nothing. --- And with respect to Dr. Lehmann's experiments I must go a step
further. For they do not even presuppose the expectation theory as a necessary
explanation. Only through a series of misunderstandings and fallacies can such
a significance be ascribed to them. I shall try to prove what I have here
asserted.

My objection does not concern Dr. Lehmann's experiments themselves. They have
their positive value apart from his unfortunate theoretical interpretation of
them. They shed light on various interesting points with respect to the
certainty and accuracy of recognition under \emph{certain} circumstances, but
what they teach about the psychological nature of recognition is very little,
and in any case not what Dr. Lehmann himself thinks. The more significance one
ascribes to experimental psychology, the more one rejoices that our knowledge
of psychical life can gain in clarity and accuracy by this route, the more
uneasy one must also be at seeing experimental research mixed up with hasty
theorizing. Psychology has not freed itself from the dominion of metaphysics
in order to be ruled by a shortsighted way of looking at things that overlooks
the multiplicity and diversity of psychical phenomena. Treatises furnished
with plates and tables easily impress uncritical readers. Not only for the
sake of the problem itself, but also for the sake of experimental psychology
itself, I have therefore undertaken a critical examination of the treatise in
question, and shall here set out what I have found.

a. Dr. Lehmann made a series of experiments in the following way. The
observers with whom the experiments were made were shown first a gray disk,
the standard disk (\textit{n}), then, after a certain interval, either a
lighter disk (\textit{l}) or a darker one (\textit{m}) or the standard disk
itself again, and the observers were then to decide and mark whether the disk
they saw last was equal to the standard disk or lighter or darker than it. ---
This very arrangement of the experiments shows that the immediate and
unprepared recognition discussed above was excluded. The recognition that
could take place under these conditions must have been
% --- p. 30 ---
\setcounter{footnote}{0}%
\opage{30}prepared in two ways. In the first place, there is so short an
interval between the appearance of the first and of the second disk that the
idea of the first must be said to be present, if not in consciousness, then
at least as a disposition that is not far from the threshold of
consciousness, and so not as much energy needs to be used for its
reproduction as if the phenomenon that was to be recognized came altogether
unexpected and unprepared. In the second place, it lies in the nature of the
case that when one sets a person up to test his ability to recognize in a
certain definite respect, his attention will involuntarily also turn in that
direction; he faces the second disk even more prepared than he is, e.g., in
the street, when he has first met one person and then happens also to meet
another, and he then involuntarily finds that this one is like that one or
different from him. The man who goes along the road and at a garden gate
unexpectedly sees a lady in whose countenance there is something familiar to
him performs the act of recognition under other conditions than the observers
in Dr. Lehmann's laboratory. The same holds of Lordat's surprise at being able
to read the title of his book. It lies in the nature of the case that one
cannot investigate immediate recognition by the experimental route. One cannot
avoid preparing it in the two ways mentioned. Immediate recognition must be
investigated by spying out in real life cases and examples that can shed light
on the phenomenon. It is the same here as with so many of the phenomena of
involuntary psychical life: experimental psychology can at most shed light on
them only indirectly, and one must proceed very cautiously in transferring
its results to this domain. What holds of self-observation in general, that
by directing attention to the psychical phenomenon it at once alters it, holds
not least of experimental psychology, which must always place and prepare its
objects of experiment (the observers) in a certain definite way, to which
involuntary psychical life need not be bound\footnote[1]{The preparation appears in a characteristic way in a remark of
Dr. Lehmann's, p. 27: ``At the moment when the recognition is [sic] to take
place, one tries by voluntary [sic] concentration of attention to recall the
earlier sensation''. Yes, that is precisely the point: one knows that one
\emph{is to} recognize; that, after all, is why one is in the laboratory, and
why one is set up before the disks, etc. As a consequence of this preparation
one therefore exerts oneself to the best of one's ability. The experiments
therefore show what one can achieve in the way of recognition by applying
``full attention'' (cf. p. 42) under certain given circumstances and with
respect to certain given objects. But they show nothing about the nature of
recognition where it is not arranged beforehand.}. It is thus clear that Dr.
Lehmann's experiments, however one may interpret their results, can prove
nothing whatever about immediate recognition, since its possibility was
excluded beforehand.

b. Dr. Lehmann starts from the view that when it is a question of the
recognition of a sensation, one must \emph{naturally} (when one does not
``start from special theoretical presuppositions'') assume that the idea
through whose comparison with the given sensation the recognition comes about
is given \emph{beforehand}. He appeals here to \emph{Wolfe}, who in his ``
% --- p. 31 ---
\setcounter{footnote}{0}%
\opage{31}Untersuchungen \emph{über} das Tongedächtniss'' (Wundt Studien. III.
p. 556) says: ``When we go more closely into the procedure in the comparison of
two tones separated by an interval, it is clear that without a memory image of
the first tone a comparison is altogether impossible''. But Dr. Lehmann, when
he read Wolfe, must himself have ``started from special theoretical
presuppositions''; otherwise, by reading two pages further in Wolfe's
treatise, he would have seen that this truly unbiased investigator expressly
says (p. 558) that it is ``as is well known'' not necessary to keep an image of
the first tone in consciousness, ``since the \emph{second tone at once calls up
an image} of the \emph{first}''. In other words: Wolfe starts from the view
opposite to the one Dr. Lehmann has him put forward as the only reasonable
one\footnote[1]{In his treatise ``Beiträge zur Theorie der sinnlichen
Aufmerksamkeit und der activen Apperception'' (Wundts Studien IV. p. 407)
N. \emph{Lange} assumes in an analogous way that in attentive contemplation of
an object it is the sensation that releases the corresponding memory image
and is thereby strengthened. --- Cf. above 5.}. Now when the second tone thus
calls up the idea of the first, this cannot happen by association by
contiguity, since the two tones have not occurred together. It can happen only
by the association by similarity that Dr. Lehmann rejects, but which here
offers the only possible explanation. --- I mention this only in order to
attach to it the remark that, even if all recognition took place in the way
Dr. Lehmann has let it take place in his laboratory, namely \emph{prepared},
the expectation theory or contiguity theory would not thereby be proved. For
the image of the disk \textit{n} does not keep itself in consciousness during
the whole intervening time until \textit{l} or \textit{m} or \textit{n} itself
comes; it is called up anew, and that it is called up can surely be due only
to the appearance of the second disk. The reproduction thus takes place
\emph{after} or upon the entry of the phenomenon that is to be recognized,
whereas the expectation theory or contiguity theory presupposes that it goes
\emph{before}. (Cf., incidentally, 11.)

c. The kind of recognition present in Lehmann's experiments now differs from
the immediate recognition that experience can acquaint us with, not only by
being prepared and half voluntary, but also by concerning such fine nuances,
such simple sensations, as can be recognized only by a special exertion of
attention. This Dr. Lehmann himself remarks. But he overlooks that between
such fine and simple phenomena, which require a given standard in order to be
recognized, and the more composite phenomena, whose perception, and so also
whose recognition, must take place successively, there lies a class of
slightly composite phenomena that can be perceived simultaneously and
recognized in an instant. (Cf. 3---4.) The further we go into the particular,
or the further we go into the intricate, the more lengthy and difficult the
act of recognition becomes. We are at once clear that it is a red color we see
on the postman's coat, but only by a more
% --- p. 32 ---
\setcounter{footnote}{0}%
\opage{32}intricate process of thought can we decide whether this red is like
or different from the red that guardsmen or royal footmen wear\footnote[1]{Cf. \emph{Spencer}: Principles of Psychology. § 115.}. Because
the recognition of guardsman's red is not immediate, the recognition of red in
general may well be so. Because the recognition of a series of letters
\emph{without} meaning is not immediate, the recognition of a series of
letters with meaning may well be so. --- It must constantly be held fast that
the dispute stands between two theories, of which the one takes all
recognition to be prepared and successive, while the other assumes (not: that
all recognition is immediate, but:) that besides that prepared and successive
recognition there is \emph{also} an immediate recognition.

d. While the two theories, as Dr. Lehmann shows, can each very well explain
several of the results that came out of the experiments, there is said, with
respect to one particular point, to be a sharp contradiction between these
results and the consequences that must be drawn from the theory of immediate
recognition. After repeated study of Dr. Lehmann's treatise I cannot see
otherwise than that he here contradicts himself in a striking way, in that in
one place he draws from his own theory precisely the same consequences that in
another place he draws from the opposing theory and finds to conflict with the
results of the experiments. I am, however, somewhat unsure here, since Dr.
Lehmann does not express himself sufficiently clearly. Even if I should not be
right that Dr. Lehmann here gets caught in a self-contradiction, it is,
incidentally, of no consequence for the matter; for I think I have
sufficiently shown that immediate recognition was excluded beforehand in the
experiments in question, and that no reasonable person could take it into his
head to hold that all recognition should be immediate.

Besides the standard disk \textit{n}, the observers in the experiments were, as
mentioned above, shown two other disks \textit{l} and \textit{m.} Each
experiment began with \textit{n} being shown, and after a small interval of
time either \textit{l} or \textit{m} or \textit{n} was then shown. \textit{n}
was thus shown far more times than \textit{l} or \textit{m.} Now Dr. Lehmann
infers: if there were immediate recognition, \textit{n} would have to be
recognized with far greater certainty than \textit{l} or \textit{m}, since it
comes back far more frequently; the theory of immediate recognition is, after
all, built on the law of practice; --- but the experiments showed that
\textit{n} was not recognized more certainly than \textit{m} and\textit{ l,}
--- and with that, then, that theory is condemned to death! The opposing
(i.e. in Dr. L.'s opinion opposing) theory, on the other hand, is said not to
come into conflict with the results of the experiments, for, it is said (p.
30), ``the recognition rests here on a comparison with a memory image that is
more or less effaced and unclear, and on which the constant repetition has
only the effect that after the lapse of a certain time it is not quite so
unclear as it would be if the practice had been less''. This is the
consequence that is drawn on p. 30 from ``the contiguity theory''. --- Quite
different is the sound of the consequence that is drawn on p. 15 from the same
theory. Here it is said, after it
% --- p. 33 ---
\setcounter{footnote}{0}%
\opage{33}has been developed how the contiguity theory explains the
recognition of successively perceived objects: ``From this exposition it is
easily understood why an \emph{idea is recognized the more easily}, \emph{the
more frequently} it has \emph{been} in \emph{consciousness}. For the memory
images must be assumed to become the more distinct, and one will therefore
also be able to judge \emph{the more certainly} that what is sensed is
congruent with the reproduced sensations''. I cannot see otherwise than that
the same consequence is here drawn from the contiguity theory as is drawn
there from what Dr. L. calls ``the similarity theory'', and if there is a
conflict with the results of the experiments, it is thus present in the case
of both theories. --- This self-contradiction I cannot get out of Dr. L.'s
treatise.

Another defect too seems to be present in the deductions on which Dr. L.
builds the interpretation of his experiments. If recognition were immediate,
then, he thinks, in experiments with two disks, the standard disk \textit{n}
and \textit{l}, \textit{n} would have to be recognized more easily than
\textit{l.} But this by no means follows. For here, after all, the choice is
not between recognition of \textit{n} and recognition of \textit{l}, but
between finding the second disk (\textit{n} or \textit{l)} equal to or
different from the first (\textit{n}). And that is another matter. It is,
after all, by no means given that the ease of recognizing \textit{n} should
increase through practice more than the ability to find a difference between
\textit{n} and \textit{l.} The perception of difference, which plays so great
a part in all our sensations, might very well grow just as much in certainty
with the lesser practice as the perception of likeness with the greater
practice. In order to perceive the difference between \textit{l} and
\textit{n} one does not need so complete and distinct a reproduction of
\textit{n} as to perceive the identity of \textit{n} and \textit{n.} It must,
moreover, be remembered that recognition and perception of difference here
take place with ``full attention''. --- This consideration can also be applied
to the experiments with all three disks; it may be relatively easier, and
therefore require less practice, to distinguish than to identify, to find
something lighter or darker than to find it equally light. ---

e. In another series of experiments Dr. L. sets out to investigate what
influence it has when the thing that is to be recognized has a definite name.
He starts from everyday usage, which designates something as known when one
can name its name. Now if it is a matter of recognizing a series of nuances,
e.g. of gray, Dr. L.'s experiments show that this recognition takes place with
greater accuracy when we have names for all the nuances. It is more accurate
with the five nuances of gray for which the language has names (white, light
gray, middle gray, dark gray, black) than with six or nine nuances. And when
Dr. L. practiced associating the first nine numerals with the nine nuances of
gray, greater certainty in recognition also showed itself than when they were
to be recognized without being marked each for itself. From this Dr. L. then
draws the conclusion that recognition most often takes place by the repeated
sensation calling up an idea of its name.

% --- p. 34 ---
\setcounter{footnote}{0}%
\opage{34}Here too I must distinguish between the experiments, interesting in
and of themselves, and the interpretation, in my eyes altogether unnatural and
unwarranted, of which they are made the object. --- Apart from what I have
already had occasion to object against the assertion that recognition is one
with the ability to name (12), I have so many objections to make against this
interpretation that it is best to number them.

1) The uncertainty of recognition must necessarily be greater when it is a
matter of six or nine nuances than when it is a matter of only five, whether
or not we have names for them. Every hypothesis must assume this, and only
long-continued practice can change anything in it.

2) It is clear that the very practicing of the connection between nuance of
color and name is a practice in recognition. It is then no wonder that
recognition becomes more certain with the named nuances than with the unnamed;
but from this it certainly does not follow that ``it is the names of the
colors that make recognition possible after a longer lapse of time''. The
naming is only a side effect; the cause both of the recognition and of the
naming is the more frequent repetition and the express attention that one
must direct to the nuance of color in order to attach the name to it. (In a
following section of this treatise I shall go more closely into this point
and show that all association by contiguity --- thus also that between
sensation and name --- presupposes a recognition of the first term.)

3) When the nuance \textit{A} has aroused the word-idea \textit{b}, then at
the moment the known name is uttered or vividly imagined, an act of
recognition takes place in the case of \textit{b}, on the occasion of which
the whole question in dispute about immediate recognition can be raised anew.
When for everyday use we so often content ourselves with being able to give a
thing a name (``give it a name and let it run'') and think that we have in this
a certain guarantee, this can simply be explained by the fact that the
quality of familiarity can be more powerful, so to speak more saturated, in
the name than in the thing itself. The same thing is at work here as
everywhere where words and concepts are confused. The security one feels with
the known name one involuntarily interprets as acquaintance with the thing,
and one does not investigate whether the known name might not be attached to
an unknown thing. It is a frequent experience that of two psychical phenomena
following one another, the subsequent one can draw as good as the whole of
attention to itself and come to determine the whole state. Yet it can also
happen that the vivid recognition of the second term (here the name) leads one
to turn attention back again to the first term (which presupposes the ability
for backward reproduction), and a \emph{more distinct} recognition of this
term will then take place than set in at the first moment. In this way that
everyday custom, of calling a thing known when its name can be named, might
then find its justification. The meaning, however, cannot be that a thing
should be known because one can blurt out a name with which one connects no
meaning. For the name, after all, arouses associations of phenomena similar
to the one before us:
% --- p. 35 ---
\setcounter{footnote}{0}%
\opage{35}the name stands --- when it is more than verbal chatter --- as the
representative of a whole series of associations by similarity. (To this point
too I shall have occasion to return in what follows.)

4) One must not confuse the greater ease in \emph{marking} the recognition
with a greater ease in recognition \emph{itself}. When the association between
nuance and name has been practiced, the ability to mark has also been
practiced, and it is then no wonder that experiments with nameable nuances
show greater certainty than those where a corresponding practicing had not
taken place. This ease in marking seems to be the most pronounced
peculiarity of this second series of Lehmann's experiments.

5) The importance of the observers' applying ``full attention'' is mentioned
in passing (p. 42). It would have been interesting to learn something more
about this attention. Was it voluntary or involuntary? And what was it
directed to? If it was directed to giving the perceived nuance the
\emph{right} name, it must be assumed to consist in a reflection on
\emph{what} it was that one saw, --- an \emph{identification}. And since (according
to pp. 39---41) everything had now been done to keep away the influence of
ideas aroused \emph{beforehand} (conscious or unconscious), might not the
immediate recognition (assimilation), contested by Dr. L. with such great
pains but with so little effect, put its head up, like the house-sprite with
the man who wanted to move away from him? --- Of attention and its relation
to recognition I have spoken in 5. ---

14. As the conclusion of this whole investigation of immediate recognition I
shall give a short \emph{systematic survey of the different kinds of
recognition} that have been discussed in the foregoing.

The different kinds of recognition are determined partly by the nature of the
\emph{object} of recognition, partly by the way in which the recognition is
\emph{prepared}. We take each principle of division by itself; it will,
however, naturally appear that certain objects can be recognized only after a
certain preparation, so that a member of the one division can coincide with a
member of the other.

A. \emph{Recognition characterized} by its \emph{object}.

1) Recognition of a relatively simple, not very composite phenomenon, whose
individual features can be perceived simultaneously: \emph{instantaneous
recognition}. (Psykologi, V B, 1. Cf. the present treatise 3---8.)

2) Recognition of more composite phenomena, which cannot be perceived
simultaneously: \emph{successive recognition}. (Psykologi V B, 2 and 4, and
C, 1---2 and C, 7 b and c. Cf. the present treatise 9---10.)

3) Recognition of elements and nuances that do not usually occur separately in
% --- p. 36 ---
\setcounter{footnote}{0}%
\opage{36}our experience, and therefore require special aids in order to be
recognized: \emph{special recognition}. (Psykologi V B, 1, p. 141. Cf. the
present treatise 13 c and e.)

B. \emph{Recognition characterized by the preparation}.

1) \emph{Unprepared recognition}. Here only a disposition to the return of the
idea is presupposed, but not necessarily any arousal of it that has taken
place in the near past. --- Unprepared instantaneous recognition is called
\emph{immediate recognition}.

2) Recognition can be \emph{prepared by voluntary}, \emph{though vain}
searching. Cf. Psykologi p. 83 f. Present treatise 3.

3) \emph{Involuntarily prepared recognition} presupposes either a) that the
idea of what is recognized has shortly before been forth in consciousness
without our having called it up ourselves, as when an \emph{expectation} has
been aroused in us (Psykologi V B, 4. Present treatise 9---10), an expectation
that can in turn be \textit{α}) \emph{conscious} or \textit{β})
\emph{unconscious}, --- or b) that ideas which each by itself are not
sufficient to make a recognition possible make one possible by their
\emph{summed effects}. (Present treatise 11.)

4) \emph{Voluntarily prepared recognition} can either a) consist merely in
one's knowing that one is to test one's ability to recognize for a certain
kind of objects (\emph{indefinite preparation}), or b) consist also in one's
knowing which idea one is to compare the entering phenomenon with
(\emph{definite preparation}). In both cases a certain voluntary attention is
present (cf. Psykologi VII A, 5. Cf. present treatise 13 a). Special
recognition presupposes as a rule a voluntary preparation. ---

The different kinds of preparation can naturally not be kept sharply apart
from one another. Thus recognition prepared by conscious expectation comes
very close indeed to voluntarily prepared recognition.

\begin{center}\large II. The Presuppositions of Association by Contiguity.\end{center}

15. Attempts have been made, as mentioned earlier, to show that all
association of ideas is really association by contiguity, i.e. comes about by
ideas that have frequently been present in consciousness simultaneously or in
immediate succession having a tendency to call each other up again. Such an
attempt at reduction will be sufficiently refuted if it can be shown that
\emph{association by contiguity itself} rests on a \emph{presupposition}
\emph{of which it does not itself contain the explanation}. --- Let \textit{A}
and \textit{B} be two sensations (or groups of sensations), and let us assume
that we have a certain number of times had \textit{A} and immediately
afterwards \textit{B.} Does it then follow without further ado that
\textit{A'}s reappearance in consciousness arouses the idea \textit{b}?

% --- p. 37 ---
\setcounter{footnote}{0}%
\opage{37}Association by contiguity is (like immediate recognition) a special
case of practice. In this case practice does not merely bring it about that
the single idea is more easily aroused; what is decisive here is that through
repetition so close a connection forms between \textit{A} and\textit{ B} that
they in reality make up one and the same process. One cannot (as is to be
developed more fully in a later section) understand the phenomenon of the
association of ideas if one does not hold fast that sensations and ideas are
not separated parts or absolutely independent elements, but psychical
activities or functions. What practice (which has here rightly been called
co-practice) brings about here is that there comes to be no leap or
intermediate term between the two activities \textit{A} and \textit{B.} When
the one of these (\textit{A}) enters through the corresponding phenomenon
being given anew, the mental activity will --- without needing any special
impulse, thus without the phenomenon \textit{B} needing to enter anew ---
continue itself further, so that we get \textit{A + b}. Here, then, the
remarkable thing has happened that what (for consciousness itself) was at
first a mere \emph{relation of succession} (that \textit{B} follows upon
\textit{A}) gives birth to a \emph{relation of causality}, in that \textit{b}
continues \textit{A} as the effect continues its cause.

But let us look a little more closely at what lies in this. If repetition and
practice have such an influence on the relation between \textit{A}
and\textit{ B,} must they not also have an influence on \textit{A} alone?
\textit{A} and \textit{B} surely cannot be subject to wholly different
conditions. That \textit{B} is reproduced we can explain only by assuming
that through frequent repetition a certain disposition or tendency has been
left behind in consciousness and brain which can be released without the
phenomenon itself needing to be given. But \textit{A}, after all, has been
repeated just as often as \textit{B}! Thus the same disposition must be
present in \textit{A'}s case as in \textit{B'}s. And should it
then be possible that the phenomenon \textit{A} could be given without this
disposition being aroused? It even seems to be an unavoidable inference that
\emph{the disposition to the idea} \textit{a} must be \emph{aroused} to
\emph{a far higher degree} \emph{when} \textit{A itself enters} \emph{than the
disposition} to \textit{b is aroused} \emph{because} an \textit{A different
from B} \textit{enters}. It is not easy to see how this consequence can be
avoided. And it is moreover confirmed by experience. For it has been shown by
experiments that homogeneous ideas combine more easily and more quickly than
non-homogeneous ones, or in other words, that the process of assimilation
goes on more quickly than the process of complication\footnote[1]{\emph{Tschisch}: Ueber die Zeitverhältnisse der Apperception
einfacher und zusammengesetzter Vorstellungen (Wundt's Studien II. Leipzig
1885. p. 625). --- Cf. already \emph{Charles Bonnet}: Essai analytique, p.
399: ``S'il faut moins de temps pour contracter les déterminations qui
constituent la simple \emph{réminiscence}, que pour contracter
\emph{l'habitude} de reproduire une \emph{suite} quelconque, c'est que la
reproduction de cette suite tient à de plus grands changements que la simple
réminiscence.''}. The connection of \textit{A} and\textit{ b} is a
complication, the connection of \textit{A} and \textit{a} an assimilation.

If the disposition to \textit{a} cannot be aroused, then the disposition to
\textit{b} cannot be
% --- p. 38 ---
\setcounter{footnote}{0}%
\opage{38}aroused either. Thus the condition for \textit{A}'s arousing
\textit{b,} is that it arouses \textit{a.} This has been expressed by saying
that all association presupposes an association by similarity. But here it
must surely be noted that \textit{A} and \textit{a} do not \emph{need} to come
forward in consciousness as independent elements. They can fuse immediately
or be assimilated, and we then get an immediate recognition as a necessary
introduction to association by contiguity. Since great weight has often been
laid on the difference between assimilation and association (cf. 7), it is
more correct to express oneself as I have done in my Psykologi (p. 180):
``\emph{Association by contiguity presupposes association by similarity}, at
\emph{least immediate recognition}''. That immediate recognition itself cannot
be explained by association by contiguity I have tried to show in the
preceding section, where it is also mentioned that recognition can take place
so quickly that it can be said to happen unconsciously. On account of the
quickness with which the first part of the process (the assimilation) as a
rule takes place in relation to the second part of the process (the
complication), it is easy to understand that it is often not noticed. It is
then, so far, no wonder that there are psychologists who deny that it takes
place at all.

16. This view is confirmed by definite experiences. For where the first term
in a relation of contiguity cannot be recognized, the second term does not
enter either. This is the counterpart of those cases in which we can know the
first term but cannot recall the second term. (See 3.) I sit in my room and
hear a rumbling in the street; at once the idea of an omnibus enters. But the
condition for this is that I recognize precisely \emph{this} peculiar
rumbling sound. If I confuse it with the sound produced by the rolling of
empty barrels, or with thunder, the idea of the omnibus does not come. If I
feel unsure, I listen attentively, and only the \emph{identification, made
with attention and clear consciousness,} of the present sensation (\textit{A})
and the idea of the sound familiar from earlier experiences (\textit{a}) then
lets the idea of the omnibus (\textit{b}) come forward with certainty.
Association by contiguity thus rests here on conscious identification.

But \emph{this conscious identification presupposes an association by
similarity}. For what should prompt me to compare the sound I now hear with
the idea of the sound I have heard before, if not precisely the likeness
between them? (Cf. 13 b.) How distinctly I become conscious of the relation of
likeness is of no consequence for the matter; here all possible degrees can be
conceived, right down to the complete and immediate fusion that we call
immediate recognition, and that we have already (7) designated as a border
form of association by similarity. --- An intermediate form between immediate
recognition and conscious identification we have where a phenomenon arouses
our attention, and where --- without an independent reproduction taking place
--- the mere attentive holding fast of
% --- p. 39 ---
% print: „han“ for „kan“ (misprint)
\setcounter{footnote}{0}%
\opage{39}the impression and absorption in it is enough for association by
contiguity to enter. So, for instance, when we want to call up a name in our
memory. We then imagine the thing as vividly as possible, turning attention
away from all other ideas; at last the name sought may then spring forth. The
water does not stream out before the auger has been drawn up, and that can
cost great effort.

Recognition can fail to come at the first term not only when there is not
sufficient disposition to the arousal of \textit{a} (sufficient preparedness)
present, but also when attention is diverted. Even if the connection between
\textit{A} and \textit{B} is ever so well practiced, \textit{A} can
nevertheless, by occurring in a context in which it is unable to arouse any
interest, remain unrecognized. So, e.g., when \textit{A} is merely a
transitional term between the two observations, each highly interesting in
itself, \textit{X} and \textit{Y. b} will in this case not be called up, although all conditions would have to be present if those were right who
assert that association by contiguity between \textit{A} and \textit{B} has
only two conditions: 1) \textit{A}'s and \textit{B'}s frequent being
together, and 2) \textit{A'}s re-entry\footnote[1]{Cf. A. \emph{Lehmann}: Om Genkendelse p. 9.}. Both conditions
are present in a case like the one described, without the association
working. The strongly emotionally toned surroundings have here made impossible
that dwelling of attention on \textit{A,} which conditions association by
contiguity in all such cases where recognition has not passed over into a
secondarily unconscious activity. The fact that a certain time and force must
be spent on the first term of association by contiguity indicates that
processes of a more composite nature take place here than might seem at first
glance.

17. Dr. Lehmann, in order to show the unreasonableness of assuming more than
the two conditions mentioned for association by contiguity, uses a comparison
that might precisely have led him to catch sight of the correct view. In
contesting that the likeness between the \textit{A} now given and the
\textit{A} given earlier is to play any part when \textit{A} by association
by contiguity calls up \textit{b,} he says (p. 9): ``The relation can be
conceived as analogous to any causal connection in the physical domain. When,
e.g., a stick of sealing wax is sufficiently near a scrap of paper, and the
stick is then rubbed, the paper flies over to the stick. This movement enters
every time two conditions are fulfilled, namely that the wax is rubbed and
that it is not too far from the paper. In a quite corresponding way one can
think that \textit{A} calls up \textit{b'}s emergence in consciousness when
the two conditions are fulfilled, that \textit{A} has earlier been together
with \textit{B,} and that \textit{A} emerges anew.'' To this I remark the
following. To the sufficient nearness of the stick of sealing wax to the paper
corresponds the sufficient repetition of \textit{A} and \textit{B} in
connection with each other. But if \textit{A} cannot be recognized, that is
precisely a sign that the repetition has not been sufficiently
% --- p. 40 ---
% print: „savne“, probably for „sanse“ (misprint)
\setcounter{footnote}{0}%
\opage{40}frequent, or that other circumstances at the moment prevent it from
making its influence felt. (Cf. 15---16.) If \textit{A} does not meet with a
disposition to the reproduction of \textit{A,} produced by sufficient
repetition, then \textit{B} will not be reproduced either. Just as little as a
rubbed stick of sealing wax on Mars attracts a piece of paper on the Earth,
just so little can I, by now seeing a man whom I saw 40 years ago, come to
think of his wife, whom he then had on his arm. And just as there is a
definite point between Mars and the piece of paper where the attraction begins
to work, so there is a definite degree of practice at which the sight of the
man will arouse the image of his wife. As there the relation of distance, so
here the relation of succession passes over into a relation of causality. But
what Dr. Lehmann has overlooked is that the presence of this definite degree
unavoidably also brings with it an act of recognition, and that what excludes
recognition therefore also excludes association by contiguity. Here the
analogy in the image adduced fails, and it therefore proves nothing.

18. The English psychologists who (after \emph{William Hamilton's}
model) have emphasized association by similarity as a necessary
condition for the entry of association by contiguity have sometimes expressed
themselves in a way that could lead to misunderstanding. They have thus said:
it is, after all, not the \textit{A,} we now sense that has been together with
\textit{B,} but an earlier \textit{A}; thus the \textit{A} now given must
reproduce the earlier \textit{A,} and it is only this \textit{A,} that calls
up \textit{b. Stuart Mill}, e.g., says (in his Note 35 in the first volume of
his edition of James Mill's ``Analysis of the human Mind''): ``My present
sensation could not remind me of those former sensations unlike itself, unless
by first reminding me of the sensations like itself, which really did coexist
with them''. A critic who wants to subject the theory he contests to the most
literal and external interpretation will here, by clinging to the words and
overlooking that all psychological designations are metaphors, have an easy
game. ``Just think: Stuart Mill holds that when I meet a man alone whom I have
earlier always seen together with his wife, I must first think of all the
earlier times I have seen him before I can come to think of his wife!'' Before
one ascribes to a man like Mill such an opinion, which, it is easy to show, is
close to absurd, one ought surely to read him accurately and not tear single
statements out of their context. For if one reads on, one finds that Mill
expressly designates that ``memory of the former sensations'' as a
presupposition of recognition, a presupposition that need not itself come
forward in consciousness. He says: ``This is \emph{implied} in what we call
recognizing a sensation as one which we have had before''.

What one misses in the English psychologists on this point is, in the first
place, a sufficiently sharp distinction between fact and theory. When
recognition is explained by reproduction, this is (see 4) because it can be
understood only on the assumption
% --- p. 41 ---
\setcounter{footnote}{0}%
\opage{41}that the same thing is at work in it as in independent and free
memory, only in a somewhat different way. But one must distinguish between
recognition and free association by similarity, even if they can serve to
shed light on each other, and even if the former can be conceived as a border
case of the latter. In the second place, the English investigators do not lay
sufficient weight on the fact that it is repetition and practice that produce
the quality of familiarity, which is, after all, the fact to be explained, and
that these causes, by continued working, must lead to the act of recognition
approaching the threshold of consciousness, indeed crossing it, and in any
case only with difficulty becoming an object of attention, since the new or
less familiar --- when no special interest is bound up with the recognition
--- more easily draws attention to itself. The more often \textit{A} has stood
as a mere introduction to \textit{B,} the more quickly it will be passed over,
and its entry can remain altogether unnoticed\footnote[1]{R. \emph{Wahle} (Beschreibung und Eintheilung der
Ideenassociationen. Vierteljahrsschrift für wissenschaftliche Philosophie IX.
p. 429) also reproaches Mill for not expressing himself definitely on whether
the reproduction of \textit{A} appears with \emph{full} clarity. For his own
part he adds: ``Es scheint mir oft, dass es in ausserordentlich
schwacher Nuance doch erscheine, und dass es nur, weil 1) der Glanz des eben
präsenten, dem gleichen Elementes dasselbe verdunkelt, und weil 2) die
Aufmerksamkeit auf die eintretende sollicitirte Vorstellung abgelenkt wird,
übersehen oder schwerer entdeckt wird... Ob \emph{immer}, wage ich nicht zu
behaupten. Und vielleicht ist es nicht ganz ungünstig für den Fortschritt der
Philosophie, wenn ihre Bearbeiter es wagen, nicht zu behaupten.'' These
remarks I can fully endorse from my own experience.}. It is in general a
law that where a process of consciousness can be saved, it is saved. The
functions bound up with consciousness stand in the economy of life as
detours that had to be traveled in order to reach the end; but by continued
activity and practice it becomes possible to reach the end by shorter ways.
In this way the time and force that every activity of consciousness lays
claim to are saved. The physiological time is shortened, and the energy can be
concentrated on new tasks. And even if the threshold of consciousness is not
crossed, the transition from one element of consciousness to a similar one
nevertheless takes place more quickly than the transition from one element of
consciousness to one different from it. Assimilation takes place more quickly
than complication. (See 15.) When we are surprised by vivid ideas whose causes
we find, only by express investigation, in impressions we have not noticed at
all, it must be assumed that such a quick or unconscious process of
assimilation has formed the transition\footnote[2]{Cf. my Psykologi p. 181.}.

\begin{center}\large III. Free Association by Similarity.\end{center}

19. By free association by similarity I understand a connection between
independent ideas, each of which stands out clearly on its own, that comes about
only on account of their likeness or kinship. Immediate recognition can, as has
been remarked several
% --- p. 42 ---
\setcounter{footnote}{0}%
\opage{42}times, be conceived as a kind of association by similarity, but it is
not free association, because the terms do not assert themselves independently in
consciousness. In immediate recognition the ideas (or the sensation and the idea)
are treated as if they were entirely identical. It is not thereby said that they
are so. A distinction must be made between psychological identity and logical
identity. Logical identity requires that in every context into which one of the
ideas concerned enters I can put the other in its place without its making any
difference. \emph{Leibniz} has aptly expressed logical identity in the following
way: Eadem sunt, qvorum unum potest substitui alteri salva veritate\footnote[1]{Non inelegans specimen demonstrandi in abstractis. (Opera
philosophica ed. Erdmann p. 94.)}. Psychologically, all that matters is whether
the ideas can fuse, be assimilated; they are then treated as identical, although,
objectively regarded, they may be extremely different. That is why normal
perception so often passes imperceptibly over into sense illusion.

What holds for the highest degree of association by similarity, the one in which
two elements can cover each other completely psychologically, and which may
therefore be said to rest on \emph{congruence}, holds no less for the more remote
relations of likeness, which, to distinguish them from identity or congruence, I
call homogeneity and analogy. Here, by the nature of the case, there are still
more elements of difference alongside the related elements. I am therefore
accustomed to express association by similarity (in all its degrees) by the
formula
\textit{a}\textsubscript{1} \textit{ + a}\textsubscript{2}, meaning by the
different indices to denote (not the order in which the ideas arose, but) the
circumstance that besides likenesses there are always elements of difference
asserting themselves as well, which can be heeded by consciousness to a greater
or lesser degree.

In examining whether there is free association by similarity\footnote[2]{It is denied by, among others, \emph{James Mill} (Analysis of the human
Mind. Chap. 3), \emph{Lotze} (Drei Bücher der Metaphysik. Leipzig 1879. p. 526;
otherwise earlier in Mikrokosmus. 2nd ed. I. p. 243) and Kroman (Kortfattet
Tænke- og Sjælelære. 2nd ed. p. 143), without, however, any of these authors
examining the question more closely. --- Lotze seems (according to a remark in
the work cited, p. 527) to assume immediate recognition, although he does not
assume free association by similarity.}, it is homogeneity and analogy that we
shall have to keep to. For some of the investigators who deny that there is free
association by similarity acknowledge that there is immediate recognition
(assimilation). (Cf. 7.)

Association by means of homogeneity I take to occur when, e.g., a color calls up
the idea of a related color, a form the idea of a similar form, a portrait the
idea of the person it represents, the idea of a man the idea of a man of similar
character, etc. Homogeneity is \emph{qualitative likeness}. It differs from
congruence in that a coincidence of the ideas is not possible; there must
necessarily be a transition from the one idea to
% --- p. 43 ---
\setcounter{footnote}{0}%
\opage{43}the other if the likeness is to form the connection between them. (The
German language is fortunate enough here to have two simple expressions, since
``Gleichheit'' means congruence, ``Aehnlichkeit'' qualitative likeness.) In the
case of congruence no transition need occur; only when it requires special
attention and investigation to establish the congruence (the identity) do we get
a successive process. (Cf. 10. 16.) --- Association by means of analogy I take to
occur where not a single part or property, but the relation between several parts
or properties presents likeness and thereby brings about a connection of ideas.
Examples of this are found in poetic images and juxtapositions, in the metaphors
of language and certain symbols. The likeness that is at work here we may call
(to distinguish it from congruence and qualitative likeness)
\emph{likeness of relation}.

20. It will now be my task to test whether one can in any way avoid acknowledging
association by similarity as an immediate and simple psychical relation. I shall
go through the various ways by which it has been attempted, or might have been
attempted, to reduce all association by similarity to association by contiguity.
But first I must make the remark that I do not assume, and have never assumed,
that there is any case where association by similarity operates pure and
unmixed. By the very formulation of association by similarity:
\textit{a}\textsubscript{1} \textit{ + a}\textsubscript{2} I have meant to
express that alongside the relation of likeness the relation of contiguity will
always assert itself. We shall in no case be able to point out in our experience
any example where one of the two relations operated entirely alone. Association
by similarity and by contiguity, together with an intermediate form (association
between part and whole), I have therefore set up as provisional laws which on
closer consideration prove to be different sides or limiting forms of one single
fundamental law: the tendency of consciousness to reproduce the whole of its
former state (the law of totality). Although all association by contiguity
presupposes a recognition, it does not therefore lose its independent
significance. ``In every connection of ideas,'' I have said, ``two laws operate
jointly... In experience hardly any case can be pointed out where likeness and
contiguity do not operate together.''\footnote[1]{Psykologi. 2nd ed. p. 181---184. --- This, however, has not prevented
Dr. Lehmann in his treatise from addressing the criticism of \emph{pure}
association by similarity particularly against my book. He does not mention with
a single word that, as the result of my investigation, I take neither association
by similarity nor association by contiguity to be entirely original and
independent laws, but trace them both back to the law of totality. --- I here
even disregard the fact that he quotes me in such a way that his readers must
believe that I want to make all association into association by similarity.} It
is here as with all our psychological expressions and concepts: we must not
confuse our abstractions, which we have to make in order to be able to
\emph{think} about the matter from all sides, with equally sharp differences in
consciousness itself. What matters in what follows, then, is to show that there
are psychical phenomena which become unintelligible
% --- p. 44 ---
\setcounter{footnote}{0}%
\opage{44}if we do not assume that relations of likeness can be an essentially
decisive and unavoidable condition for the association of ideas.

So far as I can see, only \emph{four} ways present themselves in which a
reduction of association by similarity to association by contiguity could be
carried through.

a) One might conceive the matter thus: when two phenomena resemble each other,
this means that they have some elements (parts or properties) in common, while
others are different. Let the one be \textit{A}\textsubscript{1} \textit{ =
m+n+p,} the other \textit{A}\textsubscript{2} \textit{ = x+y+p.} Now when
\textit{A}\textsubscript{1} by means of likeness seems to arouse the idea of
\textit{A}\textsubscript{2}, does not the transition in reality take place through
the common element \textit{p,} so that we run through the series
\textit{m}+\textit{n}+\textit{p}+\textit{y}+\textit{x,} and thereby imperceptibly
come over from \textit{A}\textsubscript{1} to \textit{A}\textsubscript{2}? And
if so, it would after all be nothing but relations of contiguity that were at
work.

b) In so far as \textit{p}\textsubscript{1} (\textit{p} as element in
\textit{A}\textsubscript{1}) and \textit{p}\textsubscript{2} (\textit{p} as
element in \textit{A}\textsubscript{2}) should be so different that they could
not be identified, there being only qualitative likeness, not congruence, between
them, one could assume that an imperceptible gradation took place,
\textit{p}\textsubscript{1} imperceptibly changing its character until the shade
\textit{p}\textsubscript{2} appeared. We should then get a series of shades as
intermediate terms between \textit{p}\textsubscript{1} and \textit{p}\textsubscript{2}
\textit{,} but no leaping transition from one idea to another similar one.

c) Our ideas are as a rule closely connected with linguistic expressions, and
this connection is an association by contiguity. Related ideas are commonly
attached to the same word. Could the association between \textit{a}\textsubscript{1} and
\textit{a}\textsubscript{2} not then be explained by \textit{a}\textsubscript{1}
calling up the word-idea \textit{n} and this in turn
\textit{a}\textsubscript{2} \textit{,} so that the word forms an intermediate
term between the two ideas, and the ``likeness'' has nothing whatever to do with
it here?

d) Finally, one might think that behind the apparent association by similarity
there lies hidden a more complex process. We involuntarily compare phenomena that
are either given simultaneously, or given in immediate succession, or that appear
to our ideation by association by contiguity. The phenomena we find alike we put
together, and as a consequence of this juxtaposition, brought about by
comparison, they are then afterwards associated in our consciousness.

21. a) With qualitative likeness and likeness of relation it will prove impossible
in every case to divide the phenomena, each by itself, into constituents of which
some are identical, others different. When we find the various shades of color
more or less related, e.g.\ yellow more like orange than it is like red,
when therefore persons who know nothing whatever of optics nevertheless order
the colors in the same way as they occur in the spectrum\footnote[1]{\emph{Wundt}: Die Empfindung des Lichts und der Farben (Studien
IV). p. 345 f. --- It agrees with this that it takes longer to distinguish two
colors that are neighbors in the spectrum than such as lie farther from each
other. \emph{Cattell}: Ueber die Trägheit der Netzhaut und des Sehcentrums (Studien
III). p. 105.}, or
% --- p. 45 ---
\setcounter{footnote}{0}%
\opage{45}when a thing of a certain color arouses the idea of another thing of
a related color, no identical element can be pointed out. Yellow and orange,
or the different shades of red, have no common element, and yet they remind us
of each other. When Homer compares the sight of the Greek warriors' gleaming
weapons with the sight of a distant forest fire, the colors of the two
phenomena, the metallic luster and the blaze, are after all of different
quality, and the poet's imagination must make a leap to get from the one to
the other. When a portrait reminds me of the man it represents, this
association can take place even if the shades of color in the portrait are not
the same as in reality, and a confusion would not for a moment be possible.
When the English language calls a Sommerfugl butterfly, or when the Danish
language calls a familiar flower Smørblomst, there is hardly psychological
identity of quality, but only qualitative likeness between the butter and the
yellow insect that has given the whole group its name, or the yellow flower.
This shows itself still more clearly when we keep to properties other than
color. When language has the expression ``a leaf of paper,'' can one really
suppose that the property of thinness has formed for consciousness a
\emph{free} transitional term between the idea of the leaf on a plant and the
idea of a piece of paper? \emph{Between} these two ideas there would thus have
to lie an idea of thinness, mark you, \emph{neither} of the thinness of the
plant's leaf \emph{nor} of that of the paper, but of thinness in and for itself!
Or when an apple in front of me on the table makes me think of an apple on a
copperplate engraving of Adam and Eve, is the image of the apple in my
consciousness then supposed to give way to the idea of the apple's bare form,
which in turn was replaced by the idea of the engraving's apple, which perhaps
did not even have quite the same form as the real apple? One cannot maintain
this view without renewing the old theory of abstraction, which lets the
``common'' elements of our ideas suffice to make up new, independent ideas. And
this becomes most unreasonable of all when one notices the great role such
associations by similarity play in children. A child 14 months old called the
heating apparatus in a railway compartment ``Tiktak,'' because it reminded it of
a clock. The same child, when it was 20 months old, called a red rubber balloon
first moon, then --- after some reflection --- lamp\footnote[1]{This last example is taken from \emph{Marty}
(Vierteljahrsschrift für wissenschaftliche Philosophie III. p. 324 f.), who
uses it in another connection.}.

Likeness of form forms in a way a transition from qualitative likeness to
likeness of relation, since form rests on or expresses a relation between the
parts. When a poet's imagination passes from the sight of the dancing couples in
the ballroom to the idea of the rotation of the globes in the heavens, one will
not be able to point out any constituent that can form the \emph{independent}
transitional term. Nor when the word green is transferred to designate an
inexperienced young person. Or in the many metaphors by which material
expressions pass
% --- p. 46 ---
\setcounter{footnote}{0}%
\opage{46}over to designating mental relations and activities. E.g.\
deliberation (cf.\ also the French penser, which is
properly the same as peser). The two ideas: the balance with scales and weights
moving up and down, --- and the impulses and promptings combating one another
in a person's mind, have no common constituents. The contiguity theory must
assume that consciousness always, in abstract fashion, thinks of a tertium
comparationis as intermediate term. That is what one does when one has to
explain an image or a comparison to those who do not understand it; but for
others it is unnecessary, and above all the often wonderful imagination that
manifests itself in the metaphors of language\footnote[1]{Cf.\ on these: \emph{Whitney}: Life and Growth of Language.
London 1875. p. 80 ff. and especially: A. \emph{Darmesteter}: La vie des mots. Paris
1887. --- Even if the distinction between metaphors and metonymies is not in
every case carried through with entire certainty, it will nevertheless be an
impossibility to reduce \emph{all} metaphors to metonymies.} and in poetic
comparisons has not needed this schoolmaster's method.

It thus seems to conflict completely with psychological experience that related
ideas, each by itself, should always be divided into identical and different
elements\footnote[2]{As regards ideas of hearing I refer to
\emph{Stumpf}: Tonpsychologie. Leipzig 1883. I. p. 115 ff, who on this point
has come to the same result as I. --- That such a division can \emph{sometimes}
take place I do not deny (cf.\ 22).}. I believe I can maintain that this
assumption too is an example of the logical consideration having been confused
with the psychological. (See 11, end.) Logically we can sum up a series of
marks that make up the content of the one concept, and another series of marks
that make up the content of the other concept, and when we compare these two
series we may find that their terms are partly identical, partly different. But
logical abstractions must not be confused with real ideas. To the logical
distinctions there do not correspond just as many real ideas in consciousness.

And even if one could \emph{always} explain the likeness by the two ideas having
partly identical (\textit{p})\textit{,} partly different constituents (in the
one \textit{m+n,} in the other \textit{x+y})\textit{,} not all association would
thereby be reduced to association by contiguity. Because within the one group
(\textit{m+n+p}) I have come to \textit{p,} it by no means follows that one goes
on with \textit{y + x.} That happens (as shown above, see 15---16) only when
there is a certain memory of our having had \textit{p} before. If \textit{p}
stands before me as something entirely new, it will not happen.

22. b) Even if, then, it is not always a single constituent belonging to both
ideas that is the \emph{free} intermediate term between them, one might suppose
that a certain altered, blurred and indistinct form of the first idea could form
the intermediate term, the common constituent running through a series of
shades, by which the first idea
% --- p. 47 ---
% print: „ønker“ for „ønsker“ (misprint)
\setcounter{footnote}{0}%
\opage{47}at last proved to be transformed into the second. The apple I see in
front of me on the table is red; the apple on the engraving is gray; and I
cannot have any idea of the apple-round form without \emph{some} color or other.
But perhaps the color could be blurred and left undetermined, and this
undetermined idea then form a transitional term between the idea of the red and
the idea of the gray apple.

We certainly have undetermined and provisional ideas, as regards both quality
and quantity. \emph{Berkeley}, who with such great zeal combated the old doctrine
of abstraction in psychology and maintained the necessity that every idea must
have an individual determinacy, was disinclined to admit this. When someone
promises me ``something good,'' there arises, Berkeley thinks, no undetermined
idea of good things; but the words --- without any idea whatever arising ---
immediately arouse a feeling of joy, just as a threat immediately arouses a
feeling of fear\footnote[1]{Principles of Knowledge. Introduction. § 20.}. Yet
surely, before one has become entirely used to the sound and meaning of the
words, an idea will arise and not merely a feeling. At the promise one will
think of something that has likeness to what one most wishes for, and at the
threat of something like what one most fears. In any case there will be many
intermediate forms between the clearly conscious idea of the meaning of the
promise or the threat and the purely instinctive emotion in which the promising
or threatening word acts only as a releasing impulse, without any idea attaching
to it\footnote[2]{In his interesting treatise ``On some Omissions of Introspective
Psychology'' (Mind. 1884) W. \emph{James} refers to the provisional ideas we
form before we have yet got an answer to a question or have reached a definite
view of a matter.} (cf.\ 8).

If we now consider such undetermined and provisional ideas more closely, they
must be divided into two classes, according as the undetermined property is of
an extensive or of a qualitative kind. Where it is a matter of relations of
magnitude, I can have an idea of a smaller extent and alter it by simple
enlargement, or an idea of a larger extent and alter it by simple cutting away.
In both cases I can distinguish between common and different parts, and the
indeterminacy consists merely in my not knowing where I am to set the limit.
Where, on the other hand, it is a matter of qualities, I can indeed also choose a
provisional quality, but if it is not the right one, or if doubt sets in, it
must by a new choice be exchanged for another quality, without there being any
question of common constituents. I cannot have any idea of color except by
provisionally thinking of a certain definite color, which then later perhaps
(when the answer comes) must be altered by smaller or greater leaps on the scale
of colors.

A dark object I see in the distance may be a lifeless object, an animal or a
human being. It is an undetermined idea I have of it, but I really have an
% --- p. 48 ---
\setcounter{footnote}{0}%
\opage{48}idea. Now when it comes nearer, its size is altered by enlargement; it
becomes more than a black dot. But the other properties which make me recognize
it as an animal or a human being I cannot add in this way; they come in through
qualitative changes.

As for these qualitative changes themselves, they might be thought to take place
continuously. It would then properly not be a phenomenon of association at all,
any more than the varying of after-images is association. The red color of the
apple might be thought to become less and less saturated until we came to the
gray color. If one wanted to maintain that an association took place none the
less, the idea of one shade of color calling up the idea of a less saturated
shade, etc., then this would (at least in those who were not familiar with the
scale of shades) not be an association by contiguity, but an association by means
of qualitative likeness. There is no reason to assume that shades or forms which
have great likeness to each other should for that reason occur frequently
together.

If the changes do not take place continuously, it cannot be an association by
contiguity that takes place either. At each new transition of quality it must be
a related quality that replaces the one ideated. Kinship is the only motive that
could be thought of for such a qualitative varying of the ideas.

23. c) The way by which it is most frequently attempted to reduce association by
similarity to association by contiguity is surely the third of those stated
above. It is held that the word, the common designation, serves as intermediate
term between two related ideas. This transition even seems to lie so near at
hand for some that when, e.g., they find the transition from the idea of
Napoleon to the idea of Alexander given as an example of association by
similarity, they assume without further ado that the word general has formed
the intermediate term. Indeed, one is so far captive to the assumption of the
predominance of word-ideas that one begins one's construction of the
psychological process which is assumed to have taken place with the \emph{name}
Napoleon. One supposes that this name reproduces the idea of a great general who
has won many battles and conquered many lands, and this idea then in turn
reproduces the idea of Alexander\footnote[1]{A. \emph{Lehmann}: Om Genkendelse p. 8.}. It might be of interest
here to learn whether ``idea of a great general'' means idea of a great general
in all \emph{generality} or merely idea of the great general that Napoleon was.
In the first case one would be returning to the old theory of abstraction and
ascribing to us the ability to have real abstract ideas, whereas most people
nowadays would surely agree that when they were to think of a general, they must
always
% --- p. 49 ---
\setcounter{footnote}{0}%
\opage{49}imagine a definite (historical or invented) person. If we now choose
some person or other (different from Napoleon) as psychological representative
of the concept ``great general,'' then the transition from the idea of Napoleon
to the idea of this representative is, after all, of exactly the same kind as
the transition from Napoleon to Alexander, and nothing is explained. It makes
(according to 22) no difference how undetermined the intermediate term is
thought to be. --- In the second case --- if the idea of a great general were
individualized by Napoleon himself --- we should, after all, not have got
anywhere. --- By neither of the two ways, then, can the explanation be reached.
As regards the example used, which is taken from my Psykologi, I can testify
that I relied on a self-observation in which the word-idea played no role. I am
accustomed to imagine historical personalities in concrete shape (Napoleon,
e.g., in the well-known green uniform), and indistinct images of actions and
events attach themselves to them. The transition took place from one image to
another.

There are great individual differences with regard to how great a role
word-ideas play in the ideational life of different individuals. There are those
whose memory is essentially a memory of things, and those whose memory is
essentially a memory of words\footnote[1]{Themistocles rerum, Hortensius verborum recordatione dicuntur
valuisse. L. \emph{Vives}: De anima. Brugis 1538. p. 56. --- Themistocles was a
practical statesman, Hortensius was an orator!}. And in one and the same
individual word-ideas play a different role under different circumstances. When
one lives much with books and has much literary work, or when one talks much
with other people, word-ideas will naturally become frequent and prominent and
will attach themselves as constituents to most other ideas. But the more one,
e.g.\ in a longer vacation, frees oneself from these occupations, the more the
word-ideas recede, and the images of memory and imagination move immediately
among one another. A careful and independent observer, the pathologist
\emph{Stricker}, who himself lays great weight on word-ideas, expresses himself
on this as follows: ``When I work much in a literary direction, when I shape
treatises or lectures, it costs me effort immediately after the work to be
without ideas of words or tones even for just a few minutes. On the other hand,
immediately after a walk in the open air, or after having looked at works of
art, I easily succeed in letting myself be governed by the memories of the
figures I have seen. Just as easily, immediately after a pleasant bath or after
other pleasant sense impressions, I can for some minutes give myself over
entirely to the memory of \emph{these} impressions, without ideating a single
word or a single tone''\footnote[2]{Studien über die Sprachvorstellungen. Wien 1882. p. 2.}. Artistic and
poetic imagination moves for the most part in images, not in words, and the
more so the stronger the movement it is in. It may even be with a certain
resignation that the great poet resolves to clothe his ideas in the poorer,
abstract form of the word. \emph{Goethe} thus says (in a conversation with
Eckermann) of his ballads: ``I had
% --- p. 50 ---
\setcounter{footnote}{0}%
\opage{50}already had them in my head for several years as lovely images,
beautiful dreams... I resolved reluctantly to bid them farewell by giving them
a body in the inadequate, paltry word''!

Nor is it the case for all ideas that the association with the corresponding
words is equally firm and unavoidable. Words that designate individual
phenomena are forgotten more easily than designations for abstract relations.
A highly deficient memory for names need not be a pathological phenomenon, but
may be connected with one's having little use for words, keeping essentially to
the concrete phenomena. The easier forgetting of concrete names has surely been
rightly explained by the fact that with concrete phenomena we do not need the
word, since the image of memory or of the imagination is enough for us, whereas
abstract relations are more difficult to hold fast except with the help of the
verbal sign. It follows from this that names of abstractions are far better
practiced than names of concrete phenomena\footnote[1]{\emph{Kussmaul}: Die Störungen der Sprache. Leipzig 1877. p. 164.}.

According to the now common view of the origin and development of
language\footnote[2]{\emph{Madvig}: Om Sprogets Væsen, Udvikling og Liv. Kbhvn. 1843.
p. 5. \emph{Whitney}: Life and Growth of Language. London 1875. p 278---309.}
it is owed essentially to the urge to communicate. Only after it has been
developed as a means of communication does it become, secondarily, for the
individual a means of holding fast and distinguishing ideas. It is wrongly that
older thinkers in particular reversed the relation and let the word be a mark
(nota) for the individual in his consciousness before it became a sign (signum)
by which he made himself understood to others\footnote[3]{Thus \emph{Hobbes}: Logica. II, 3.}. We can --- as both
philologists and pathologists insist\footnote[4]{A. \emph{Darmesteter}: La vie des mots. p. 37. ---
\emph{Kussmaul}: Die Störungen der Sprache. p. 17.} --- very well have ideas
without words. Above, in another connection (12), use has already been made of
the fact that we can recognize a phenomenon without remembering its name; what
holds for recognition (bound memory) holds also for free memory. The development
of the meanings of words shows clearly enough, moreover, that ideational
activity and the connection of ideas can hurry on ahead of the use of words.
When a word is transferred from being the designation of one phenomenon to being
the designation of a similar phenomenon, this transition presupposes a
connection between the ideas of the two phenomena, a connection that can only
have been brought about by the likeness between them. The word general was
formed to designate such men as Napoleon and Alexander; even if one wanted to
say that it is the word-idea general that now ties the idea of Napoleon and the
idea of Alexander together, this association by contiguity nevertheless
presupposes a preceding, primary activity of consciousness by which Napoleon and
Alexander (or similar men) were put together for the first time.
% --- p. 51 ---
\setcounter{footnote}{0}%
\opage{51}Where else should the common word come from? --- In scientific
classification one seeks to arrange the names in such a way that they point to
the place the phenomenon occupies in relation to other phenomena according to
its likeness or its difference with respect to them. What in science happens
with full consciousness and on the basis of thorough analysis is done at an
earlier stage of development, in its rough outlines, involuntarily on the basis
of associations by similarity.

It is by no means to be denied, however, that word-ideas --- whatever their more
precise psychological character may be\footnote[1]{Here too individual differences assert themselves. In
some (most) it is the sounding word, in others the visible word, in others again
the word as an idea of movement that predominates. \emph{Stricker} has in his ``Studien
über die Sprachvorstellungen'' laid so much weight on the purely motor side of
word-ideas that a critic in ``Revue philosophique'' (1886. II. p.
335) has supposed an abnormality in him: ``Le cas de M. Stricker est peut-être
pathologique. II y a peut-être des sourds chez ses ascendants. Son procédé
intellectuel est-il un fait d'atavisme?'' Instead of asking who is normal and
who is abnormal, it is more correct to acknowledge the individual differences.
--- Cf.\ on this question: \emph{Stumpf}: Tonpsychologie.
I. p. 153 ff. \emph{Ribot}: Psychologie de l'Attention. Paris 1889. p. 84 ff.}
--- play an extraordinarily great role. The word becomes a means of
communication because it is itself from the very first a natural outlet of the
effect the phenomenon makes on us: sound or gesture are consequences of the
emotion our experiences bring with them. Sense impression, feeling of pleasure
or pain, and audible or visible outburst: that is the primitive series of
psychical elements that every phenomenon of significance for the individual's
existence will call up. This series is gradually shortened, since the outburst
can be held back when it is not to be used precisely as a means of
communication. Our bare sensations and ideas are artificial products in so far
as they were originally only transitional terms to outburst (or action). A
process of inhibition is the condition for sensations and ideas being able to
become the main thing in the individual's relation to the objects it confronts.
But although the primitive connection between idea and outburst\footnote[2]{Cf.\ Psykologi. p. 179 and 197.} is thus inhibited, it is
strengthened on the other hand by the urge to communicate, which brings it about
that word-ideas are important elements in a great many of our connections of
ideas, and by the urge to fix the ideas of the common properties of things. Very
often, therefore, besides the association by similarity proceeding from the idea
of the thing, an association by contiguity proceeding from the name of the thing
will take place. As mentioned above, there is hardly any pure association by
similarity; rather, both motives work together. They may lead in the same, but
also in different directions. When the bare association of words gains the upper
hand, and is not controlled by association by similarity, it will be a sign of
thoughtlessness, if not of mental illness.

It may be difficult in individual cases to decide how much is due to association
by similarity proceeding from the intuitive principal idea, and how much is due
to
% --- p. 52 ---
\setcounter{footnote}{0}%
\opage{52}association by contiguity proceeding from the simpler, more mechanical
word-idea. --- During a visit to a small provincial town I remarked one day that
one seldom saw anyone go into a shop. The next day I sat several miles away on
the coast and watched the gulls fishing. I then remarked to my companion that one
almost never saw any of them catch a fish. At the same moment I came to think of
the small town's empty streets and unvisited shops, and the characteristic
juxtaposition of the two images (the gulls that got no fish, and the merchants
that got no customers) struck me at once, although it was vacation time and I was
not particularly occupied with psychological theories. --- I am sure that the
analogy between the two phenomena played the greatest role in the association,
but cannot deny that I may have used the same words on both occasions. The words,
so commonplace in themselves, would not, however, be sufficient to explain the
association; there can only be a question here of a contribution. ---

Since there in fact are ideas that need not be attached to words, --- since the
use of the same word for related phenomena presupposes association by similarity
between the corresponding ideas, --- and since the association of words is
evidently often only a contributing factor, --- it will be impossible to explain
\emph{all} associations by similarity as associations by contiguity that come
about with the help of the word. ---

24. d) As regards, finally, the fourth hypothesis\footnote[1]{It has been put forward by, among others, \emph{Victor Brochard} in his treatise:
La loi de similarité (Revue philosophique. 1880. I. p. 252 f.) and by
\emph{James Ward} in the article ``Psychology'' in the Encyclopaedia Britannica.
9\textsuperscript{th} ed. Vol. XX. p. 60 b.}, it could lead to an investigation
of the relation between the association of ideas and comparison. This
investigation, however, I shall postpone to the next section and here only remark
provisionally that there is hardly reason to assume any difference in principle
between the activity of consciousness that manifests itself as association of
ideas and that which manifests itself as free comparison. On the other hand, I
shall dwell here on another point, which seems to be decisive for this hypothesis
and which will at the same time give me occasion to go somewhat more closely into
the psychology and physiology of association by similarity.

For the question has been raised: what does likeness between a present and a
non-present idea properly mean? How can likeness cause an idea to emerge when
consciousness has not beforehand compared this idea with the present one which is
assumed to call it up?

It is certainly a remarkable property of consciousness that we are confronting
here; but is it properly more remarkable than the one that shows itself in
association by contiguity and in all memory? It is the peculiarity of the higher
forms of life that they are independent of the single moment and the single place
and can exercise functions that point far out
% --- p. 53 ---
\setcounter{footnote}{0}%
\opage{53}beyond the limitation of time and space. In all memory and in all
activity of thought, in every instinct and in every impulse, this peculiarity,
or if one will, this ideality, shows itself. In association by contiguity
consciousness connects what is different in character, although it belongs
together in time and space. This ``belonging together'' in time and space does
not, however, abolish the difference between the ideas; and it is only a fancy
when one thinks one can better understand how \textit{a} can arouse
\textit{b} than how \textit{a}\textsubscript{1} can arouse
\textit{a}\textsubscript{2} \textit{.} In association by similarity there is
difference in time and space, but the character is identical, homogeneous or
analogous; what is put together here is certainly something that has not been
put together in consciousness before; but should not the likeness of the content
of the ideas be able to weigh up against nearness in time and space? Just as well
as one asks how consciousness can ideate something merely because it is like
something else, when it has not compared the two beforehand, just as well one
might ask how consciousness can ideate something merely because it formerly
occurred simultaneously with or immediately after something else, when it does
not first remember this former occurrence. There are, after all, a great many
cases where we are surprised at our own associations by contiguity and only on
closer investigation become clear that they are reproductions of something we
have formerly experienced; often one must even resort to waking or real dreams to
find the explanation. --- It will in any case be worth while to test whether one
could not show that it is just as ``natural'' to pass from an idea to a similar
one as to pass from an idea to one different from it that has occurred together
with it.

The life and range of the association of ideas depend on the energy at the
disposal of consciousness and the brain. There may be so little energy present
that the process of association as good as stops, or that at any rate
associations that would otherwise be ready to hand do not take place. Through
overexertion and exhaustion one may lose knowledge one otherwise has at one's
disposal. On the other hand, the association of ideas gathers speed when there
is particularly much energy available, and when it is particularly easily
released. In states of fever, under the influence of opium and in a strongly
agitated mood, associations of ideas and acts of recognition become possible
that could not otherwise take place. In physiological respects one must here
recall the decisive significance that the flow of blood and its fluctuations have
for the activity of the brain. According to Maudsley a fifth of the whole body's
quantity of blood flows through the brain, and every decrease or increase at once
has an influence on the activity of the brain.

If there is ample energy present, there will be an urge to use it. The spring is
strongly tensed, and the very slightest influence is enough to make it strike
out. Just as animals and human beings in a healthy and vigorous state
involuntarily move merely to find an outlet for their accumulated force, so that
even if there is an external occasion for
% --- p. 54 ---
\setcounter{footnote}{0}%
\opage{54}movement, this external cause is at any rate entirely vanishing in
relation to the internal cause, so there is also often an urge to ideational
activity which is not directed at any definite end, but which physiologically is
likewise connected with the urge to discharge the energy stored up in the brain.
One can speak of spontaneous ideational activity, just as one speaks of
spontaneous movements. And just as spontaneous movement naturally strikes into
the directions and paths the individual is used to going, so too, as regards
spontaneous ideational activity, there will be certain series of ideas that are
reproduced particularly easily; it is the accustomed circle of ideas that
involuntarily presents itself when, after sleep, we again orient ourselves with
such great speed in the world of waking reality. \emph{Spontaneous activity goes}
in the \emph{direction} \emph{where the least resistance is met}, \emph{and where
the discharge therefore goes} on most easily. Association by contiguity finds its
simple explanation in this. But in this same principle lies also the explanation
of association by similarity. For if a certain
idea or group of ideas is given, and if it is embraced with strong interest and
thus sought to be held fast as long as possible, while at the same time the
strong movement in consciousness and brain seeks an outlet, then two tasks seem
to be set here, or two needs to be nourished, which conflict with each other:
interest in what is ideated leads to holding it above the threshold of
consciousness and becoming absorbed in it as much as possible, but the strong
movement leads to passing over to other ideas, to which is added that (as already
shown above, see 5.) attention cannot operate entirely continuously. The point
is: at one and the same time to dwell as much as possible on, or at any rate to
depart as little as possible from, the given idea, and yet to give satisfaction
to the urge to ideational activity\footnote[1]{To make the matter more evident I have supposed a case of
particularly ample energy of consciousness or of the brain. But under ordinary
circumstances too there is so much energy present, and it is released so
constantly and unavoidably, that we are not able to stop the stream of ideas in
us, but must content ourselves with forming, by exertion of attention, centers of
association around which the stream can circle. Already \emph{Locke} remarks: I would have any
one try whether he can keep one unvaried single idea in his mind, without any
other, for any considerable time together..... He will find it difficult to
keep all other ideas out of his mind, but that some, \emph{either of another
kind}, \emph{or various considerations of that idea} (each of which
considerations is a new idea) will constantly succeed one onother (Essay, II,
14, 13---14).}. In association by similarity both demands are satisfied, and the
apparently contradictory task is solved. Consciousness here circles around a
given idea, leaving it only to pass over to homogeneous (qualitatively alike) or
analogous (alike in relation) ideas. The urge to change is satisfied; \emph{other}
ideas come in place of the given one; but the change, on account of the
homogeneity or the analogy, becomes the \emph{least} possible. Here too, then, it
holds that the urge to activity leads in the direction in which there is least
resistance to overcome.

% --- p. 55 ---
\setcounter{footnote}{0}%
\opage{55}One will not be able to understand the working of the poetic
imagination without taking a way of looking at things like the one put forward
as one's basis. The poet's heartfelt and strong sympathy leads him, as epic and
dramatic poet, to imagine the object carefully and faithfully in all its details
and with all its features. But it is the same urge that leads him, as lyric
poet, to express his experience and his state by seeking parallels and analogies
to it. He cannot express his whole mood merely by dwelling on the idea of the
object itself; and on the other hand he will not let go of this idea. Then the
expedient is seized of seeking other ideas that present similar properties. He
regards, so to speak, the whole of existence from the point of view that the
idea attached to his mood gives, and finds a reflection of it everywhere. In the
Homeric similes one finds the satisfaction of this urge to let the peculiarity
of the ideated thing come fully into its own --- precisely by imagining another
thing. The urge to symbols arises in part from this cause. (Only in part, for
many symbols have formed purely historically, thus by way of association by
contiguity.) --- Let me take a single example. It says in an old Little Russian
poem\footnote[1]{Quoted in G. \emph{Brandes}: Indtryk fra Rusland. p. 277.}:

A hop vine alone in the garden
% En Humleranke alene i Haven

twined along the ground.
% slynged sig langs med Jorden.

A young little girl among the people
% En ung lille Pige blandt Menneskene

wept most bitterly.
% græd ganske bitterligt.

O green blossoming hop, say,
% O grønne blomstrende Humle, sig,

why do you not climb upward?
% hvorfor ranker du dig ikke opad?

O dear young girl, why
% O kære unge Pige, hvorfor

do you grieve over your fate?
% sørger du over din Skæbne?

The weeping girl and the vine creeping along the ground are tied together because
the mood needs more than one idea to find vent, and because the first idea is
nevertheless not wholly abandoned when an analogous idea steps into its place.
The irrepressibility of the grief stands out more clearly when it is put
together with the impossibility for the hop of climbing upward without support.
The same urge is at work here as that which, in a strongly agitated mood, leads
us to repeat the same word several times. Wonder, inwardness and the feeling of
the importance of the matter find vent in this way. (A great, great man. I was
very, very glad. Verily, verily I say unto you.) This is connected with the slow
development of feeling and with its tendency to spread in consciousness\footnote[2]{Cf.\ my ``Formel Logik''. 2nd ed. p. 27 f. and the sections of my
Psykologi cited there.}. Association by similarity too is a kind of repetition,
a repetition with variation, but a variation that does not compel a change of
the original mood.

% --- p. 56 ---
\setcounter{footnote}{0}%
\opage{56}In thinking proper, too, association by similarity, the transition to
related ideas, plays an essential role. In treating a problem the thinker
starts from a provisional idea and seeks to illuminate it by as many parallels
as possible. A great discovery often rests on an ingenious juxtaposition that
comes about only by means of association by similarity. Prior to the voluntary
and methodical working through and testing of the connection of ideas there
goes an involuntary ideational activity that supplies the possibilities which
the conscious methods can test. Without such an involuntary process as basis,
knowledge would not get anywhere.

The history of the sciences shows us, e.g., how an involuntary classification
leads to ``naming natural objects that are extremely \emph{similar}, but stand
in \emph{no visible connection}, with the same name,'' and how the first step
toward a scientific investigation of the relations of phenomena was thereby
taken\footnote[1]{Cf.\ R. \emph{Körner}: Die logischen Grundlagen der Systematik
der Organismen (Wundt's Studien. II). p. 208.}. At the same time scientific
concept formation, by which homogeneous and analogous phenomena are comprehended
under one representative idea, points to the same motive that is at work in all
association by similarity: the tendency to alter the present state as little as
possible, to hold fast the given point of view, even though there is a passing
over to new ideas. This tendency is satisfied now by our adapting our earlier
ideas to the new experiences (instead of forming entirely new ideas), now by our
conceiving the new experiences, as far as possible, in the light of our already
given ideas. The continuity in the development of knowledge rests on this
\emph{psychological} ``law of parsimony,'' which, by the way, does not always
satisfy what the \emph{methodological} law of parsimony demands, for the
phenomena can often be conceived and explained in a simpler way than we, led by
the urge to preserve the continuity of consciousness, are for the time being able
to\footnote[2]{Cf.\ on this E. \emph{Mach}: Über Umbildung und Anpassung im
naturwissenschaftlichen Denken. Wien 1884. --- \emph{Avenarius} has in a highly
interesting way carried through the thought put forward above in his treatise:
``Philosophie als Denken der Welt gemäss dem Princip des kleinsten
Kraftmasses''. Leipzig 1876. --- Already \emph{Locke} remarked that we form
general concepts and propositions and shape them in words only in order to ease
our memory and save time both when we think and when we want to communicate with
others. Essay III, 6, 30; IV, 12, 3.}. ---

Naturally neither the poet nor the thinker can, by association by similarity,
produce ideas they have never had before. Preparatory and supplementary
processes play an important role here, no less than in immediate recognition
(11). All association is a reproduction; but what is peculiar to association by
similarity is that in it an idea is reproduced by another with which it has not
previously been in connection. Herein lies what is new, what is ``creative.''
The possibility of this creative
% --- p. 57 ---
\setcounter{footnote}{0}%
\opage{57}activity is, as I have tried to show, given in the tendency to
preserve the continuity of consciousness with as few variations as possible.

Since in the foregoing strong weight has been laid on the influence of feeling, I
will merely remark quite briefly that it would not be correct \emph{always} to
regard the element of feeling as the intermediate term between the two ideas that
seek to associate themselves by means of likeness. In some cases it is no doubt
so, e.g.\ in ``analogies of sensation,'' in which ideas from different domains of
sense are associated, which only seems explicable by their each having a
tendency to arouse the same feeling. When we put a high tone together with an
active color (red, orange or yellow), this can only come from the related effect
on feeling\footnote[1]{Psykologi p. 334. --- The possibility cannot, however, be said to be
entirely excluded that there could be a real likeness present between sensations
from different domains of sense. \emph{Fick}: Die Lehre von der
Lichtempfindung (Hermann's physiol. Handbuch, III, 1) p. 166 maintains this in
opposition to Helmholtz's sharp distinction between the different
``modalities.''}. In the phenomena I mentioned above, on the other hand, there is
(apart from the element of feeling) something homogeneous and analogous, and the
element of feeling only contributes in so far as it drives toward seeking related
ideas. The one idea (e.g.\ of the hop creeping along the ground) would not in and
of itself arouse precisely such a mood as the one present in the example; it is
precisely only through its relation of kinship to the other idea that it can
serve to express the mood. --- One could no doubt point out cases in which both
--- the kinship of the content of the ideas and that of the effect on feeling ---
worked together.

It has been thought that a quite peculiar difficulty was bound up with giving a
physiological explanation of association by similarity. Whereas for association
by contiguity one refers to the analogy with a series of practiced movements
which on repetition involuntarily follow in the same order as before, it is held
that for association by similarity no corresponding process can be pointed out.
The matter seems, however, --- after what has been developed above --- to present
no difficulty in principle. If we suppose that to every idea there corresponds in
one way or another certain vibrations in the small parts of the brain (or of a
certain part of the brain), then the brain, when there is in consciousness the
tendency to association by similarity described above, will have a tendency to
carry out a series of vibrations that resemble those corresponding to the idea
that is the starting point of the association by similarity. The same form of
vibration or direction of vibration (however one may now think of the matter) is
held fast with certain changes. One might perhaps think of the matter in analogy
with the series of waves of essentially the same form and direction that
successively form when a stone has been thrown into the water. --- This mode of
explanation is by no means more hypothetical than the corresponding explanation
of association by contiguity in terms of the physiology of the brain.

% --- p. 58 ---
\setcounter{footnote}{0}%
\opage{58}25. If all association is an expression of the energy of consciousness,
we must also expect to find again in association the general fundamental form of
the life of consciousness. We must be able to point out in the process of
association the comprehending, unifying activity which is the mark of
consciousness, and whose traces can be found in every phenomenon of
consciousness. (Cf.\ Psykologi, II, 5; V B, 5.) And so indeed it is. Association
is, after all, as the word itself shows, a connecting activity. What happens in
every association is properly this: starting from the single given element of
consciousness, consciousness constructs, according to the cooperating principles
of likeness and contiguity, a larger connection in which the given element of
consciousness makes up a single term. \emph{Association} between part and
\emph{whole} therefore becomes the fundamental law of the association of ideas
(see Psykologi, V B, 8 c). The less different the parts are from one another, the
more we approach a mere association by similarity; the more different, the more
we approach a mere association by contiguity. But in reality both tendencies
always assert themselves in union. The urge to activity finds vent both in the
supplementing of what is given fragmentarily and in movement in circles around
it. Both through the relation of likeness and through the relation of contiguity
the single element gets its place in the whole psychical order of the world.

Not only from the formal side, but also from the real side, we find in the law of
totality the proper fundamental law of the association of ideas. As regards the
\emph{content} of consciousness too a concentration takes place, all ideas being
as far as possible brought into relation with, and ordered according to their
kinship with, the idea or ideas that are most intimately attached to the
prevailing feeling and direction of will in the individual.

It is the Scot \emph{William Hamilton} who (after earlier hints in Kant and Fries)
has paved the way for this reduction of the empirical laws of association to one
fundamental law.

When we speak of association by similarity and association by contiguity
separately, it is thus abstractions we are operating with. They cannot be reduced
to each other; but under some circumstances the one, under others the other,
indicates the most important point of view. If one wants to extend the concept of
association by similarity so as to include immediate recognition as well, it
turns out that the relation of likeness plays a role in every transition of
ideas, since (as I have tried to show in ch. 2) all association by contiguity
presupposes a recognition. We have here, then, a psychological starting point for
understanding the role the principle of identity plays in logical thinking.
Involuntary and voluntary activity of thought are thereby brought into natural
connection with each other.

It was the one-sidedness of the older (English and Herbartian) association
psychology that it regarded the individual sensations and ideas as entirely
independent things that came into connection with one another only in a wholly
external fashion. For the most consistent representative of this whole
tendency, David Hume, the very
% --- p. 59 ---
\setcounter{footnote}{0}%
\opage{59}phenomenon of association therefore became an insoluble riddle. The
mutual difference and independence of the ideas had been so firmly established
that it now became impossible to understand their connection. The unity of
consciousness was conceived merely as a \emph{result} of association, and it was
not seen that it is precisely a \emph{presupposition} of all association.

\begin{center}\large IV. The Association of Ideas and Comparison.\end{center}

26. The theory of the association of ideas is, in its historical development,
closely bound up with the way in which psychical life as a whole has been
conceived at different times and from different points of view. The
investigators who particularly endeavored to demonstrate definite, natural laws
in the domain of psychical life had to lay great weight on the association of
ideas, since this phenomenon shows most clearly of all that a conformity to law
prevails in the inner as well as in the outer nature. They thereby came into a
position of opposition especially to the spiritualistic and mystical views, for
which it was a degradation of psychical life to be subject to a definite
conformity to law. Through this sharp opposition to a tendency that asserted the
freedom of psychical life from all natural order, they were led not only to
conceive psychical life in its close connection with material processes, but
also to regard it as more or less analogous to them, which showed itself in
particular in the placing of the elements of psychical life side by side with
the material atoms. It is simply the case that the expressions of language for
psychical phenomena are metaphors, words which were originally used of material
phenomena (cf. 21); and one must therefore take good care not to let oneself be
led by incorrect analogies. The distinctions that we make as regards psychical
phenomena, and designate by the word elements, must not be turned into
independent real beings. Ideas, e.g., can rightly be called elements of
consciousness; but not in quite the same sense in which we call the atoms the
elements of matter. Every idea is an activity, a particular direction of
ideational activity, and this in turn a particular side of the whole activity
of consciousness at a given moment. The psychical elements thus have their
material parallels not in material atoms but in material processes. Through an
incorrect analogy the so-called ``psychology of association'' then came to
conceive association as an external, mechanical process that goes on between
ideas each independent by itself and at bottom independent of one another. This
one-sidedness reached its height in \emph{Hume}, who (see 25) carried it through
so far that the phenomenon of association itself became incomprehensible to
him: for how --- he asked, quite consistently --- can there be a connection
between elements that are wholly independent of one another? Ideas had been
hypostatized, and it had been overlooked that they are only directions or
expressions of psychical activity. This activity
% --- p. 60 ---
\setcounter{footnote}{0}%
\opage{60}was pushed wholly aside, and where one sought, as \emph{Stuart Mill and Bain}
attempted, to remedy this one-sidedness by recognizing a psychical activity that
expressed itself in attention and comparison, no connection was nevertheless
demonstrated between the passive process that goes on in sensation and
association and the active process in attention and comparison\footnote[1]{Already in my book ``Den engelske Filosofi i vor Tid'' (1874) p.
37 and 87 I have drawn attention to this. --- It is not correct to present
English associationist psychology as a closed unity. A definite change and
development, with attempts at correction, can be demonstrated, as I have sought
to show in the work cited; but full clarity and consistency was not attained by
the English school.}.

No fewer difficulties presented themselves to the opposing, predominantly
spiritualistic tendency in psychology. It did indeed --- often reluctantly ---
come round to recognizing the laws of association, but it sought to conceive
them as the expression of a lower side of psychical life. Just as, in general,
one distinguished sharply between soul and body as two different beings, so in
particular, as regards each single psychical phenomenon, one asked whether it
was due to the action of the body on the soul or to the soul's own action. The
process of association one was inclined to assign to that part of psychical
life which is determined immediately by the material processes in the brain.
According to Charles \emph{Bonnet}, association properly takes place between
the movements in the ``fibers'' of the brain, and to this association between
the causes in the brain there corresponds the association of ideas as an
association between the corresponding effects. Conversely, in the act of
attention the soul sets one brain fiber or another in motion in an active
manner. In our own day \emph{Lotze} has set up ``the comprehending and
relating (`beziehende') consciousness'' as a ``new'' expression of the
soul's activity, a mental activity of a higher order in contrast to
reproduction and association. That higher action does indeed presuppose the
lower, but it does not develop out of it\footnote[2]{Ch. \emph{Bonnet}: Essai analytique p. 106. 255 ff. ---
\emph{Lotze}: Mikrokosmus. 2nd ed. I. p. 251 f. Drei Bücher der Metaphysik. p.
531. --- While \emph{Wundt}'s earlier statements on this point could be
understood as meaning that he regarded association and apperception (the
activity of consciousness in attention) as two different processes, he has now,
in the 3rd edition of ``Physiologische Psychologie'' (II. p. 378), declared
that the distinction between association and apperception rests merely on ``an
abstraction to which reality can more or less approximate.''}.

There is thus every reason to enter somewhat more closely into the relation
between the association of ideas and comparison.

27. Just as with recognition and association, we must also with comparison
distinguish between different stages and forms. To compare is to find likeness
or difference or both. This activity is best known to us in its higher forms;
but it goes on uninterruptedly at all stages of the development of the life of
consciousness.

% --- p. 61 ---
\setcounter{footnote}{0}%
\opage{61}Two objects lie before me on the table in such a way that I can see
them both at a single glance. What psychological process now takes place when I
compare them? The two objects (\textit{A} and \textit{B}) are certainly both
within my \emph{field of vision}; but when I compare them, my \emph{attention}
moves back and forth between them. For although at each moment a certain number
of elements can be present in my consciousness, they do not all appear at each
moment with equal life and distinctness, even if the external impressions are
equally strong. Now the one, now the other stands at the center of
consciousness, as attention is turned toward it. In its simplest form this
appears when the eye is surfeited with an impression of color and therefore
seeks another that forms a contrast to it, or else seeks out a more neutral
impression. With objects given simultaneously, the one object need not now fall
out of the \emph{field of vision} because \emph{attention} directs itself to the
other. But we notice a peculiar change as a consequence of attention's passing
from \textit{A} to \textit{B}. This transition is as a rule bound up with a
certain sensation of effort, which is had in general in attentive comparison,
and which is probably to be explained by a certain exertion of the muscles of
the sense organ taking place, which is necessary in order that the impression
from the object may fall on the ``yellow spot''\footnote[1]{Cf. Psykologi. 2nd ed. p. 136---138. 363---365.}. The result of
this transition of attention, the perception of the difference or the likeness
between \textit{A} and \textit{B,} springs forth in consciousness as something
that we have indeed prepared, but whose character we are nevertheless not
masters of. By the shifting of attention we make an experiment, but we do not
control its outcome.

The properly active element in comparison thus lies in attention, which is
directed in turn to the objects that are compared. Here again we stand before
two concepts that can be kept apart from each other only by artificial
abstraction, which is also why attention and comparison have sometimes been
declared to be one and the same\footnote[2]{Thus \emph{Bonnet}: Essai analytique. p. 203. \emph{Lotze}:
Drei Bücher der Metaphysik. p. 540.}. Yet it has its importance to attend to
the difference between the two elements in our cognizing activity, especially
since we should otherwise not understand why a successive turning of attention
to the different objects is necessary in order to attain a clear perception of
likeness or difference. This successive turning can only have the significance
that the sensation is intensified, i.e., becomes more distinct and livelier, and
this in turn can only be explained by reproduction and immediate recognition
taking part. (Cf. 5.) Even where attention is involuntary, i.e., does not
presuppose any prior idea of what is observed, acts of recognition will not be
absent; for since attention works not continuously but rhythmically,
reproductions will occur at all the later turnings to the same object.

% --- p. 62 ---
\setcounter{footnote}{0}%
\opage{62}If the objects that are compared are not present to consciousness
simultaneously, as belonging to one and the same field of vision, but succeed
one another, a process similar to the one just described takes place. If I
sense first \textit{A,} then \textit{B,} I must, if I am to be able to compare
them, either keep the idea of \textit{A} in consciousness or reproduce
\textit{A,} after I have come to \textit{B.} In this case too, then, a
simultaneity is brought about between the elements that are compared. Here too
there takes place at the same time (at any rate in the more difficult acts of
comparison) a successive turning of attention from the \textit{B} given to
sensation to the idea of \textit{A} and back again. --- The same naturally
holds when it is two ideas that are to be compared. Ideas can, after all, be
recognized just as well as sensations (6).

If in the transition of attention from \textit{A} to \textit{B} I notice no
difference at all (apart from the exertion of the transition itself), or notice
a difference only partly or in slight degree, I conceive their relation as
identity, or as composite likeness (in which some elements are identical,
others different), or as qualitative likeness or as likeness in relation.

The kind of comparison described here can be called free and
\emph{involuntary}: free, because the elements, either simultaneously or
successively, appear as independent members in consciousness, --- involuntary,
because no intention, no prior idea, is presupposed. One sees clearly here how
the active and the passive side work together: an object attracts attention to
itself and thus arouses an exertion of force and a reproduction, and as the
effect of this active turning the perception of likeness or difference arises.
It is likewise seen clearly that the general comprehending fundamental activity
of consciousness is a necessary presupposition of comparison\footnote[1]{In my book ``Den engelske Filosofi i vor Tid'' p. 37 it is
incorrectly stated that comparison is the condition of comprehending. The
relation is the reverse.}. Only when several elements are held together is
comparison possible.

From this free, involuntary comparison we distinguish two forms or stages,
which lie one on each side of it, the one being of a more simple and
uncompounded, the other of a more intricate and compound nature.

28. The most primitive elements of cognizing consciousness are the sensations.
No more than the other elements of consciousness are they wholly independent
and self-standing in relation to one another. Their arising and their character
depend, as the more recent investigations especially have shown (cf. Psykologi
VA), on the mutual relation between them, or more precisely between the
conditions under which they come into being. What makes itself known to us in
the single sensation is properly the difference between two states
% --- p. 63 ---
\setcounter{footnote}{0}%
\opage{63}(in psychical life and brain), a certain relation of difference or
opposition between our states being the condition of the sensation's coming
into being and of its definite character. We can express this by saying that
every sensation comes into being through a \emph{discriminating} or a
perception of difference.

It might seem unjustified to speak of a comparison here, since the two members
that stand in relation to each other are not objects of consciousness, as in
proper, free comparison. Indeed it has often been sought to distinguish between
the sensation itself and its ``appraisal,'' i.e., its comparison with other
sensations. But such a distinction is not justified. For whether the sensation
comes into being, and with \emph{what} character, depends precisely on the
relation between the different states (or --- with simultaneous impressions ---
the different parts of the same state). The sensation can thus be said to be
appraised before it arises. Its being determined by its relation to the other
elements is not something that happens after its coming into being, but is
precisely what conditions this. When it has come into being, it can, in the
manner described in 27, be made an object of free comparison. But what we wish
to bring out here is the important circumstance that the \emph{single sensation
springs forth} in \emph{consciousness in the same way} \emph{as the perception
of likeness} \emph{or difference springs forth}, \emph{when} by free
\emph{comparison we hold two elements of consciousness together}.

We find ourselves here at the border of unconscious psychical life (if this
expression may be allowed), to which self-observation cannot immediately
penetrate, and for which language has not formed wholly fitting expressions.
But it seems to be fully justified to place the border phenomena side by side
with what the clearer, more developed consciousness teaches us, if only we pay
attention to the differences that follow from the different relations. It is a
justified analogy to call sensation a discriminating; for the fundamental
relation is the same, only that in free comparison the members are conscious,
while in this \emph{elementary comparison} they are unconscious. Here, then, an
activity is put forth whose result appears to consciousness, although it is
itself exercised unconsciously. We get the conclusion without the premises (to
use this image), whereas in free comparison we have both premises and
conclusion. If one wished to speak of an unconscious attention as the condition
of all sensation, this expression would not lack justification; for all
attention is an adjustment by which a perception becomes possible; and the
adjustment can take place not only involuntarily but also unconsciously. On this
point too, then, active and passive processes interlock.

Such an immediate discriminating takes place in every comparison that
establishes difference. All that I can do in free comparison is to hold
together the attentively regarded members; the difference comes forth
immediately when the members are placed side by side.

% --- p. 64 ---
\setcounter{footnote}{0}%
\opage{64}The most primitive perception of likeness we have in the immediate
recognition described earlier. Essentially the same process takes place here as
in free comparison that establishes likeness, only that here the members do not
(any more than in elementary discriminating) come forth independently in
consciousness. This process might be called \emph{bound comparison}, because
the two members are here immediately fused. In free comparison likeness springs
forth just as immediately as recognition in the cases described above (3---4).

29. Out of free but involuntary comparison (27) there develops \emph{free},
\emph{voluntary comparison}, which presupposes the ability to set oneself a
purpose, to set oneself a task, to fix a standard\footnote[1]{I use the word voluntariness in a somewhat different sense from,
e.g., \emph{Wundt} (Zur Lehre vom Willen. Studien I, p. 352 cf. Physiol.
Psychologie. 3. Aufl. II. p. 243, 500), who speaks of voluntary action only
where a choice between several possibilities takes place. I call every action
voluntary in which an idea of the purpose is determining, even if only one
single possibility appears. Impulse too (not merely intention and resolution)
is in my usage a voluntary activity.}. Use is made here of attention --- not, as
before, merely to \emph{perceive} given sensations or ideas, but first and
foremost --- to \emph{hold fast} an idea as continuously as possible at the
center of consciousness in order to hold the following sensations or ideas
together with it. This requires a greater energy and development than
involuntary attention and comparison, which follow the given material step by
step.

The transition from involuntary to voluntary comparison is analogous to the
transition from reflex movement and instinct to impulse. The essential change is
that an idea of its purpose makes itself felt prior to the activity. And here
there is in reality more than an analogy. Voluntary comparison, thinking
proper, is an expression of will: it sets an end and seeks the means to realize
it. At first it is only practical motives that set thinking in motion; the
external means to the relief of an urge are sought. Even a practical purpose
gives the connection of ideas a certain order and consistency, since only those
ideas become objects of attention that accord with the idea of the purpose, in
that they point out means to its attainment. \emph{Thomas Hobbes} aptly
explained the lack of coherence in the ideas of dreams by our having no fixed
purpose in sleep\footnote[2]{Qvia ordo et cohærentia omnis a freqventi ad finem respectatione,
id est, a consilio oritur, necesse est, amissà per somnum cogitatione finis, ut
phantasmata alia aliis succedant, non amplius eo ordine, qvi ad finem tendit,
sed ut contingit. \emph{Hobbes}: De corpore. XXV, 9.}; therefore the ideas move
in all possible directions and have no center about which they could arrange
themselves. Every idea, on the other hand, to which a strong and lasting
feeling attaches itself, is
% --- p. 65 ---
\setcounter{footnote}{0}%
\opage{65}held fast thereby, and all associations that go out from it and favor
the ruling feeling are held fast. It is made a \emph{center of association}.
Besides association by similarity and by contiguity, the association between
feelings and ideas (the law of interest) is also particularly at work here.
When the element of feeling belonging to an idea is very weak, because there is
the definite assurance that the phenomenon to which the idea corresponds will
not occur, the phenomenon may remain unnoticed even if it occurs. (Negative
attention.) Thus when a hypnotized individual is forbidden to see a certain
person or object which is nevertheless actually present\footnote[1]{Cf. Revue philosophique. Mai 1887. p. 455. --- The expression
``negative attention'' (which is used by \emph{Féré}: Sensation et Mouvement.
Paris 1887. p. 18) is more correct than ``negative hallucination'' (which is
used by Bernheim and Pierre Janet). For the individual \emph{does get} a
sensation, as is seen from the fact that a contrast sensation sometimes arises
(so that the somnambulist sees a greenish spot where it had been forbidden to
see anything red). It is thus the turning away of attention that is the proper
cause.}.

The more clearly and definitely I am able to represent to myself the relation in which the
ideas sought are to stand to the ideas I have, the more surely and precisely
thinking proceeds. When in mathematics I put a problem into an equation, this
means precisely that I make exact the demands placed on the ideas sought, ---
that I represent to myself the end I wish to reach, the standard I wish to apply, with
perfect precision. Here, therefore, we can force out the idea sought by a
methodical procedure. In other domains we are obliged to feel our way forward by
trial. In its simplest form this happens when we ``bore'' an idea forth by
energetic concentration of attention on the given idea with which the idea
sought is known to stand in some connection or other. It is often by
association by similarity that the idea sought springs forth, as in discoveries
of homologies and analogies between organs in different species. The other
member of the comparison can here spring forth while the investigator is
absorbed in the single given phenomenon. As at a stroke, the kinship between the
phenomenon at hand and others far removed from it in time and space then dawns
on consciousness. What is at work here is the circling about a ruling idea
mentioned earlier (24). ---

According to \emph{Aristotle}\footnote[2]{De anima III, 10.} practical reason differs from theoretical
reason in that it starts from a purpose. By this doctrine he laid the foundation
for a dualistic distinction between theory and practice, between knowledge and
will, which has had an unfortunate influence on the psychological view of both
these sides of psychical life. In particular it was overlooked that in all
thinking a purpose, a striving, an interest makes itself felt. Only through the
thought of a purpose do order and coherence enter into our connections of
ideas, so that a \emph{valid} knowledge can be reached through them. The
purpose we set ourselves need not be an external one; it can be that of
understanding nature
% --- p. 66 ---
\setcounter{footnote}{0}%
\opage{66}or of preserving consistency and excluding contradictions among our
ideas. In all thinking a choice takes place, in that out of all the sensations
and ideas that present themselves we hold fast those that satisfy the demands
set up. This selecting procedure shows itself in the law of \emph{dichotomy},
which, as \emph{Wundt} has shown\footnote[1]{Logik. Leipzig 1880. I. p. 557.}, is peculiar to the connections of
ideas that come about through thinking proper (voluntary comparison), as
distinct from the involuntary connections of ideas, which spread out into the
indefinite. The dichotomy arises from the question being raised at every point
of the movement of thought: does the idea now emerging accord with the idea of
the conditions that the idea sought is to fulfil? Here an Either---Or is posed
at each moment, and the comparison takes place at each single moment between
two ideas at a time (of which the one is the standard), because attention can at
each moment be concentrated on only one idea. ---

In voluntary comparison there are four moments to distinguish, of which some
have a more passive, others a more active character: 1) The starting point is
the impetus which, through the energetic \emph{concentration} on a single idea,
is given to the process of association in a certain definite direction. 2) The
\emph{process of association} itself runs its course involuntarily, although
single members of it will from the very outset be brought into prominence by
virtue of their relation to the ruling interest. This second moment is related
to the first roughly as the entry of air into the lungs is related to our
active, voluntary expansion of the chest. 3) \emph{Recognition} and holding
fast of the ideas that seem to accord with the idea of the purpose or the
standard. 4) \emph{Verification} of this recognition through voluntary attempts
at congruence, which lead to immediate perception of likeness or immediate
discriminating, and then perhaps to a closer determination of the degree of the
likeness or the difference.

Where it is a series of successively occurring phenomena that are the object of
voluntary comparison, the series of observations takes the place (in the second
moment) of the series of associations. The difference is only this, that with
the observations the cause of the emergence of the new members lies outside us,
with association it lies in ourselves. If we can have any influence on the order
in which the observations come, the associations will in this case too play a
part.

If this description of thinking proper or voluntary comparison is correct, there
is no reason to assume a faculty of thought (faculty of appraisal, reason, or
whatever one may wish to call it) wholly different from the faculty of
association. Attention, the element of will which makes thinking into thinking
proper, is found also (as involuntary) in all sensation and all association of
ideas. In thinking proper it takes on
% --- p. 67 ---
\setcounter{footnote}{0}%
\opage{67}only a peculiar form. However significant the transition from the
involuntary to the voluntary may be, both as regards knowledge and as regards
the will itself, it is nevertheless not justified to assume that a wholly new
faculty should here come into force.

One of the older psychologists who laid much weight on the element of will
(attention) in thinking (at the same time as he developed the doctrine of
association with great skill), \emph{Charles Bonnet}, has said: ``L'attention est
la mère du génie''\footnote[1]{Essai analytique. p. 313.}. He thus finds intellectual genius
especially in the ability to concentrate attention on an object. But however
important this moment may be, it is nevertheless only the introduction, the
impetus. And it is no less important that a lively process be released by this
impetus, just as it is not enough that we can expand the chest when the air is
for some reason or other prevented from streaming in. It then depends on whether
the idea that has been made the center of association can release a long and
well-connected series of associations. And in this involuntary, instinct-like
side of thinking, genius perhaps expresses itself quite as much as in the first
moment. There is thus no reason to regard the process of association as a lower
form of psychical life. Precisely the highest psychical content can be developed
along this path.

And finally it must always be remembered that in our psychological terminologies
and theories we cannot avoid bringing out the difference between the involuntary
and the voluntary far more strongly than it appears in reality. \emph{The very
transition from involuntary to voluntary connection of ideas takes place} by
\emph{virtue} of the \emph{involuntary association between a feeling} and an
\emph{idea}, \emph{whereby an impulse}, a striving guided by an \emph{idea} of
an \emph{end}, \emph{becomes possible}. The continuity of psychical life shows
itself to us ever anew when, from the differences that provisional observation
presents, we penetrate to the inner course of life.

\begin{center}\large V. On the Concept of Psychical Activity.\end{center}

30. It goes with the concept of activity in the psychical domain as with the
concept of motion in the domain of external nature. Just as one was long caught
in the assumption, apparently supported by observation, that there was absolute
rest, so one is also inclined to think that certain forms and stages of the life
of consciousness are absolutely passive. And just as it has there turned out
that everything is in motion, or rather, that motion and rest are relative
concepts, so too it will turn out here on closer consideration that the nature
of psychical life consists in activity, and that it is only different
% --- p. 68 ---
\setcounter{footnote}{0}%
\opage{68}kinds and degrees of activity that we set in opposition to one another
when we distinguish between passive and active psychical phenomena.

The activity of consciousness has been described in the foregoing from two
sides. From the formal side it expresses itself in the comprehending and
combining activity which all phenomena of consciousness bear the stamp of, and
whose results they are. From the real side it expresses itself in the attention
at work in all association and comparison, which is in its nature one with the
will. If one wishes to designate all psychical activity as a willing, one can
thus distinguish between a formal willing, which constitutes the proper essence
of consciousness, so far as we know it, and which is the expression of its urge
to self-preservation, --- and a real willing, which shows itself in the
particular direction that the life of consciousness takes at the single
moment\footnote[1]{Cf. Psykologi. 2nd ed. p. 113, 159, 367.}.

31. To what extent and in what way do we now become conscious of this psychical
activity?

While sensualism and associationist psychology sought to push the concept of
psychical activity wholly aside, denying that psychological experience showed us
anything other than passively received sensations that were passively
transformed by association, other psychological views have ascribed to us an
\emph{immediate} consciousness of activity as constituting our own innermost
nature. According to this view it is a primitive psychological fact, attested
by self-observation, that we \emph{will}, i.e., exert ourselves, apply force.
This sensation of exertion and of putting forth force (effort) is then supposed
to be the fundamental fact of psychological experience, which cannot be doubted
but also cannot be further explained, and in which our proper nature makes
itself known to us. This thought \emph{Maine de Biran} developed in his various
writings. Against Hume's doctrine that there was not, in any activity whatever
of a bodily or mental kind, anything that could lead to the idea of force or of
necessary connection, he set the claim of an immediate consciousness of the self
itself as cause\footnote[2]{Cf. e.g.: Oeuvres philosophiques. Paris 1841. IV. p. 75. Le moi
se manifeste à lui-même par le sens intime, dans l'effort ou le mouvement
volontaire que l'âme aperçoit intérieurement comme un produit de son activité,
comme un effet dont sa volonté est cause.}.

Both the contending parties appeal to immediate psychological experience. This
experience must then, on the one side or the other, be mixed up with unjustified
inferences. We shall examine more closely the so-called consciousness of will,
which in its simplest form appears in the sensation of exertion to which it is
Maine de Biran's merit to have pointed so energetically, even if he perhaps did
not analyze it with sufficient precision. There is all the more reason to
undertake a thorough analysis of this psychical phenomenon, since recent
psychology in reality
% --- p. 69 ---
\setcounter{footnote}{0}%
\opage{69}confirms Maine de Biran's claim that we here stand before a
fundamental psychical fact. For while the old faculty theory distinguished
sharply between knowledge and will, recent psychology lays strong weight on
there being a constant co-operation of the will in all sensation, association,
and comparison. Thinking is conceived as a matter of will, no less than external
action.
Thereby the question of the consciousness of will is not only moved up into the
first rank among psychological problems, but also becomes easier to elucidate.
What is essentially at issue is the right psychological view and understanding
of \emph{attention}.

If we seek out the very simplest phenomena in which attention clearly shows
itself, it turns out that it expresses itself through a movement by which the
sense organ is involuntarily brought into a suitable relation to the
stimulating object. (Cf. Psykologi V A, 7.) There is thus no doubt here that
attention makes itself known in consciousness through \emph{sensations of
movement}. And here the difference between inner and outer action falls away,
or both stand in the closest connection, the outer action (the adjustment of the
sense organ) being the condition of the inner (the perception of the object).
It might seem otherwise when attention is directed to an idea, and not to an
external object. How can movement and the sensation of movement play a part
here?

Yet here too self-observation shows clearly enough that all energetic holding
fast or bringing forth of ideas is bound up with a sensation of exertion or
effort. And it is likewise easy to show that in every such mental effort a
certain muscular tension takes place. The very appearance of a person absorbed
in thought shows this. And the connection between eager ideational activity and
muscular tension will not be difficult to demonstrate. In the first place, all
ideas, however abstract and ideal they may be, rest in the end on sensations
and reproductions of sensations. Now if attention in sensation expresses itself
through movement, this must also be the case with the attention attached to the
idea, since every idea is in the end only reproduced sensation. In the bringing
forth and holding fast of the idea, the idea of movement (the reproduced
sensation of movement) will therefore play a great part. The same thing will
happen in a somewhat weaker degree as when it is external objects that attract
attention. In the second place, all attention is conditioned by a feeling of
pleasure and pain, and every state of feeling is in greater or lesser degree
bound up with movement. Since strained attention is an intensive state of
feeling, muscular tension will thus always be present, and consequently also
sensations of movement as an essential constituent of the whole state of
consciousness.

Attention, and with it the consciousness of will in general (for in external
action too it is first and foremost a matter of holding fast ideas: ideas of
% --- p. 70 ---
\setcounter{footnote}{0}%
\opage{70}ends and of means), could then, it seems, be traced back to sensations
of movement, partly original, partly reproduced. The two concepts, psychical
activity and tendency to movement, would be identical. One could then, without
the mysticism that often comes out in Maine de Biran, maintain the proposition
that we become immediately conscious of our activity, since consciousness of our
activity would be the same as consciousness of movement or tendency to movement.
Thus \emph{Wundt} remarks that all attention (or in his terminology:
apperception) is a conscious activity, and that we immediately sense it as an
inner action\footnote[1]{Physiologische Psychologie. 3. Aufl. I. p. 535. II. p. 466 ff.}.
And he states expressly that when he calls apperception an activity, it is
chiefly on account of the innervations that take place in it\footnote[2]{``In diesen Innervationsempfindungen liegt der nächste sinnliche
Anlass dafür, dass wir die Apperception oder Aufmerksamkeit eine
\emph{Thätigkeit} nennen''. Zur Theorie des Willens (Studien I) p. 347.}.

This identification of the consciousness of activity and the consciousness of
movement could, however, be made the object of critical misgivings. It must
necessarily assume that there are sensations of movement which correspond to the
very departure of the motor impulse from the motor centers in the brain, or to
innervation itself, the initiation of the nerve process. But this is still a
disputed point. If we keep to self-observation, there certainly seems to be good
reason to assume two classes of sensations of movement, of which the one could
be said to have a more active character, since it is assumed to point to the
very initiation of the movement in the brain center, while the other is of a
more passive kind, since through it we get information about the state of the
muscles as a consequence of the contractions initiated or carried out. One
might call the first class sensations of force, the second muscular sensations.
The former would thus correspond to a process going out from the brain
(efferent), the latter to a process coming to the brain (afferent). It is clear
that if the consciousness of psychical activity is to be one with sensation of
movement, it must be with ``efferent'' sensations. For if there were only
``afferent'' sensations, thus only ``muscular sensations'' and not ``sensations
of force,'' we should after all have in the sensations of movement only
consciousness of the result of an activity, not of the activity itself. And if
one nevertheless wished to maintain that there is an immediate consciousness of
psychical activity, this consciousness could not be consciousness of movement,
since the latter (if all sensations of movement were ``afferent'') would not be
more active than all other sensations.

In his treatise ``The Feeling of Effort''\footnote[3]{Anniversary Memoirs of the Boston Society of Natural History.
1880.} the American psychologist \emph{William James} regards the physical
sensation of effort as ``a compound afferent sensation that comes from the
tensed muscles, the strained tendons, the compressed joints, the tensed chest,
the closed glottis, the knitted brows, the clenched jaws, etc.'' And from this
sensation of physical effort he sharply distinguishes
% --- p. 71 ---
\setcounter{footnote}{0}%
\opage{71}another kind of sensation of effort, which is not motor but arises in
``the effort to remember, to make a resolution, or to give oneself up to an
unpleasant work.'' This last kind of effort, he thinks, is characteristic of
willing proper, which is a purely inner, psychical phenomenon that exists wholly
independently of whether movement occurs or not. This willing proper is supposed
to express itself in the assent, the adherence, that we give to a thought, in
the \emph{Fiat}! with which we settle that an action is to take place. What
demands effort and exertion of will is the holding fast of a thought, even if it
brings pain or suffering with it. This sensation of effort, bound up with all
belief and all willing, with most expressions of mental action, stands according
to James in opposition to all ``afferent'' sensations (thus also to all
sensations of movement). On the physiological process to which it is attached he
will not express himself more closely; he is only certain that it is the centers
of the idea, not of movement, the ideational, not the motor centers, that play a
part here.

I shall not here go further into the discussion of the question whether all
sensations of movement are peripheral (afferent), or whether there are also
central (efferent) sensations of movement (sensations of innervation)\footnote[1]{In the 3rd edition of his ``Physiologische Psychologie'' (I. p. 405),
\emph{Wundt} has replied to James's criticism of the doctrine of sensations of
innervation. As an experimentum crucis that was to confirm the assumption of
central (efferent) sensations of innervation, it had been adduced
(\emph{Helmholtz}, \emph{Gräfe}, \emph{Wundt}) that patients suffering from
paresis of a certain eye muscle (m. rectus externus) locate objects farther to
the side than they really are, which is then explained by the stronger (vain)
central exertion to move the eye. Against this \emph{James} remarks (following
\emph{Ferrier}) that it has been overlooked that the other eye, which was
covered in the experiment, involuntarily moves in the direction of the object,
and with all the greater energy since the diseased eye cannot be moved. From
that healthy eye there now come afferent influences, and the displacement to the
side can thus very well be motivated by peripheral sensations (On the Feeling of
Effort. p. 12). In his reply \emph{Wundt} draws attention to the fact that the
displacement to the side in paresis occurs only when the normal eye is closed;
when it is open and can take part in the localization of direction, the
illusion does not occur, or \emph{if} double images arise, it is only as
regards the image of the paralyzed eye.
--- Adhuc sub judice lis est.}. For I believe that too great weight has been
laid on this question from both sides.

In the first place, \emph{ideas of movement} could very well make themselves
felt in an act of attention and of will, even if all \emph{sensations of
movement} were peripheral. For because the sensation of movement is due to an
influence from the muscle to the brain, it can very well, like every other
sensation, be reproduced, and then corresponds essentially to a central process,
which nevertheless, like all stronger central processes, has a tendency to
spread and thus to reproduce completely the original, peripherally called-up
process. Like every other idea, the idea of movement too has a tendency to
approach as nearly as possible the original sensation
% --- p. 72 ---
\setcounter{footnote}{0}%
\opage{72}of which it is the reproduction. The fact that we can have motor
hallucinations, just as we can have visual or auditory hallucinations, points in
this direction.

In the second place, the sensation of movement, even if it always corresponds to
peripheral processes, can very well occur before the movement is completed, and
can thus very well correspond to the \emph{initiation} of the movement. The
contraction of the muscle takes some time, and what we psychologically call the
sensation of force might be assumed to correspond to the very first stage of the
contraction, where it is not yet visible externally. What we psychologically
call the muscular sensation might then correspond to the \emph{completed}
contraction, and could to that extent rightly be said to have a more passive
character.

But whether one now conceives the matter in the one way or in the other, the
sensations of movement would not be sufficient to ground the inward and absolute
consciousness of activity that Maine de Biran has sought to demonstrate. Whether
our consciousness of will has its physiological correlate in processes that are
released immediately in the brain itself or in such as are released there only
by virtue of impressions from the muscles, it will in both cases correspond to
processes that are themselves effects of preceding processes. There will in
itself be no reason why these processes in particular --- the release of
impulses to motor nerves --- should be called active in a higher degree than all
the other processes of release that go on in the nervous system.
Sensations of movement have their own peculiar stamp, by which they are
distinguished from other kinds of sensation; but it is not thereby given that
this stamp is particularly to be designated as activity, and still less is it
given that all psychical activity should make itself known through this kind of
sensation. \emph{Through the parallelism with physiological processes} one will
therefore hardly be able to elucidate why certain psychical states seem to us to
have a particularly active character.

If we now consider the matter from the \emph{purely psychological side}, the
question arises whether we can become immediately conscious of activity at all.
All that we find immediately in our consciousness by self-observation --- also
in the phenomena of exertion so strongly emphasized by Maine de Biran and James
--- is after all only sensations, ideas, and feelings of different qualities and
in different degrees of composition. There is besides a great difference in the
degree to which the single elements come to rule in consciousness, in the power
they have in relation to other elements. But can we notice this power
\emph{immediately}? In that case the consciousness of activity would have to be
quite simple and uncompounded, as our simple sense sensations are; it would
appear with a peculiar quality that would be just as little open to mistake as
the quality of our sensations of color. \emph{The question} is whether
\emph{the quality of activity really is attached in this way to quite simple
states of consciousness}.

% --- p. 73 ---
\setcounter{footnote}{0}%
\opage{73}32. The peculiar state in which we are at every resolution, every
decision we believe we have come to, or even where the decision has not taken
place but where we seek to suppress some elements in our consciousness and to
hold fast and favor others, cannot be described otherwise than as a
concentration of idea and feeling as far as possible on one point, namely on the
idea of the action we will carry out, or the assumption we acknowledge. A center
of association has formed (cf. 29), which we regard as the most essential
content of our self at the moment, and with which we make ourselves one as far
as possible. Even if other ideas and feelings rise up and lay claim to a part of
our energy of consciousness, they nevertheless stand over against us as more
foreign; we do not recognize ourselves in them with the clarity with which we do
in that center; they stand before us as something new and forced on us from
outside. We are approximately in a state of monoïdeïsm: one idea (with its
accompanying feeling) masters us. We have in a way hypnotized ourselves.

The reason why we designate such states as active lies, no doubt, in two
circumstances. \emph{Partly} the reason must be sought in the concentrated,
pointed form itself that our life of consciousness takes on, and which excludes,
or tends to exclude, all inner resistance on the part of other elements of
consciousness. As soon, on the other hand, as several different elements make
themselves felt and lay claim to our energy without being able to fuse more or
less inwardly with one ruling idea, we have a sensation of dividedness, discord,
or at any rate of neutrality or equilibrium. \emph{Partly} the reason is to be
sought in this concentration's being followed by inner or outer changes that
can with unusual distinctness be referred to the concentrated state as their
cause. Every state we are in gives rise to inner and outer changes as its
effects. But the more manifold and heterogeneous its content is, the less
distinctly and strikingly does the relation of causality between it and what
follows appear. It is then no wonder that we sense ourselves as particularly
active, as cause, when our consciousness is strongly concentrated.

If the consciousness of activity were wholly immediate, the first reason would
have to be sufficient. The second reason presupposes that besides our immediate
consciousness we also take account of what follows. When we find the criterion
of activity in what follows, the consciousness is not quite immediate; a
comparison and reflection, a kind of inference, makes itself felt. The second
reason is related to the first as verification to hypothesis. It is incorrect to
ascribe to us an immediate consciousness of being the cause of something in us
or outside us; such a consciousness can, by the nature of the case, never be
immediate. \emph{Causality} can only \emph{be discovered by inference},
\emph{not by intuition}. Hume remains in the right: observation (and that
immediate consciousness is a self-observation) can only show us a relation of
succession, never a relation of causality. It is therefore mysticism when Maine
de Biran and other psychologists ascribe to us an immediate consciousness of
activity in the sense that we should immediately notice that we are cause.

% --- p. 74 ---
% print: „at at“ (misprint)
\setcounter{footnote}{0}%
\opage{74}Let us then keep to the first reason and examine whether it is
sufficient.

It is easily seen that it is, or can be, extraordinarily difficult to decide
whether the concentration spoken of has taken place or not. A certain
concentration is always present in consciousness and is a necessary condition of
the existence of consciousness. The more or less manifold content of
consciousness must always be comprehended and held together; the character of
consciousness as synthesis can never wholly deny itself so long as the life of
consciousness still exists. There are countless degrees of the energy with which
the synthesis works; but a certain concentration follows necessarily from the
general nature of consciousness. And apart from this formal coherence there will
also always be a certain real coherence present: feeling and attention will at
each moment be directed to certain ideas more than to others. At \emph{what
point} on the \emph{scale of degrees of concentration} that we \emph{can set up}
here \emph{are} we now to \emph{think} that the \emph{state lies} \emph{that
should have a special claim} to \emph{be called consciousness of activity}?
Should there be a particular ``tone color'' by which precisely this definite
degree of concentration made itself known to consciousness itself? --- If there
were such an immediate and unfailing criterion, some of the greatest
difficulties of practical self-knowledge would fall away. On the one hand there
can be thoughts and moods that we play with and regard as harmless, and which
can nevertheless later turn out to have taken root and to have become a willing
in us. On the other hand it can be wholly impossible for us to come to clarity
about the character of our will where it does not get occasion to show itself
in action. Conversely, one can torment oneself by confusing a moment's ignoble
mood or wish with a concluded willing, and thereby judge oneself unjustly. At
all times deeper moral and religious experiences have shown that here lies one
of the greatest problems in the domain of psychical life. In psychological
theory it is easy enough to distinguish in abstracto between passivity and
activity; actual experience does not find it so easy, and it warns precisely
and definitely against trusting immediate consciousness: Noli credere affectui
tuo, qvi nunc est; cito mutabitur in aliud (De imitatione Christi. III, 33). ---
There is also clear enough experience of how difficult it can be to state the
definite point in time at which the inner decision was made\footnote[1]{Cf. the case analyzed in my Psykologi p. 396.}.

There is one more circumstance that causes self-observation particular
difficulty here, which I shall point out. What makes itself felt in
consciousness with the greatest vehemence and violence, and brings about the
greatest concentration on a single point, need not on that account be what is
in reality the strongest and the victorious. For there are two kinds of
psychical strength; we may designate them as the strength of violence and the
strength of inwardness. Apart from single sudden fluctuations, the whole
direction of psychical life will be predominantly
% --- p. 75 ---
\setcounter{footnote}{0}%
\opage{75}determined by the latter kind of strength. But this works distributed
over many elements of consciousness and many moments of consciousness rather
than concentrated in a single element and a single moment. Self-observation,
which keeps to the most prominent points, has more difficulty in becoming clear
about relations of strength of the more inward or distributed kind. And yet our
\emph{activity}, \emph{our willing can be strongest} at \emph{such times} and at
\emph{such points} \emph{where we notice it least} in our
\emph{self-observation}. And when it comes to our saying: ``I will!'' (or with
W. James: ``Fiat!'') the real turning point most often does not lie here at all,
but far earlier and at quite another point. What happens in our clearly
conscious resolution is then only an establishing, a so to speak official
conclusion of what has in reality been decided\footnote[1]{Cf. on this Psykologi p. 108, 326 ff., 380, 397 f. --- See also
\emph{Ehrenfels}: Ueber Fühlen und Wollen. Wien 1887. p. 76 and \emph{Ribot}:
Les maladies de la volonté. Paris 1883. p. 175.}.

To this is further added that the highest degree of exertion and concentration
to which self-observation bears witness comes to stand, for the subject himself,
as a state of surrender, of suffering, or of passivity. It is not accidental that
two states so opposite as the hypnotic, in which the individual has lost mastery
over himself, and the mystical, in which the individual with the greatest
self-exertion and self-activity absorbs himself in one single thought, are,
psychologically regarded, so closely akin. They both have the character of
monoïdeïsm, and in all absorption in one thought, or in general in one element of
consciousness, activity and passivity pass over into each other. Exertion forms
the introduction to rapture, and what the mystics call ecstasy designates the
point where the more active state gives place to the more passive. The direction
the mystics give for reaching this point amounts simply to an increasing
abolition of all manifoldness and change in idea and memory. The will (``the
ability to love'') holds the object fast and permits no ideational activity that
leads away from it; every rising memory of other things is inhibited. That this
ecstatic state is nevertheless more active than the hypnotic is seen from the
fact that while the hypnotized individual falls in with every idea that
corresponds to the postures one makes it take or the gestures one makes it
carry out, the ecstatic mystic holds fast the thought on which he has
concentrated himself, and takes the postures and shows the gestures that
correspond to it\footnote[2]{See \emph{Maury}: Le sommeil et les rêves. Paris 1865. p. 247.}.
The former is always determined from without inward, the latter from within
outward. And yet the mystic is fully convinced that the highest stage of
concentration is reached not by voluntary willing but by an ``impulse of grace,''
by an inspiration, toward which he stands in a receptive relation. He no longer
grasps, but is grasped\footnote[3]{In \emph{Ribot}: (Les maladies de la volonté. Paris 1883.
123---35 and La psychologie de l'attention. Paris 1889. p. 138---50) are found
interesting contributions to the psychology of mystical ecstasy, built especially
on St. Theresa's writings.}.

% --- p. 76 ---
\setcounter{footnote}{0}%
\opage{76}What thus takes place in mystical ecstasy, which is the height of inner
concentration, takes place in a lower degree in every resolution. Anyone who has
taken a decisive and significant resolution will have had an immediate sensation
that there was something that drove him forward, a force that seized him, and
that his inner state was only imperfectly expressed by the word activity. The
expression that so often recurs in a serious decision of will, ``I cannot do
otherwise!'' (cf. e.g. Luther at Worms), is characteristic in this respect. One
appeals not to one's own self-activity but to an objective impossibility of
willing otherwise. We not only grasp what is resolved, but are ourselves grasped
by it. Where the active ends and where the passive begins, no self-observation
can show us.

What now follows from all this? It might seem to be the natural consequence that
we should have to give up the usual threefold division in psychology, since the
will seems no longer to be able to be recognized as an independent side of the
life of consciousness. But such a conclusion can be drawn only on the basis of a
misunderstanding of the psychological distinctions. These must not be confused
with real separations between different phenomena given to immediate
observation. Otherwise we could not rightly distinguish between knowledge and
feeling either, for this distinction too rests not on immediate observation, but
partly on the fact that the same sensations and ideas can under different
circumstances be bound up with different feeling, and conversely, partly on the
two sides of our states of consciousness turning out to be subject to different
conditions as regards development, reproduction, and repetition. Just as the
chemical elements cannot be established by immediate external observation but
are discovered by analysis, so the psychical elements, the mutually irreducible
characters of our states of consciousness, must also be discovered by analysis,
naturally on the basis of self-observation. As regards the active (willing) side
of consciousness, it now follows of itself that there can be no question here of
an immediate observation of it. For activity and causality we discover by
inference from the relations of succession given to observation; in a single
observation, an immediate state, they cannot be given. This holds in the
psychical domain just as well as in the physical domain. In the phenomena of
consciousness given to self-observation we always have the results of activity,
not activity itself.

The more inwardly activity is bound to psychical life, the more evident it is
that it cannot be the object of a single and immediate act of consciousness.
Precisely because psychical life is an uninterrupted willing, it can become
difficult to find the will in a special state. Just as little as we can build
our concept of the I or the self on a single sensation or idea, just as little
is this possible in the case of the will. So far, then, is the denial of a
special consciousness of activity from leading to a denial of the reality of
psychical activity, that it is grounded precisely on the conviction that
% --- p. 77 ---
\setcounter{footnote}{0}%
\opage{77}psychical life subsists under continuous activity, of whose results
self-observation is able to grasp only single prominent points, from which we
then, by way of construction, form for ourselves a concept of the activity that
lies at the bottom. Our knowledge here, as everywhere, goes from present, more or
less sporadic and heterogeneous experiences to the conception of the inner,
continuous coherence.

33. A certain obscurity nevertheless seems still to cling to the concept of
activity. We have emphasized that psychical life is active through and through.
This activity can appear more or less concentrated, so that (in accordance with
Fechner's rule) the lower, more elementary degree of activity, such, e.g., as
shows itself in sense sensation and reflex movement, can appear as a passivity in
relation to a higher degree of activity, as in thinking proper and conscious
resolution. But what then do we understand by activity at all? We have, after
all, so far used this word without defining it.

If one wished to say that this concept cannot be defined, much might seem to
speak for such a claim. For it is not only the most fundamental concept in
psychology; it is also the fundamental concept in every thoroughly thought-out
view of the world, since all passive or resting being turns out on closer
investigation to be an activity. Newton in his time warned against the sensuous
view of external phenomena that distinguishes between rest and motion, since
``it could after all be that no body is in reality at rest.'' This proposition
we could extend to hold not only of motion but of all activity. But this would
lead us too far here, if it were to be sufficiently developed. My meaning here is
only that one could, with just as good right in philosophy, declare the concept
of activity indefinable as one can in physics declare the concept of motion
indefinable\footnote[1]{In ``Drei Bücher der Metaphysik'' (1879) p. 534 \emph{Lotze} calls
``den Begriff der Thätigkeit'' ``einen Begriff, von dem wir glauben werden, dass
er etwas Eigenthümliches und in der Welt wirklich Befindliches bezeichne,
\emph{obwohl} wir es \emph{ganz unmöglich finden}, das, was wir mit ihm im
\emph{Gegensatze zum blossen Geschehen meinen}, \emph{auf irgend eine
mechanischer Construction sich annähernder Weise zu bezeichnen}''. ---
\emph{James Ward} says in Mind. 1887. p. 569: ``I doubt if activity can be
defined in terms that do not already imply it''.}.

Yet one could give a definition of the concept of activity by holding fast its
close connection with the concept of causality. \emph{Spinoza} has accordingly
defined activity in the following way: ``I say that we act (\emph{agere}) when
something happens in us or outside us of which we are the complete
\emph{cause}, i.e., when something follows from our nature, in us or outside us,
that can be clearly and distinctly understood through it alone. On the other
hand, I say that we \emph{suffer} (pati) when something happens in us, or
follows from our nature, of which we are only in part the cause''\footnote[2]{Ethica III. def. 2. --- A similar definition is found in
\emph{Charles Bonnet}: Essai analytique. p. 93.}. This definition has the
merit of distinguishing
% --- p. 78 ---
\setcounter{footnote}{0}%
\opage{78}between the phenomena, the given, ``what happens in us or outside us,''
and its explanation. It thus does not appeal to any special kind of immediate
consciousness, but proceeds from the view that every such thing needs an
explanation. It indicates at the same time that a part of the cause of what
happens in us or follows from our nature always lies in ourselves, so that we
are never absolutely passive; but it would restrict the application of the
concept of activity to such cases in which the whole cause lies in ourselves.

The question is whether it can ever lie completely in ourselves, whether any
state of consciousness and any action can be understood from our preceding life
of consciousness alone, and whether there are not always other contributing
circumstances. We therefore alter the definition cited as follows: The
\emph{greater} the part of the \emph{cause} of what happens in us or outside us
that lies in ourselves, the \emph{more} active we are. And to this it must be
expressly added that even where we behave passively, the kind and degree of the
effect are nevertheless always determined also by our own nature, not merely by
the character of what acts upon us. --- In accordance with what has been
developed in the foregoing, it is a question by itself whether we become
conscious of our activity; and the consciousness we can have of it is in any
case not always reliable.

But this definition then calls up in turn the question: \emph{what is meant} by
``our\emph{selves}'' \emph{or} ``\emph{our nature}''? --- Can this question find
any answer, and what answer?

A frequent view gives the following answer. It lies in the nature of the case
that every activity must be exercised by someone or something. We must
distinguish between that which acts and the \emph{activity} itself. In particular,
in the psychological domain we must distinguish between that which thinks, feels,
and wills, and thinking, feeling, and willing themselves. --- According to this
view we should thus be the more active (in accordance with the definition given)
the more the psychical phenomena spring from the very \emph{bearer} of all our
thoughts, feelings, and resolutions. Of this something that ``is'' conscious or
``has'' consciousness we know, to be sure, nothing; it is declared to be an
\textit{x. ---} This is the \emph{monadological}-\emph{spiritualistic} theory. I
am unable to see either its justification or its possibility.

In the first place, the rise of the theory mentioned can simply be explained by
the unjustified influence partly of linguistic expressions, partly of material
analogies.

Language distinguishes between subject and predicate and as a rule designates
each of them by a word of its own. A word is then formed for the acting being,
another for its activity. Even where it is difficult for intuition to make this
distinction between the acting thing and the activity, language nevertheless
employs a formal subject; thus in impersonal sentences (``it rains''). There is
expressed here an urge to analyze the immediately given phenomena and to keep the
different elements in them apart from one another. But separate linguistic
expressions and logical distinctions must not without further ado be made one
with real differences. The necessities we feel when we wish to think about the
phenomena and
% --- p. 79 ---
\setcounter{footnote}{0}%
\opage{79}express our thoughts are one thing; real existence itself and its
differences are quite another. --- That distinction between the acting thing and
the activity, --- between what ``has'' consciousness and consciousness itself,
may perhaps find its explanation in linguistic usage.

In material phenomena (which have had by far the most important influence on the
mode of expression of language) we easily distinguish between the starting point
of the activity (here: of the motion) and the activity (the motion) itself. We
always seek to trace the origin of the motion back to a body located at the place
in space from which the motion is assumed to come. We often see this body remain
apparently at rest in its place, whether the motion it can cause actually takes
place or not. Whether the stone falls from the roof or not, the earth remains
apparently in its place. When one billiard ball sets the other in motion by
striking it, we call the first active in this relation, and we are not then at a
loss here about the origin of the activity; we can, after all, see it and
picture it, quite apart from whether it is at the moment calling up any change or
not. In an analogous way, then, the soul too is thought of as separate from the
activity that is assumed to go out from it. One is afraid, as it were, that
psychical life will evaporate if it is not hung up on a soul-substance as on a
secure hook, although one is not so consistent as to examine where this hook is
fastened. The soul-substance seems just as much to be able to need a ``bearer''
as the psychical activities. It is intuitive imagination, not thinking proper,
that does the talking here. It is spatial images and symbols that are confused
with psychological concepts. The urge for fixed, resting images makes it
impossible to hold fast the concept of activity in its peculiarity and
consistency. The spiritualistic theory can thus be explained by the involuntary
--- \emph{materialistic} tendency of the imagination. It is not only (cf. 26) the
one-sided associationist psychology that has let itself be ruled too much by the
analogy with the material atoms; this is also the case with the
monadological-spiritualistic theory. The difference is only that the former
regards the single ideas and feelings as atoms, while the latter makes the
conscious being itself into a kind of atom.

In the second place, the concept of activity gets only a highly mystical
application according to the theory mentioned, since this theory itself admits
that we do not know the bearer of consciousness. We cannot, then, derive any
inner or outer phenomenon from it in such a way that we thereby gain a clear
understanding of it. We get a scientific explanation only where the phenomena are
derived from something whose nature we already know from experience. But
experiences can, after all, never show us that \textit{x.} This
\textit{x} might then in and for itself very well be something other than
``ourselves,'' another being, or other beings, from which the phenomena spring.
According to the theory mentioned we properly know, after all, only so much,
that certain phenomena can\emph{not} find their explanation in external, material
causes; that they are derived from \textit{x,} means, after all, only that their
% --- p. 80 ---
\setcounter{footnote}{0}%
\opage{80}origin is unknown. And even if we knew it, it would stand as an
insoluble riddle how the activity could spring from something that was not itself
active.

It is completely correct that we cannot see through the origin, course, forms,
and laws of psychical activity, nor its relation to other kinds of activity. But
none of these riddles is solved in the least by assuming a particular being
behind the activity. --- If by assuming that something lies behind consciousness
one merely wishes to express precisely this, that there is a very great deal in
psychical activity that we do not understand, then I have no objection to make
to this usage, and I employ it myself. But metaphysical views ought not to be
built on such a foundation.

In the third place, we could very well say what we understand by ``ourselves''
or ``our nature'' without entering into the monadological-spiritualistic
hypothesis. Keeping to experience, we see psychical life unfold as an activity
that passes through different stages and forms. We seek to describe these stages
and forms and to explain the transitions between them. At every given stage of
development we then have psychical life before us with certain properties and
dispositions, and the task is then to understand what, under certain conditions,
can develop out of them. ``Our\emph{selves}'' \emph{or} ``\emph{our nature}''
is, at \emph{every given stage, the properties and dispositions} \emph{that our
phenomena of consciousness show} themselves to be in \emph{possession} of. This
is a wholly clear and intelligible concept, and we need not go back to mystical
depths. Where we are compelled to assume unconscious intermediate members and
dispositions, we make them clear to ourselves by the analogy with the conscious
activities, just as in the material domain we make clear to ourselves the state
of equilibrium as a state in which opposite tendencies to movement mutually
inhibit one another. The concept of property or disposition is certainly more
difficult to make clear and to apply in the psychical domain than in the material
domain\footnote[1]{See on this Psykologi p. 162 cf. p. 76 and 96.}; and there are a great
many dispositions that escape our attention. Therefore our psychological
knowledge will always remain defective. But all the greater importance, then,
attaches to determining the fundamental psychological concepts in as close
agreement with experience as possible.

The more, now, the continued psychical development is conditioned by the given
properties and dispositions, the more activity we ascribe to psychical life; the
more it is conditioned by influences coming from outside (arising independently
of the past of psychical life), the greater passivity we ascribe to it. This
relation between an inner and an outer extends through all the stages of life and
makes itself felt as early as we find occasion at all to assume a psychical life.
When psychical life as a whole is called active, it is first and foremost on
account of the form of unity with which all its content is stamped, and which
cannot be explained by external influence. But psychical life is active also
because every sensation
% --- p. 81 ---
\setcounter{footnote}{0}%
\opage{81}appears with an independence and a quality that is conditioned not
only by the single, momentary impression but also by the simultaneous and
preceding state of consciousness, --- because in the process of association it
spins out, on the basis of an often vanishingly small hint, a long series of
ideas, --- because in thinking proper, which can be defined as association
advancing with attention, it pursues a single connection in all its
consequences, --- because as willing (in the narrower sense) it concentrates its
whole energy in the thought of the realization of a purpose. In all such
psychical expressions the cause lies \emph{predominantly} in a preceding
psychical activity and the dispositions it has produced.

Just as the globes of our world system, according to the Kant-Laplace
hypothesis, had a certain initial velocity and initial direction when they broke
loose from the great nebula, so our life of consciousness --- at the point where
it awakens from the night of unconsciousness --- begins by being active according
to the same laws that it follows later. It begins with dispositions which not
only condition what it can later become, but also give it a certain independence
over against external influences. Development does not begin from an absolute
zero point, but advances from more elementary to more compound forms of activity.
The individual differences depend on the degree of the activity and on the
direction in which it moves from the very first.

\begin{center}\large VI. Psychical and Physical Activity.\end{center}

34. The question how psychical activity is related to physical
activity I have treated at length in my Psykologi (2nd ed. p. 62---80,
92---98), where I have sought to ground the \emph{identity hypothesis}, according to
which physical and psychical activity differ not in their innermost ground but only in
their phenomenal appearance. I shall return to this
question here only in order to answer the objections that Prof. \emph{Kroman}, in the 2nd
edition of his ``Tænke- og Sjælelære'', has raised against this hypothesis, which he
calls, by a not very happy name, Duplicism. (This name is not very happy for
the reason that it diverts attention from what is essential in
the hypothesis, namely the assertion that it is one and the \emph{same} being, one and the
\emph{same} activity, that appears to us under the two forms.) The objections raised
are essentially three: a) There is supposed to be no ground in natural science
for conceiving the relation between the psychical and the material in the
way the identity hypothesis states. b) It is supposed, when the identity hypothesis
is taken as the basis, to be inexplicable how we can know anything of the
material world around us. c) It is supposed to be inexplicable how there can be
identity between two things as different as the outer multiplicity
(of which the body consists) and
% --- p. 82 ---
\setcounter{footnote}{0}%
\opage{82}the inner unity (which is the distinctive mark of psychical life). Each of these
objections we shall test by itself. This testing will give occasion to throw light on
the problem from new points of view.

35. a) If we wish to form a hypothesis at all about a relation that lies
so far on the border of our knowledge as the relation between consciousness and
matter, such a hypothesis can have some significance only if it keeps as
strictly as possible to the principles and laws that already stand
firm in science. A hypothesis that begins by drawing a line through all the
points of view and methods by which we otherwise orient ourselves in our research has
no value whatever, regardless of whether it can be directly refuted or not. The
scientific significance that a hypothesis can have at a point where --- as here
--- verification proper is impossible can only be that of showing what logical
consequences, with regard to this point, we can draw from what already
stands firm for us. A hypothesis about the relation between soul and body that would
begin by setting aside the law of inertia, the first proposition of the
natural science we have so far possessed, such a hypothesis would have no
value. \emph{If} it is \emph{justified, in the problem of the relation between
soul and body, wholly} to \emph{disregard the law of inertia}, then \emph{there exists
no problem at all} of the \emph{kind stated}; in that case, in the
many discussions of this matter, people have tormented themselves with imagined difficulties, and it
becomes inconsistent to want to go on discussing the various hypotheses.
The matter in that case presents no greater difficulties than are involved
in understanding how the movement of one billiard ball can cause the movement of the
other.

But --- it is said --- natural science lays down the law of inertia only as holding
for the relation of bodies to one another. This law then naturally finds no
application here, where what is in question is the relation between soul and body. --- This
view of the law of inertia is not defensible. This law is laid down as a
principle for the explanation of the movement of bodies in general and contains the
demand that when a body changes its state of rest or rectilinear
motion, an \emph{external} cause is to be sought. It thus states
\emph{where} we are to seek the cause when a material change (a
change of motion) takes place. This was the most essential advance that was made
by the laying down of the law. \emph{Galilæi}, the first discoverer of the law, wished precisely to
reject \emph{internal} causes of a body's changes of motion\footnote[1]{Cf. \emph{Wohlwill}: Die Entdeckung des Beharrungsgesetzes
(Zeitschrift für Völkerpsychologie. XV). p. 114.}. \emph{Newton}, to whose
formulation of the law of inertia I keep in particular, says in his ``first
law of motion'': ``A body remains in the state in which it is, rest
or uniform progress, so long as no \emph{external forces} compel it to
change its
% --- p. 83 ---
\setcounter{footnote}{0}%
\opage{83}state'', and the concept of external force (vis impressa) he explains
more closely in the following way (in the fourth definition): ``External forces can be
of different origin; they can arise from impact, pressure or attraction.''
A more recent physicist, \emph{Maxwell}, says (in his work ``Matter and Motion'', § 41),
after having set out the law of inertia: ``The experimental proof of this law's
truth lies in the fact that, every time we meet with a change in a body's
state of motion, we can trace this change back to an \emph{action
between this body} and \emph{another}, i.e. to an \emph{external force}''.

The law of inertia states the method that natural science has followed for as long as
it has existed, and which lies at the basis of \emph{all} the results it
has reached. It is therefore no hasty assertion that we know no
explanation of a material phenomenon that does not consist in showing another
material phenomenon to be its cause. Now it is an acknowledged rule of research
(which Newton has expressed in his third and fourth ``regula philosophandi'') that the
laws experience has clearly established we must accept as valid until new
experiences restrict them. Thus we must also accept the law of inertia as valid
for the material phenomena to which the life of consciousness is bound. Since we
cannot explain any material change in the brain by an intervention of the soul
without coming into sharp conflict with the law of inertia, it is an unavoidable necessity
to look about us for a hypothesis that does not involve such an intervention and thus
does not break the law of inertia. The identity hypothesis fulfills this condition, in that
it conceives the material movements in the brain as external, sensibly intuitable
forms of what goes on in the soul.

Prof. Kroman, to be sure, expresses the proposition of inertia otherwise than Newton and
Maxwell, in that he merely has it state that ``every change of motion
is due to a force''\footnote[1]{Tænke- og Sjælelære. 2nd ed. p. 63.}. He thus leaves out the word
``external''. Definitions are every man's own affair. But this alteration is connected
with his conceiving the proposition of inertia as a simple corollary of the
general proposition of causality\footnote[2]{Vor Naturerkendelse. Kbhvn. 1883. p. 292 f.}. It then states only
that all material phenomena have their cause. The peculiar significance
the proposition of inertia has had in the course of development of modern natural science ---
particularly every time a new hypothesis was put forward or a new discovery was made, ---
the significance, namely, of prompting researchers to look out for definite
\emph{material} causes, one does not understand on Kroman's view.
If the proposition of inertia, in the form in which Galilæi and Newton conceived it,
had not been the constant (more or less clearly conscious) presupposition,
Kant, Laplace and Darwin would not have been led to their hypotheses. ---
But it has also, as I believe, been rightly urged against Kroman's view
of the law of inertia as a mere paraphrase of the general law of causality that it
conflicts with ``the conviction of pretty well all physicists and theorists
% --- p. 84 ---
\setcounter{footnote}{0}%
\opage{84}of knowledge in our day''\footnote[1]{. Vierteljahrsschrift für wissenschaftliche Philosophie. IX. p.
366.}. By authorities, naturally, this question cannot be decided. What
I wish to establish is only that the view of the law of inertia I have kept
to is not only the natural one but also the one that has the most significance for the
special problems, and that in particular it is the only one that makes it intelligible
how the relation between soul and body could become a problem at all. ---

If I have kept here to the law of inertia, it is because I become more and more
convinced that it is in this law that the decisive point lies in the whole
problem. This law is --- when it is understood as Galilæi and Newton understood it
--- the typical expression of the method of natural science, and it comprises within itself,
as in a sum total, all that natural science has taught us about the material
world. It must therefore come especially into consideration when what is in question is to establish
the relation between natural science and the mental sciences.

In the most recent times it is mostly the law of the conservation of energy that people (the first time
\emph{Fechner} in his ``Elemente der Psychophysik'' 1860) have started from in
discussing the problem of the relation between soul and body. It is indeed also
clear that this law cannot hold for the material phenomena that are assumed to
be brought about by, or to bring about, psychical states. Here, however, one could (as
I have already pointed out in my Psykologi. 2nd ed. p. 64)
point out that because the \emph{direction} in which a movement goes is changed,
the sum of energy is not thereby increased. The discussion must then here return to
the law of inertia, which expressly demands an external cause also for changes in
the direction of the movement.

Meanwhile the law of the conservation of energy naturally has a great interest in
the treatment of our problem, and it is no accident that the discussion in the
last decades has turned on it to so high a degree. In the explanation of a
natural phenomenon the task is indeed (according to Newton's regulæ philosophandi) to keep
to such causes as are shown in actual experience and are
\emph{sufficient} to explain the phenomenon. It is therefore justified to start
from the assumption that the material changes in the brain, like all the material
changes that have been most precisely investigated, are subject to the law of
the conservation of energy. We cannot, naturally, assert that it is so before
it has been shown in detail; but if we wish to form a
hypothesis at all, this is still the one that lies nearest. The burden of proof seems to have to
rest on him who would here posit an exception. It is an unscientific
hypothesis that is built up on nothing but such exceptions.

It would now, to be sure, be necessary to posit exceptions if
actual experience compelled us to do so. Before an actual observation no theory,
no law and no principle can stand. But the point is precisely that observation
is altogether impossible here. No one
% --- p. 85 ---
\setcounter{footnote}{0}%
\opage{85}can \emph{observe} how the relation stands between what goes on
in his brain and what goes on in his consciousness, and just as little can
anyone \emph{observe} this relation in others. The ``experiences'' that are
appealed to here are even cruder than those that were appealed to when
it was asserted that experience taught that the sun went round the earth. ---

It seems clear that every view, whatever hypothesis it may finally
land in, or even if it will put forward none at all, must attach great weight to the
considerations advanced here, since it is in any case in the highest
degree remarkable that we are so placed that, starting from the knowledge we possess
of material nature, we not only cannot understand at all the
coming about of the phenomena of consciousness, but are even led to the consequence that they
could be wholly dispensed with in existence, since the \emph{sufficient} causes of
all material phenomena must be found within material nature itself,
\emph{if} natural science goes on advancing in the same way as
hitherto. To throw a veil over this peculiar circumstance does not seem to me
to be good philosophy, since it is one of the most important tasks of philosophy to bring
the consequence of our knowledge to our full and clear consciousness. --- The possibility that
the material phenomena not yet cleared up (to which, first and foremost, so
many of the most important phenomena of organic life belong) should necessitate wholly
different principles and laws from those that have hitherto exclusively stood their test
as guiding thoughts cannot, naturally, be excluded. But this possibility is too
airy a ground to build any hypothesis on. In the gaps that science must at
any time leave open one can build whatever one likes; there
imagination can erect its castles in the air without fear of coming into conflict with
science --- for the time being. --- On the other hand, the circumstance that
these gaps are there means that on this point we cannot reach more than a
hypothesis, and probably can never reach more. All that matters is, in our
forming of hypotheses, to posit as few exceptions as possible to the laws that otherwise
hold, and best of all to posit none at all. ---

From the \emph{difference} between the psychical and the
material phenomena so sharply stressed here it by no means follows that they are ``\emph{absolutely} heterogeneous''\footnote[1]{\emph{Kroman}: Tænke- og Sjælelære. 2nd ed. p. 116. Having first
made the identity hypothesis into ``Duplicism'', he here lets it advance to
dualism.}. A difference, irreducible for us, between two properties or
forms of manifestation need not be a difference in the innermost being: this is indeed
precisely what the identity hypothesis asserts. If we knew the being, the
fundamental activity, that reveals itself in so different a way to our inner and
to our outer sense, we would probably also be able to understand how this double
appearance is necessary. Such knowledge we do not possess, and we will surely be
cut off from attaining it forever. But for that reason we have no ground to
think that there lie
% --- p. 86 ---
\setcounter{footnote}{0}%
\opage{86}two different beings at the basis. The form of state and the color of sulfur
vary with the degree of heat; we cannot derive the form from the color or conversely,
but we can see that both depend on the degree of heat. In the relation between the
psychical and the bodily phenomena we do not know the ultimate cause of the
connection of the two; but we find that they vary in a corresponding way, although they
cannot be derived from one another, and we therefore conclude that they are manifestations of one and the same
being. The usual (monadological-spiritualistic) view
unjustifiably makes the difference so great that it regards soul and body as two
different beings. It stresses the opposition between soul and body where there is
no ground for it, and overlooks it there where the great difficulty precisely lies.

The identity hypothesis faces two ways. It maintains against materialism
that one cannot derive the phenomena of consciousness from matter according to the concept
natural science gives us of matter. It maintains against spiritualism
(Descartes and Lotze) that in this irreducibility there lies no ground
for setting a limit to the validity of the principles of natural science for all
material phenomena, the human and animal brain included. In order to
understand how the life of consciousness can come into being at a certain stage of material
development, one must then assume that in all material activity, besides the
properties and manifestations natural science ascertains, there also cooperates a property not
appearing to external sensation, and that it is from this property that consciousness
comes into being. One can then think of it by analogy with consciousness, or as a
psychical life of vanishing intensity in comparison with the psychical life we know
from the phenomena of our own consciousness. When the throw of a stone calls up sensation and
pain, then it is thus really \textit{a} + \textit{α} that calls up
\textit{b + β,} where \textit{a} means the physical movement in the throw of the stone,
\textit{α} the psychical analogue corresponding to it, \textit{b} the brain process
produced by the throw and \textit{β} the state of consciousness corresponding to this
(sensation and pain). --- By conceiving the matter in this
way we preserve continuity in our conception of nature, at the same time as we
acknowledge the peculiarity of the psychical phenomena as against the material.

The principles and laws of natural science, on this hypothesis, acquaint us only with
the outer side of existence, open to sensible intuition; but existence
has also an inner side, which makes itself known to us in self-observation, and which we
can, everywhere that self-observation does not reach, think of by analogy with what
it shows us. We must extend our concept of matter, put more into it than
natural science, which always builds only on the outer sense, has occasion to.
Such an extension is necessary if the arising of the psychical phenomena is to
become explicable. This extension of the concept of matter does not in the least shake
its validity and application, and seems to be more justified than the
assumption that a material movement should suddenly --- in conflict with all that can be
inferred from the principles and laws that have stood firm hitherto --- be interrupted and
replaced by a psychical process.

% --- p. 87 ---
\setcounter{footnote}{0}%
\opage{87}36. b) With the last consideration we have already prepared the answer
to the second objection against the identity hypothesis. This objection amounted to this,
that our knowledge of the bodily world would become inexplicable if there is no
interaction between soul and body.

This objection has surprised me, since it seems to me that the very
wording of the hypothesis, or its mere name (when one does not rename it), contains an answer.
If there is to be identity of consciousness and brain activity, such that
brain activity is the outer, sensible expression of what is at the same time going
on in consciousness, then it is clear that when a certain brain activity
(\textit{b}) is caused by a certain physical stimulation (\textit{a}), there arises
in the same moment, besides \textit{b}, also a state of consciousness
(\textit{β}), and that if \textit{b} corresponds to \textit{a,} the \textit{β}, in its
innermost ground identical or belonging together with \textit{b}, must also
correspond to it. Whether I say, with the usual spiritualistic-dualistic
view, that \textit{a} produces \textit{b,} and that \textit{b} thereupon
produces \textit{β}, or I say, with the identity hypothesis, that
\textit{a} produces \textit{b,} which is in its ground identical or\footnote[1]{Which mode of expression one will use is indifferent. \textit{b} and
\textit{β} are assumed to be \emph{phenomenally different}, but in \emph{their
ground identical}. On the first point of view we can use the formula \textit{b}
+ \textit{β}, on the second the formula \textit{b = β.}} belonging together with
\textit{β}, must in this connection come to one and the same thing. \emph{Common to both
views} is that \textit{b} does not \emph{come about} \emph{without} \textit{β
also coming about}. Whether one explains the relation between \textit{b} and\textit{ β}
as a relation of causality or as a relation of identity makes, in this
connection, no difference whatever. The identity hypothesis nevertheless has an advantage here too;
for while spiritualism must posit two transitions: from \textit{a} to
\textit{b,} and from \textit{b} to \textit{β}, the identity hypothesis needs only
one transition: from \textit{a} to \textit{(b + β). ---} The view developed in 35 (toward the end)
shows us how we can understand that \textit{a} can
produce not only \textit{b,} but (\textit{b + β}); this becomes possible, namely, only
on the presupposition that \textit{a} in reality means
(\textit{a} + α)\textit{. ---} It is thus quite possible for the identity hypothesis
to hold fast to the continuity between the phenomena of the material world around us and
the phenomena of consciousness that correspond to them. The objection falls away as soon as
one has sufficiently thought one's way into the hypothesis. ---

But besides this \emph{psychological}-\emph{physiological} (psychophysical)
consideration the objection also gives occasion for an
\emph{epistemological} consideration, which is of significance in this
connection. --- All that we know about the material world we know from our
sensations. That these sensations in turn derive from the material world
is an assertion that is no doubt plausible from a popular-dogmatic standpoint, but
which, regarded from an epistemological point of view, rests on a
circular inference. That we --- so far as we adhere to the principle of causality --- cannot
think that our sensations arise ``of themselves'', or that we ourselves produce
them out of nothing (i.e. without \emph{given} occasion), follows of itself. But
it is, epistemologically regarded, self-contradictory to declare the material
world, which
% --- p. 88 ---
\setcounter{footnote}{0}%
\opage{88}we know only from our sense sensations, to be the cause of these. Our
objects of knowledge are constituents of our knowledge itself, of our knowing
consciousness, and the warrant for assuming that there exist things with the same
qualities as those with which they appear in our consciousness, this warrant has
first to be established. Under the influence of the stimulations the knowing subject
receives, there arise in it sensations and ideas, which are worked up according to
the laws of consciousness in such a way that the image of the extended, colored
and sounding world forms itself. But in this image we have by no means the right
without further ado to see an adequate image of the causes whose influence on
the subject sets and keeps its activity of knowing in motion.

The objection mentioned seems without further ado to make the \emph{object} of knowledge
(the object) one with its \emph{cause}. Our objects are the images that
form themselves in consciousness under the influence of certain acting causes; but
immediately the objects and the causes of knowledge have nothing to do with
each other. To use an image (already employed by Lotze): our knowledge
of the material world can be compared to the spark that flies out when
the steel strikes against the flint. But the spark does not resemble the steel, and it would
yet be most unreasonable to assert that one could not understand the spark's
coming about from the flint if it was not a spark that acted on the flint! ---
The whole question, as it is put in the objection mentioned, seems to me
to be wrongly put. It rests on a confusion of the
epistemological relation between subject and object with the psychophysical
relation between consciousness and matter. These two relations must be kept well apart
from each other\footnote[1]{In my work ``Spinozas Liv og Lære'' (Kbhvn. 1877) p. 100 f. I have
shown how this confusion appears in Spinoza. \emph{Windelband} too
(Geschichte der neueren Philosophie. I. Leipzig. 1878. p.
211) and \emph{Tönnies} (Studie zur Entwickelungsgeschichte der Spinoza.
Vierteljahrsschrift für wissenschaftliche Philosophie. 1883. p. 176 f.) have
made this remark. --- In Schelling and Hegel this confusion plays a
fateful role.}. The first relation concerns all our knowledge, not merely
the knowledge of the material world; the second relation is essentially a
border relation between psychology and physiology. This last is thus far more
special than the former, and moreover of a wholly different character.

In our special research and our practical consciousness we keep to the
\emph{objects of knowledge} and investigate their connection with one another. We then form,
on the basis of our sensations, as exact an image of the material
world as possible. But the psychical states in ourselves and others are also
objects of knowledge, although they do not appear to outer sense intuition or in
the form of space. When we now wish to decide how, within our objective
picture of the world, we are to conceive the relation between the material and the psychical,
we can build only on what natural science on the one side and psychology on the other
side teach us, and we reach our hypothesis by drawing
% --- p. 89 ---
\setcounter{footnote}{0}%
\opage{89}the ultimate consequences of these teachings. According to the identity hypothesis,
e.g., the relation is conceived in such a way that the psychical states are related to
the corresponding brain processes in the same way as the concave side of a circular arc
is related to its convex side. Existence thus comes to stand before us
as an arc that can be seen from two sides. --- This is the psychophysical consideration.

\emph{Epistemologically}, on the other hand, we ask on what given basis and
with the help of what principles we have formed this picture of the world, which is the objects of our
knowledge, ordered in a definite connection. The given basis is
in the end our sensations, whereas their cause, and with it the cause of
our whole psychophysical picture of the world, is never itself \emph{given}, but is only
found with the help of a metaphysical hypothesis. In the determination of this cause
the various metaphysical systems come into opposition to one another, in that
each of them determines it in its own way. In psychology and physiology we disregard
this problem and fit our sensations in as links in the
connection of the world that we build up on the basis of these very sensations. But
epistemologically our sensation (just as well as the principles of our thinking)
is something ultimate, an original presupposition, which naturally cannot be
explained by what is built up on the basis of this very presupposition. ---

For the rejection of the objection mentioned it is sufficient to keep to
the psychophysical consideration given above. But since the question can have a
greater range and can easily give occasion to misunderstanding, I have gone
into the epistemological side of the matter.

37. c) The problem of the relation between the psychical and the bodily would not
belong to the greatest and most central problems of human research if
any hypothesis could already now be assumed to have given a final solution. We
cannot get further than to pilot our way forward in these waters, which are filled with
so many skerries, using as seamarks the principles and results that
have hitherto stood their test in our knowledge. As regards the identity hypothesis
in particular, it was only the dogmatic-metaphysical cast of thought
of its great founder that could lead to the belief that the problem had been
completely solved by it. There constantly remains material enough for the wonder
that sets thought in motion. For my part I have, to be sure, the ---
as I believe, not unfounded --- conviction that there is a sharp
corner here, a station that our view of the world must pass. It must here, as at
so many other points, in its conception of existence, abandon the
apparently immediately given way of regarding things. But there always remains
the great strangeness that existence presents two such different sides: on
the one side the \emph{inward unity} in \emph{consciousness}, on the other
side the \emph{outer multiplicity} in \emph{matter}.

% --- p. 90 ---
\setcounter{footnote}{0}%
\opage{90}It must first and foremost be noted here that the difficulty that may
lie in this is common to all theories, and not least makes itself felt in the
usual spiritualistic theory. For how will one be able to connect any
intuitable idea with the interaction that is assumed to take place between
the extended manifold and the non-spatial unity? The riddle here is at the
least just as great, whether one assumes a relation of causality or a
relation of identity. And the riddle is really the same in both cases. All
causality is continuity. What we call cause and effect stands for
scientific knowledge not as two different things, but as two links in one and
the same continuous process. All scientific theories aim at formulating
such a continuity between phenomena that stand for sensory perception
as different things or events. The atomic theory, e.g., makes it possible to
disregard the qualitative changes and the discontinuity following from them in
our sensory experience of the changes of substances, and forms the idea of
a process of exchange and rearrangement by which material particles
move continuously from one point in space to another. If there
now existed a relation of causality between bodily and mental phenomena, then
one and the same continuous process would have to lead from the multiplicity of matter to
the unity of consciousness. I will not now make the demand of the spokesmen of popular
spiritualism that they show us how this transition goes
on in reality. I will only ask whether we can form any
idea whatever of its \emph{possibility}. A break in continuity must in
any case come there where the physiological process stops in order to be replaced
(continued one cannot call it) by the psychical. Let us think of
the resultant of the physical and chemical processes in the brain as a straight line.
This line must then suddenly be broken off; an observer would from a certain moment
on see nothing more. The spatial movement would suddenly be ``absorbed''. ---
Such a break in continuity the identity hypothesis does not involve.

The difficulty that remains for the identity hypothesis I believe I have
already stressed as strongly as possible in my earlier exposition\footnote[1]{Cf. Psykologi. 2nd ed. p. 76, 96, 98, 409.}. The material
world appears to us (in any case in our hypotheses) with a wholly
different continuity from what we are able to think of in the case of the psychical
world. On the other hand, individuality, the
closed unity and whole, appears in return in a quite different, more marked way in the domain of psychical life,
even if it would be dogmatism here to speak of absolute unity, since unity in
our experience appears in all possible degrees, all the way from the highest concentration
to such phenomena where unity seems almost lost in the purely sporadically
emerging elements of consciousness. Within each single individual unity it is
in turn attended with difficulty to maintain continuity; there is a swinging between
% --- p. 91 ---
\setcounter{footnote}{0}%
\opage{91}consciousness and unconsciousness, and the psychical elements seem to
vanish without equivalents being demonstrable. We have insuperable
difficulties in thinking of a law of the conservation of energy in the case of psychical
existence. Both within the single consciousness and between the
various consciousnesses continuity seems broken, and the hypothetical supplement indicated
above (35 and 36) reaches no further than a most vague indication
of something analogous to the psychical elements that could fill the lacunae. If
we were able in our knowledge of existence to connect completely these two
points of view: the continuity of phenomena and their individuality, then
all riddles would really be solved.

It seems to me to be an advantage of the identity hypothesis that it places
both the connection and the opposition between material and psychical nature
in so clear and simple a light. It joins on to the childlike and practical
view that regards the body as the outer expression of man's
personality, whereas on spiritualism the body is a being different from man's
personality. In scientific respects it permits
physiology and psychology each to carry through its own method and yet to work
hand in hand with each other. It anticipates no result. It by no means
asserts that everything in the brain is explained according to the law of inertia and the law of the conservation of matter and
force. But it asserts that no scientific result has hitherto been
reached that did not confirm these laws, and that a hypothesis that begins by
rejecting them can therefore have no warrant or significance. ---

Psychical activity does not lose its real significance in the world because in
its being it is identical with a physical activity. It comes to count for just as much
as the nerve process that is its expression accessible to outer sensation.
The history of development shows us that the nervous system is one of the parts of the organism
that first appear distinctly in the rudiment, and that from the first it occupies
a central position in relation to the other parts of the organism\footnote[1]{See on this: \emph{Kölliker}: Entwickelungsgeschichte des Menschen
und der höheren Thiere. Leipzig 1861. Vorlesung VII, VIII, XXV. --- R. S.
\emph{Berg}: Forelæsninger over den almindelige Udviklingshistorie. Copenhagen
1887. p. 102 --- 104.}. There is here an originally acting rudiment, which shows that
the functions of the nervous system, whose correlates the psychical activities are, are not
mere effects of outer influence, but determined in their nature and
direction at the individual's first coming into being.

Whether in either of the two forms of manifestation we have the absolute nature of existence given is
a question that cannot be answered. Spirit and matter stand, to be sure, as
the chief diversities in our experience; we know no third besides them. But it
is a merely factual duality that we stand before here. There is no contradictory
relation between them, such that a third or a fourth form of manifestation would be a
logical impossibility. Nor is there, as with the three dimensions of space, any
impossibility of intuition present to the assumption of
% --- p. 92 ---
\setcounter{footnote}{0}%
\opage{92}more forms. It is therefore no loose assertion, but a remark
that throws a light on the factual limitation of our knowledge, when we point
out that the solution of the problem of existence may lie far deeper than both the
spiritualistic and the materialistic theories imagine it. It would be
short-sighted dogmatism if here, at this border point, one did not have an open
eye for the various possibilities that present themselves. Only slowly and
reluctantly has the human view of the world reconciled itself to the thought of
the infinity of the physical universe; it refuses in a similar way to
admit how little we are able to see through the inner connection of existence, but
is inclined to think itself sufficiently equipped to be able to establish a
conclusive knowledge.

\begin{center}\large VII. The Conformity to Law of Psychical Activity.\end{center}

38. For psychology, as for every real science, the acknowledgment of the principle of causality
is a condition of life. Only when it is acknowledged that the psychical
phenomena have their causes do the psychological problems come into being. In any
case psychology could never ascertain that the law of causality did not
hold; it could only establish that we could as yet find no causes
for certain phenomena.

Yet it could be conceived that psychological research came upon phenomena
whose reality it had to acknowledge, but whose peculiar nature was precisely
conditioned by the law of causality's not holding, or not holding completely,
so that one had the choice between rejecting the reality of such phenomena
and acknowledging the limitation of the law of causality. Since a fact always has precedence
over a principle, the choice could not be doubtful.

A phenomenon of this kind the moral feelings (particularly the feeling of remorse
or of guilt) are held by many to be, and one then presents the choice either of rejecting
such feelings as grounded in illusion or of assuming causelessness. For
with the feeling of remorse there is supposed to be bound up an inescapable idea that
one in the given case \emph{could have} willed otherwise than one willed, and
the validity of this idea is supposed to be the condition to which the
continued existence of the feeling of remorse is tied.

An assertion such as this confuses the standpoint of psychology with that of ethics. The psychologist
asks merely about the factual arising and connection of the phenomena of consciousness. He
therefore also investigates the nature of that feeling more closely. First he investigates whether, in
the form stated, bound to the idea of having been able to will otherwise, it
is \emph{found} in all men. Now since there are in fact those who know
remorse well enough but have never had the clear consciousness that they could have willed
otherwise when their whole
% --- p. 93 ---
\setcounter{footnote}{0}%
\opage{93}state of mind had not been altered, --- the psychologist concludes that where
that feeling is found in the stated form, it probably owes its arising to quite peculiar
conditions, and he will then seek an explanation of how so
strange an idea as the one mentioned can have fixed itself in consciousness and
connected itself with the feeling of remorse. Indeed, even if the form of the feeling of remorse mentioned
were found in all men, the psychologist as a theoretical researcher would not be able to
take a different position. The ideas that connect themselves with the feelings and
give them their special character need by no means contain the
\emph{real} cause of the feelings. It is sufficient if the individual
who has the feeling regards the idea bound up with it as valid; so
long as he does so, the feeling will keep its character. In the psychology of the feelings
the validity of the ideas that determine the feeling's character
and direction is not tested; if only an association of idea and feeling is entered into, it
is, purely psychologically, of no decisive significance whether the idea is valid
or not.

39. By several different paths one could show how the idea of
the causelessness of an act of will can arise.

In the first place, it can be pointed out that the more a psychical state lays
claim to our interest and attention, the less we notice the
preceding and conditioning circumstances. In a strongly agitated state of feeling
we do not think of its having been prepared and introduced in many ways, not only
by definite outer experiences, but also by preceding states of consciousness
partly of a related, partly of a dissimilar, perhaps even opposite
nature. The whole state comes to stand before us as having set in suddenly, of
itself, as having come into being out of nothing, and we are disinclined to entertain
any genetic explanation of it whatever. This holds most of all in a
serious decision of will. When, after a long inner debate, the decisive
resolution pushes itself forward and in one moment sweeps away all opposing
misgivings, there seems to be a break in continuity. The
new state, in which our whole mind is gathered in the thought of the action, stands
in such opposition to the inner strife and discord of deliberation that, even if we
keep a lively memory of the latter, we still have difficulty in holding fast
to the knowledge that there is a relation of causality between the two states.
Achilles's sudden and decisive self-mastery in the violent scene
with Agamemnon (in the first book of the Iliad) the ancient Greeks could explain
only by Athena's coming down from Olympus and --- visible only to the angry hero ---
pulling him by his locks and holding him back when he had already begun to
draw his sword from its sheath. The ancient Greek consciousness appealed in such
cases to a divine intervention; a modern indeterminist would here have
appealed to ``free will''.

In the second place, one could point out that the decision of will, the resolution, brings with it
% --- p. 94 ---
\setcounter{footnote}{0}%
\opage{94}a peculiar sensation of unrestraint and force, owing to the fact that
the mind is almost completely absorbed in one thought. All inner resistance, all constraint has
fallen away. All thoughts and feelings go in the same direction. No considerations or
barriers are heeded. In the resolution there takes place a breaking through of our real self, of
the circle of elements of consciousness that forms the relatively constant part of the
content of our consciousness, and we thereby have the same sensation as in a full and
deep breath, a sensation that we have now truly come to ourselves, a heightened
self-feeling, which pushes all thought that this state itself nevertheless came into being only
as a consequence of other states far into the background. We feel ourselves sovereign,
unconditioned. But all being caused is being conditioned, limitation, boundness. No
wonder that we cannot in the same moment be wholly absorbed in the present state
and clearly hold fast the idea of its dependence on preceding states!

In the third place, the memory of the state of deliberation --- \emph{if} it
can make itself felt --- may lead one to assume the objective possibility
of a resolution other than the one actually taken. During deliberation various
possible actions presented themselves to our gaze. We entered into
each of them and were torn away from each of them only because the idea of
the others made itself felt. But the idea of a rejected possibility may
have come forward so vividly that it has almost acquired for us the stamp of a psychical reality.
(Cf. on this 32). We indeed often call ourselves to account merely
because we have been able to \emph{think} of performing a certain action. And in the
vivid idea of an action we are --- since activity does not begin at
a certain definite point --- always some way along the road to carrying out the action. Such
an action then does not come to stand in vivid memory as a merely
subjective, pale possibility; we attribute a certain objectivity to it --- and with that the
illusion of a real ability to will, in the single moment, otherwise than
we actually will is strongly prepared. The more clearly and intensely the rejected
possibility stands side by side with the preferred one, the greater the temptation
to attribute to them roughly the same validity. The image that naturally pushes itself
forward shows us two roads; we go only along the one; but the road we do not go along
nevertheless goes on standing before us as an objectively given road, a reality. This
image of the highroad we transfer to the subjective possibilities we have struggled with.

In the fourth place. Even if remorse is not always conditioned by the idea of
\emph{having been able to will} otherwise, it nevertheless contains, where it stirs
with strength, a lively wish that one \emph{had willed} otherwise. The
state in which, in the moment of remorse, one looks back on one's action is
most different from the state in which one took the resolution to
carry out the action. One has trouble identifying the two I's that express themselves in
the two states. Self-recognition is difficult. The identification then easily takes place
in such a way that one transfers something from the present state to the
earlier one. One then attributes to the earlier I at one and the same time the force that
% --- p. 95 ---
\setcounter{footnote}{0}%
\opage{95}led to the reprehensible resolution, and the force that in the
present moment leads to rejecting it. What in the earlier moment
was an Either---Or one makes, by a peculiar joining together, into a
\emph{Both}---And. While remorse does not necessarily \emph{presuppose}
the idea of having been able to will otherwise, it can thus be shown that it
can \emph{produce} such an idea by a quite natural psychological
process. For such displacements go on uninterruptedly in our life of consciousness.

Finally, in the fifth place, an illusion similar to the one last mentioned can make
itself felt when, during a serious deliberation, one thinks of the future\footnote[1]{Cf. on this \emph{Darlu} in Revue philosophique. Juin 1887. p.
574---576.}. When I look ahead in this way, it is clear to me that the future
--- for the domain my will has influence on --- becomes different according
as I choose the one or the other of the actions I am discussing.
The future thus comes to stand before me in a double light: partly as
it would be in consequence of a certain definite action, partly as it
would be without this. There then easily takes place here a pushing together, a kind of
interference, of the two images of the future. The resulting
image of the future comes to stand before me as indefinite, shifting, wavering. It
takes on the appearance that the events of the future are in themselves undetermined,
because we could hold fast no clearly definite image of them, other
images intervening disturbingly. Here a constant oscillation makes itself felt. ---
This oscillation, which in many people expresses itself with regard to the part of the future
their willing can determine, expressed itself for the thinkers of antiquity, even the greatest
of them, also as regards the future as a whole. \emph{Aristotle} taught that
to our problematic judgments there correspond objective possibilities in nature, just as
to our apodictic judgments there correspond necessary relations in nature. The course of the world
was for him undetermined as regards the lower spheres of the universe (the sublunary world).
In the work de interpretatione (cap. 9), which at any rate has been
\emph{attributed} to Aristotle, it is taught that the proposition that of two contradictorily
opposed judgments one must be true, the other false, holds only for the past
and the present, not for the future. When we say: either there will be a sea battle
or not, then it is --- according to this teaching --- \emph{neither} necessary that
there will be a sea battle \emph{nor} that there will be none. Neither of the two can
be true beforehand. As regards the past and the present the possibilities are
exhausted, but as regards the future they stand, not only for investigating
thought, but in real existence itself (however one may more precisely have thought
of this!), still open!\footnote[2]{Cf. on this Aristotelian teaching \emph{George Grote}: Aristotle.
I. p. 183 ff. --- In the most recent times W. \emph{Wundt} has maintained that the principle of
the equivalence of cause and effect does not hold in the case of psychical causality,
but only for physical causality. In opposition to the
principle of equivalence he puts forward the principle of the ever-growing psychical energy. He
connects this assumption with the assertion that our causal explanation
in the psychical domain concerns only the past, not the future. (Ethik. Leipzig
1886. p. 399 f. --- Physiologische Psychologie. 3. Aufl. II. p. 183 f.) It can
naturally not be denied that the mental sciences are able to predict the future in a far lesser degree than
natural science. But is it not the
most natural thing to derive this circumstance from our most deficient and little
exact knowledge in the mental domain? Because meteorological predictions
can far from measure up to astronomical predictions, we do not after all think
that in the meteorological domain another kind of causality makes itself felt
than in the astronomical domain. Our relative ignorance in the psychological
domain has moreover its parallel in no less an ignorance in the domain of
brain physiology, so that what holds for the mental phenomena
would (according to the identity hypothesis, to which Wundt adheres) also have to hold for
the brain functions. --- When new energy seems to come into being in the mental
domain, this is to be explained by analogy with the way in which the organism
during its growth comes to command greater energy. It is the advancing
differentiation of organs and functions that makes a more many-sided taking up and
use of energy possible. But the principle of equivalence is not thereby broken. ---
For the rest I have, as I believe, shown with sufficient emphasis what
difficulties a carrying through of the principle of equivalence in the mental
domain suffers from. --- It must be remarked that Wundt (Ethik p. 412) concedes
that where a real \emph{character} has formed itself, we are able not
merely to explain a man's actions after the fact, but also to predict them.
Thus it is precisely for the highest form of psychical causality that that assertion
does not hold.} --- It is natural that the conviction of
% --- p. 96 ---
\setcounter{footnote}{0}%
\opage{96}firm conformity to law forms itself later in the case of mental life than
in the case of outer nature. And yet it is mostly lack of attention that
is responsible when mental life presents the image of irregular change, and
in practice we in any case rely quite as much on the conformity to law of human nature
as on that of outer nature. We are indeed no less dependent on
men than on nature, and could not live among men if we could not to
a certain degree predict their manner of acting.

40. Psychological research cannot take a different position toward the assertion
of exceptions to the law of causality in the case of certain mental phenomena. Like
all real research, it stands and falls with the principle of causality. It is accordingly also
now fairly generally conceded that, purely theoretically, there would be no ground
for maintaining indeterminism. On the other hand it is still often asserted that moral
motives ought to lead us at certain points to keep our zeal for research in check and
take the understanding captive (if one may otherwise use these expressions here; for the abandonment
of the principle of causality is sometimes depicted as if it were the easiest thing in
the world and cost no struggle). \emph{With this}, as should well be noted,
the \emph{whole question is moved from psychology over} into \emph{ethics}.

I must here permit myself, for my own part, to point to how I have
in all my earlier pronouncements on the deterministic problem maintained
this point of view, that it is an ethical, not a psychological question. That it is
an ethical question does not now, naturally, mean that we can appeal to the
current moral ideas, which are as little sufficiently clear and grounded in an ethical as in a psychological
respect. If the question is really
to be discussed, one must go back to the principles that lie at the basis of all
ethical valuation, and on which therefore all talk of ethical significance and
% --- p. 97 ---
\setcounter{footnote}{0}%
\opage{97}warrant must rest. It must be asked: what ethical significance and
warrant do the feelings (particularly the feeling of remorse or of guilt) have on which
the necessity of assuming the causelessness of the will is often grounded? And it must
then be investigated whether the postulate of causelessness is really an unavoidable logical
consequence of what constitutes the ethical value of those feelings, or whether the
indeterministic assertion cannot very well fall away without this ethical
value suffering any loss. If one does not undertake such an investigation,
it is of no significance whatever to appeal to moral motives.

In my first treatise on this point (the chapter on the freedom of the will in
``Grundlaget for den humane Etik''. 1876) I already sought to show how
``remorse becomes intelligible and is seen in its great ethical significance, even if we
do not assume any absolute freedom'', and how it would be altogether
superfluous ``if it were to presuppose that we could also have acted
otherwise under the given circumstances'' (p. 124). This standpoint I have since
constantly taken as my basis, most recently in the new investigation of this question in my
Etik (1887). My fundamental thought is this: that the explanation of the feeling of remorse that is possible
on the basis of determinism agrees entirely with the significance that must be attributed to remorse from an ethical
standpoint. On my view psychology and ethics meet
in the finest way in the same result. I therefore do not \emph{defend}
determinism. But I \emph{attack} certain current and traditional
ideas that start from an ethically untenable concept of remorse and
imputability and thereby (besides other drawbacks) produce disharmony between
the psychological and the ethical way of regarding things.

Remorse and the feeling of guilt have, as I have sought to show, significance only by
arousing the impulse to change and progress in the will. As mere brooding over what
cannot be undone they are unethical, since they slacken the will and cow courage.
Remorse is produced by the sharp conflict between what my own conviction
acknowledges as right and what memory holds up to me about the nature of my own
willing. When it has been held that, according to determinism, the pain felt
over having willed wrongly must be essentially of the same kind as that felt
at not being more beautiful or more gifted, one overlooks, on the one hand, that in the will
man's real self finds its most heightened expression, so that a
self-condemnation works most strongly and deeply here, and on the other, that while one cannot
do much to help one's gifts and one's beauty (especially not one's beauty),
a change of the will is, on the contrary, not impossible, precisely by virtue of psychological
conformity to law. The feeling of pain thereby acquires here a particularly active character, which
aesthetic or intellectual dissatisfaction with oneself lacks. The
pain and unrest of remorse derive from the cooperation of the two moments adduced, with the
lively feeling of what is right as background. Remorse can be described and explained
just as naturally as any feeling at all can be with the means that stand
at the disposal of psychology. And with this very explanation the
ethical significance is really already given.

% --- p. 98 ---
\setcounter{footnote}{0}%
\opage{98}If now the one on whom an ethically condemning judgment is passed says: ``Such
was my will, once and for all! I could not in the given moment will otherwise!''
--- then the answer is the following: ``That is not what is in question. We merely establish that
the will you displayed was bad. And we tell you so in order energetically
to direct your attention to it, so that your will may become different in
the future.'' --- We have no right at all to pronounce ethical judgments on
others except in order to arouse a better willing in them. \emph{Ethical judgments must themselves
be ethically motivated}. This might seem to follow of itself, but is
(amid the great desire to moralize and operate with moral arguments) all
too much overlooked. Ethical judgments are not justified merely as outbursts in which
one gives oneself air. Even if it is not Pharisaic self-feeling that in this
way finds its outlet, it is still overlooked that being the object of ethical
disapproval is a suffering which --- like all suffering --- must not be inflicted on anyone
except as a necessary means to progress. This follows directly from the first
ethical principles\footnote[1]{See my Etik (1887). p. 73---75.}. Everyone who pronounces an ethical
judgment thereby incurs an ethical responsibility. He operates with forces that belong
among the most strongly and deeply moving in the world, and which therefore ought not to be used
blindly. --- Toward the insane and the unconscious we therefore pronounce no
ethical judgments either. And him who is cowed by the feeling of his powerlessness to keep himself
up in ethical respects we do not cow still more by directing ethical judgments at
him; on the contrary, we seek precisely to strengthen his self-confidence and his self-respect. The
teleological or pedagogical character of ethical judgments becomes manifest here.

When we say to a man: ``You ought not to have willed thus!'' this
statement by no means rests on indeterministic presuppositions. We merely hold up to
him the opposition between his actual willing and the willing he himself must
acknowledge as the right one. In this sharp opposition there lies a sting, a stimulus
for the will. ---

This consideration throws a clear light on the assertion that determinism and ethics are
incompatible. The ethical law lives in consciousness as a motive, and ethical judgments
are pronounced --- in an ethical way --- in order to become motives. And determinism is the doctrine
of the complete motivatedness of the will. It is then not easy to see how there
should be any incompatibility here.

The problem of the ``freedom of the will'' is therefore by no means insoluble. It finds its
clear and definite decision in the settling of the ethical principles. From these
indeterminism must either follow or not follow. I do not underestimate the
difficulties of settling these principles; I dare say that I have
sought to place these difficulties in as clear a light as possible; but if one
wishes at all to take definite principles as the basis of one's ethical reflections,
one must surely be able to become clear about whether indeterminism
% --- p. 99 ---
\setcounter{footnote}{0}%
\opage{99}is a consequence of them or not. And what value an appeal to
the ethical can have that does not take definite principles as its basis, I do not understand.

41. In the new treatment of the question of the freedom of the will that Prof.
\emph{Kroman} has given in the 2nd edition of his Tænke- og Sjælelære, a
similar way of regarding things is now also indicated. He has stated here (p. 310) that
determinism and indeterminism each start from their own concept of guilt, and that the
question remains: which concept of guilt is the right one? --- With this
question the whole debate ought to have begun. It would then have been
more fruitful. But this way of regarding things in any case marks a great
advance in relation to the treatment of the same question in the author's work
``Vor Naturerkendelse'' (Kbhvn. 1883), where the possibility of determinism was conceded only
if one rejects the feeling of responsibility. No attempt was made in that work
to bring forward the possibility that the only ethically tenable
concept of responsibility and remorse might be one that precisely determinism both
could and had to take as its basis. By ``\emph{real} responsibility, \emph{real}
guilt and remorse'' (Vor Naturerkendelse. p. 246) were understood in the earlier work
such forms of these feelings as only indeterminism could acknowledge.
--- Even in his new exposition, however, the author does not yet do equal justice to
both parties. At least his expressions can be understood as if the
concept of guilt compatible with determinism were more superficial than that of indeterminism.
Once one has declared it an ethically open question which concept of guilt
is the right one, one must also take all the consequences. It will then
in particular be clear that \emph{if} there should be ethical misgivings
about indeterminism's concept of guilt, these misgivings become all the greater
the ``deeper'' this concept is.

It has surely greatly surprised most readers of ``Vor Naturerkendelse'' that
the author in his latest exposition protests so decidedly against being
designated an indeterminist\footnote[1]{In the same note in which my esteemed colleague objects to
my having counted him among the indeterminists, he replies to some
remarks I have made in my Etik on the relation of indeterminism to the
religious (particularly the theistic) view of the world. I sought to show that
indeterminism is incompatible with theism, thus with the current religious
view (to which reference had been made), because it abolishes God's omnipotence and
omniscience. Kroman's answer amounts to the following: 1) ``The designation `almighty'
does not mean, after all, that the deity is the \emph{only} cause, but that all other
causes are posited by him and can again be abolished''. 2) The indeterminist can only
not believe in an omniscient God ``if he assumes the deity's knowledge to be
\emph{homogeneous} with human knowledge. But such a view is hardly held nowadays
by very many''. (p. 312.) To this I must remark: 1) Omnipotence
can neither posit nor abolish \emph{anything} \emph{that has no cause}. For
both to posit and to abolish is to be a cause, and the domain of power extends
only as far as causes act; if there is something that has no cause, it
is eo ipso withdrawn from power, and the latter is thus not omnipotence. 2) What
idea can we form of a knowledge that does not consist in insight into the
\emph{grounded connection} of the matter? If divine knowledge is to be so different
from human knowledge that it does not build on the principle of sufficient reason,
then it is only a word we use, not a thought, when we call it knowledge.
(And we are indeed precisely in the domain where words are wont to step into
the place of concepts.) I had said: ``If one maintains that we here stand before a
`mystery', then one must see to how one can distinguish between mystery and
self-contradiction''. It does not seem to me that the rejoinder has given
any answer to this. My protest too against determinism's being particularly the
outcome of an antireligious tendency I must maintain. The history of the Church shows
that the truly religious natures have always inclined to determinism, while
natures that were led more in an ethical than in a religious direction (by a
misunderstanding) maintained indeterminism. Not all determinism is religious; but
all true religiosity must be deterministic, since it must maintain the inward
unity between man and his God, so that man feels himself shone through
by the deity and acknowledges that he can only will and accomplish what the
divine activity brings to fulfillment within him. ---}. The whole defense of
indeterminism, and not least the light in which
% --- p. 100 ---
\setcounter{footnote}{0}%
\opage{100}determinism is placed (see especially ``Vor Naturerk.'' p. 250---253), is of
such a nature that it is in any case excusable if one has thought that
the author stood in a far more intimate relation to indeterminism than that of
being its advocate. However, he must after all be his own best interpreter. And
the main thing is that the new exposition in any case brings out something that
was not stated in the earlier one. --- There are some points of general
interest that I shall here go into somewhat more closely.

In the new exposition it is stated that the greater part of the psychology of the will
can be set out without one's needing to enter into the deterministic
problem. (See p. 304, 312.) Earlier the author seems not to have been so
sure of this. In ``Vor Naturerkendelse'' (p. 244) he was not without misgivings
at a determinist's using expressions such as, e.g., to ``reach''. If every
expression that denotes impulse and striving were to contain the whole problem within itself, then
the author's own exposition of the psychology of the will would have had to acquire a wholly
different character. It is to the decided advantage of his psychology that he has
overcome this misgiving. --- I bring forward this point here only in order
to attach a remark to it. Prof. Kroman connects the circumstance that so
great a part of psychology can be treated without the question of ``freedom''
needing to come up with the slight exactness of psychological research. But when one reduces indeterminism, or rather
indetermination, to so vanishingly small a magnitude as Kroman does, then
one will --- without danger of being refuted --- be able to introduce it even into the
most exact empirical sciences. No chemist will ever be able to disprove
the assertion that when water is formed from oxygen and hydrogen, a billionth part is created out of
nothing, or that this formation of water is by no means a continuous process, but that
in a vanishingly small moment (e.g. a billionth of a second)
the oxygen and the hydrogen are annihilated, and water is immediately created in their place. But that
kind of assertion every empirical science, psychology included, has the right
calmly to leave alone, leaving the burden of proof to those who wish to maintain them.

% --- p. 101 ---
\setcounter{footnote}{0}%
\opage{101}It is the self-contradiction in my esteemed colleague's standpoint on this
question that he has reduced indeterminism to something so vanishing that
it no longer has either any theoretical or any practical significance. Meanwhile
he seeks in his new exposition to show that the very slight part of
the act of will that according to his theory is to be (or can be assumed to be)
causeless is nevertheless not without significance. --- The part of the act of will that is not
causally determined is (it says on p. 308) nevertheless not ``emerged out of nothing'': the I
merely employs voluntarily a certain part of the energy that it possesses
on the whole. --- To this I must answer the following.

Only by using figurative language can one distinguish between the I and its energy.
The I itself expresses itself in a more or less vigorous activity, and it is precisely
the awakening of this activity at the right moment and in the right direction that
matters. If one now wishes to say that a distinction must be drawn between the I itself and
its activity, then the I can mean only the potential energy that can pass
over into actual energy; but this transition itself cannot take place without a
release taking place. If one denies that an impulse is needed for this
release to take place, one eo ipso asserts that something ``emerges out of nothing'',
namely the very release of energy at this moment and in this direction. Even if
one lays ever so great weight on the energy of the human will (and that the
determinist can quite well do), one must still not forget that at the least ``a
differential of a push'' is always needed for this energy to be
employed. Without such impulses our will can act just as little as our
organic life can subsist without breathing. The impulses need not at every moment
come from outside. By virtue of the general psychological laws, partly about the relation
of ideas to one another, partly about the relation between ideas and
feelings, releasing motives will be able to develop out of the very interior of the I. In
this, but also only in this sense, one can say that the I releases its own
energy. But it must then not be forgotten that all ideational activity and
movement of feeling is itself an employment of energy and does not run its course independently
of a willing. We never get, as long as we think clearly and definitely in a psychological
manner, outside a system of interactions partly between inner and outer
conditions, partly between inner conditions among themselves. Only by keeping
the idea of the I in mystical indefiniteness can this be avoided.

When it is said that the I ``\emph{voluntarily}'' employs a part of its energy,
what then is understood by the word ``voluntarily''? If it is taken in the meaning usual
in psychology (cf. 29), it means that in consciousness, prior to
the act, an impulse or a resolution makes itself felt. But then the whole question
indeed comes back, since it must be asked whether this impulse or resolution
arises according to psychological laws or not. ---

An ethical misgiving is aroused by the indeterministic theory at this point, in that
it tempts us to rely on the undetermined possibilities which we cannot beforehand
make ourselves
% --- p. 102 ---
\setcounter{footnote}{0}%
\opage{102}masters of. Whether that incomprehensible minimal addition to the naturally
developed willing will set in at a given moment is, after all, not in our power beforehand.
If we could determine our own future willing, it would indeed be
determined! Determinism, on the other hand, earnestly urges us to lay a good
and firm foundation, in that it shows us that each of our acts of will, indeed each
of our thoughts and feelings, lays a stone in the building of our future
character, and that we cannot at every moment begin afresh. While
indeterminism teaches that we can reap without sowing (and without anything at all
having been sown), determinism enjoins that as we have sown, so shall we also reap.

In the new edition of ``Tænke- og Sjælelære'' the following remark is made in answer to this objection.
``Every time the subject acts in a certain way, its
total nature will be altered a little in the same direction, and it thus
daily cooperates itself in the formation of the character whose momentary nature
gives the single action its most essential stamp'' (p. 308). We then stand here
before the question whether the indeterminist can acknowledge the law of practice in the case of the
will. I cannot see how that is possible without self-contradiction.
Every causeless willing (or causeless part of a willing) is, by the nature of the
case, not determined by anything preceding whatever. Even if a whole habitual
tendency has in some way or other been produced, the undetermined
additions can at every moment strike out in new directions. It is a \emph{pure
accident} whether such a minimal addition acts in the same direction as an
earlier one. \emph{Between the undetermined additions among themselves} there can
\emph{no relation of causality} take \emph{place}. If my ``free'' willing an
hour hence went more easily in a certain definite direction because my ``free'' willing at the
present moment goes in this direction, then it would not be ``free''.

Just as little can I see how, on indeterministic
presuppositions, one can maintain a certain responsibility for the effects of the undetermined additions.
These effects are, after all, necessary, and if responsibility and
causelessness are to stand and fall with each other, then responsibility cannot
extend further than the causeless willing itself goes. As soon as
the law of causality comes into force after the indeterministic intermezzo, and it does so
as soon as the effects of the act of will develop, responsibility too must cease.

It will perhaps be answered to this that deterministic ethics does after all make us
responsible for more than we have directly willed, also for the effects of our resolutions and
actions, which we need not, after all, have willed, although we
willed the action itself. This rejoinder would overlook that what (according to
deterministic ethics) is laid to our charge, and what is impressed on us with
emphasis through punishment, when the talk is of effects that were not
intended, is that we have acted short-sightedly and frivolously. Only for such
effects, therefore, ought we to be called to account as we might possibly have come
to foresee and to avoid by more thorough deliberation. It is therefore altogether
consistent when responsibility in
% --- p. 103 ---
\setcounter{footnote}{0}%
\opage{103}deterministic ethics is extended to include also more remote
effects of the action. But for indeterminism this possibility is cut off,
since it has once and for all tied its concept of guilt to the postulate of
causelessness. If the decisive moment in the decision of will is to be
causeless, then it can never be determined by any thinking about
the possible effects of the act of will. Every warrant for holding anyone responsible
for these is thus lacking.

\begin{center}\rule{2cm}{0.4pt}\end{center}

42. The series of questions that have been discussed in this treatise is not
arbitrarily put together. Step by step the one question has led us
to the other. The guiding thread has been the lawful connection of psychical life and
the forms under which it makes itself known in the domain of ideas and of the will.
The possibility is, as I have also pointed out, not
excluded that one who agrees with me on one point does not for that reason agree with me
on all other points. One may, e.g., grant that I am right as regards the theory of recognition,
and yet disagree with me with regard to the explanation of
association by similarity. But this does not remove the warrant for joining
the preceding investigations into a whole, even if I must concede that the outer
occasion for treating precisely these questions together here was that
researchers for whom I had respect had come to results other than mine with regard to them. A dialogical moment has thereby entered
the treatment, which in the discussion of questions of this kind does no harm. At
least I willingly concede that several of the objections that have been made
against my views have given me occasion to develop them to greater
clarity and precision, or at any rate to attempt to do so.

\begin{center}\rule{2cm}{0.4pt}\end{center}



\end{document}
